\documentclass[12pt,oneside]{book}

% ============================================================
% METADATOS
% ============================================================
\newcommand{\universidad}{Universidad de Buenos Aires (UBA)}
\newcommand{\facultad}{Facultad de Derecho}
\newcommand{\materia}{Contratos Inmobiliarios}
\newcommand{\titulo}{Contratos Inmobiliarios: derechos reales, sociedades y tecnología jurídica}
\newcommand{\titulocorto}{Contratos Inmobiliarios}
\newcommand{\autor}{Isaac Edgar Camacho Ocampo}
\newcommand{\anio}{2026}
% TODO: verificar consistencia entre fuentes originales. Los archivos fuente consignan
% distintos rótulos de materia (Derecho Societario, Derechos Reales, Contratos Inmobiliarios,
% Derecho Notarial y Registral, entre otros), distintas instituciones (Universidad de Buenos Aires;
% Colegio de Escribanos en un caso) y el año 2024 en dos de ellos. Se unificó bajo el rótulo
% "Contratos Inmobiliarios", UBA, Facultad de Derecho, 2026, que es el que se repite en los archivos más recientes.

% ============================================================
% PAQUETES
% ============================================================
\usepackage[utf8]{inputenc}
\usepackage[T1]{fontenc}
% Entorno sin paquete babel en español (no disponible en este sistema):
% se localizan manualmente los nombres generados automáticamente por LaTeX.
\usepackage{babel}
\renewcommand{\contentsname}{\'Indice general}
\renewcommand{\chaptername}{Cap\'itulo}
\renewcommand{\partname}{Parte}
\renewcommand{\figurename}{Figura}
\renewcommand{\tablename}{Tabla}
\renewcommand{\appendixname}{Ap\'endice}
\renewcommand{\today}{\number\day\ de \ifcase\month\or enero\or febrero\or marzo\or abril\or mayo\or junio\or julio\or agosto\or septiembre\or octubre\or noviembre\or diciembre\fi\ de \number\year}

\usepackage[
    top=2.5cm,
    bottom=2.5cm,
    left=3cm,
    right=2.5cm,
    headheight=15pt
]{geometry}

\usepackage{amsmath}
\usepackage{amssymb}
\usepackage{mathtools}

\usepackage{graphicx}
\usepackage{float}
\usepackage{caption}
\usepackage{subcaption}

\usepackage{booktabs}
\usepackage{array}
\usepackage{longtable}
\usepackage{tabularx}

\usepackage{enumitem}

\usepackage[table]{xcolor}
\usepackage[most]{tcolorbox}

\usepackage{fancyhdr}
\usepackage{ragged2e}

% ============================================================
% CONFIGURACIÓN GENERAL
% ============================================================
\graphicspath{{./img/}{./graficos/}{./figuras/}}

\setcounter{tocdepth}{2}
\setcounter{secnumdepth}{2}

\setlength{\parindent}{0pt}
\setlength{\parskip}{0.5em}

\setlist[itemize]{topsep=0.3em,itemsep=0.2em}
\setlist[enumerate]{topsep=0.3em,itemsep=0.2em}

\clubpenalty=10000
\widowpenalty=10000

\definecolor{azulacad}{RGB}{31,56,100}

\captionsetup{font=small,labelfont=bf,skip=4pt}

% ---------- Títulos de parte y capítulo (sobrios) ----------
\makeatletter
\renewcommand{\@makechapterhead}[1]{%
  \vspace*{0.4cm}%
  {\parindent\z@ \raggedright \normalfont
   \ifnum\c@secnumdepth>\m@ne
     {\large\scshape\color{azulacad}\@chapapp\space\thechapter}\par\nobreak\vskip 4\p@
   \fi
   \interlinepenalty\@M
   {\LARGE\bfseries #1}\par\nobreak
   \vskip 4\p@
   {\color{azulacad}\rule{\textwidth}{0.6pt}}\par\nobreak
   \vskip 18\p@}}
\renewcommand{\@makeschapterhead}[1]{%
  \vspace*{0.4cm}%
  {\parindent\z@ \raggedright \normalfont
   \interlinepenalty\@M
   {\LARGE\bfseries #1}\par\nobreak
   \vskip 4\p@
   {\color{azulacad}\rule{\textwidth}{0.6pt}}\par\nobreak
   \vskip 18\p@}}
\def\part{%
  \if@openright\cleardoublepage\else\clearpage\fi
  \thispagestyle{plain}%
  \if@twocolumn \onecolumn \@tempswatrue \else \@tempswafalse \fi
  \null\vfil
  \secdef\@part\@spart}
\def\@part[#1]#2{%
    \ifnum \c@secnumdepth >-2\relax
      \refstepcounter{part}%
      \addcontentsline{toc}{part}{\thepart\hspace{1em}#1}%
    \else
      \addcontentsline{toc}{part}{#1}%
    \fi
    \markboth{}{}%
    {\centering
     \interlinepenalty \@M
     \normalfont
     \ifnum \c@secnumdepth >-2\relax
       {\Large\scshape\color{azulacad}\partname\nobreakspace\thepart}\par
       \vskip 14\p@
     \fi
     {\huge\bfseries #2}\par}%
    \@endpart}
\def\@spart#1{%
    {\centering
     \interlinepenalty \@M
     \normalfont
     {\huge\bfseries #1}\par}%
    \@endpart}
\makeatother

% ---------- Encabezados y pies ----------
\pagestyle{fancy}
\fancyhf{}
\renewcommand{\chaptermark}[1]{\markboth{\thechapter.\ #1}{}}
\fancyhead[L]{\nouppercase{\leftmark}}
\fancyhead[R]{\titulocorto}
\fancyfoot[C]{\thepage}
\renewcommand{\headrulewidth}{0.4pt}
\renewcommand{\footrulewidth}{0pt}
\fancypagestyle{plain}{%
    \fancyhf{}
    \fancyfoot[C]{\thepage}
    \renewcommand{\headrulewidth}{0pt}
}

% ============================================================
% COMANDOS MATEMÁTICOS
% ============================================================
\newcommand{\abs}[1]{\left|#1\right|}
\newcommand{\norm}[1]{\left\lVert#1\right\rVert}
\newcommand{\vect}[1]{\mathbf{#1}}
\newcommand{\deriv}[2]{\frac{d#1}{d#2}}
\newcommand{\pderiv}[2]{\frac{\partial#1}{\partial#2}}

% ============================================================
% ENTORNOS / CAJAS ACADÉMICAS
% ============================================================
\newtcolorbox{definition}[1]{
    enhanced, breakable, title={Definición: #1}, fonttitle=\bfseries,
    colback=blue!3, colframe=blue!50!black, boxrule=0.8pt, arc=2mm
}
\newtcolorbox{concept}[1]{
    enhanced, breakable, title={Concepto clave: #1}, fonttitle=\bfseries,
    colback=teal!4, colframe=teal!50!black, boxrule=0.8pt, arc=2mm
}
\newtcolorbox{important}{
    enhanced, breakable, title={Importante}, fonttitle=\bfseries,
    colback=yellow!8, colframe=orange!70!black, boxrule=0.8pt, arc=2mm
}
\newtcolorbox{warning}{
    enhanced, breakable, title={Advertencia}, fonttitle=\bfseries,
    colback=red!4, colframe=red!60!black, boxrule=0.8pt, arc=2mm
}
\newtcolorbox{example}[1]{
    enhanced, breakable, title={Ejemplo: #1}, fonttitle=\bfseries,
    colback=green!4, colframe=green!40!black, boxrule=0.8pt, arc=2mm
}
\newtcolorbox{formula}[1]{
    enhanced, breakable, title={Regla: #1}, fonttitle=\bfseries,
    colback=purple!4, colframe=purple!50!black, boxrule=0.8pt, arc=2mm
}
\newtcolorbox{exercise}[1]{
    enhanced, breakable, title={Ejercicio: #1}, fonttitle=\bfseries,
    colback=gray!5, colframe=gray!60!black, boxrule=0.8pt, arc=2mm
}
\newtcolorbox{note}{
    enhanced, breakable, title={Nota}, fonttitle=\bfseries,
    colback=white, colframe=gray!50, boxrule=0.6pt, arc=2mm
}
\newtcolorbox{synthesis}{
    enhanced, breakable, title={Síntesis}, fonttitle=\bfseries,
    colback=azulacad!5, colframe=azulacad, boxrule=0.8pt, arc=2mm
}

% ============================================================
% CUADRO CORNELL
% Uso: \cornell{Título}{filas}{Resumen}
% Cada fila: \cq{Pregunta o palabra clave} & Respuesta o nota \\ \hline
% ============================================================
\newcommand{\cq}[1]{\textbf{#1}}
\newcommand{\cornell}[3]{%
  \par\medskip\noindent\begingroup
  \renewcommand{\arraystretch}{1.25}%
  \begin{tabularx}{\textwidth}{|>{\raggedright\arraybackslash}p{0.27\textwidth}|>{\raggedright\arraybackslash}X|}
    \hline
    \multicolumn{2}{|>{\centering\arraybackslash}X|}{\cellcolor{azulacad!12}\textbf{Cuadro Cornell --- #1}}\\ \hline
    \cellcolor{gray!10}\textbf{Preguntas / palabras clave} & \cellcolor{gray!10}\textbf{Notas / respuestas} \\ \hline
    #2
    \multicolumn{2}{|>{\raggedright\arraybackslash}X|}{\cellcolor{azulacad!6}\textbf{Resumen.} #3}\\ \hline
  \end{tabularx}%
  \endgroup\par\medskip}

% ============================================================
% HIPERVÍNCULOS Y METADATOS PDF
% ============================================================
\usepackage[
    colorlinks=true,
    linkcolor=blue!50!black,
    citecolor=green!40!black,
    urlcolor=blue!60!black,
    pdftitle={\titulo},
    pdfauthor={\autor},
    pdfsubject={\materia},
    bookmarks=true
]{hyperref}

% ============================================================
% DOCUMENTO
% ============================================================
\begin{document}

% ---------------- PORTADA ----------------
\begin{titlepage}
\thispagestyle{empty}
\begin{center}
\vspace*{1cm}
{\large \textsc{\universidad}}\\[0.2cm]
{\large \textsc{\facultad}}

\vspace{3cm}
{\Huge \textbf{\materia}}

\vspace{1cm}
{\LARGE Derechos reales, sociedades y tecnología jurídica}

\vspace{1.2cm}
\rule{0.70\textwidth}{0.5pt}\\[0.5cm]
{\large Apuntes de estudio}\\[0.5cm]
\rule{0.70\textwidth}{0.5pt}

\vfill
{\large \autor}\\[0.3cm]
{\large \anio}
\end{center}
\vfill
\begin{flushleft}
\textit{Este es un modesto aporte a los estudiantes y no pretende sustituir el material oficial de la materia.}
\end{flushleft}
\end{titlepage}

% ---------------- ÍNDICE ÚNICO ----------------
\tableofcontents

% ============================================================
% PARTE I — MARCO GENERAL
% ============================================================
\part{Marco general del curso}
\label{parte:marco}

% ------------------------------------------------------------
\chapter{Orientación general: objetivos, criterio de estudio y estructura}
\label{cap:orientacion}

Este material reúne, en un único desarrollo, los contenidos de un curso de \materia{} que combina derechos reales, derecho societario y tecnología jurídica. El primer capítulo explica para qué se estudian estos temas, con qué criterio se los ordena y cómo está organizado el libro, de modo que el lector sepa en cada momento por qué aparece cada concepto y hacia dónde conduce la lectura.

\section{Objetivo del curso: de la organización patrimonial a la tokenización}

El curso se orienta a un objetivo final común: la \textbf{organización patrimonial} y, en particular, la \textbf{tokenización de activos}. Para llegar a él es necesario recorrer temas notariales y de derecho civil, porque sin ellos no se comprenderían la importancia y el porqué de la tokenización: para comprender sus ventajas hay que entender en profundidad qué es lo «no tokenizado». La estructura del curso responde a dos orientaciones simultáneas: una tradicional, correspondiente a los contenidos del derecho inmobiliario, y otra tecnológica. Como todavía no se sabe con certeza hacia dónde evolucionará la implementación tecnológica, es necesario enseñar ambos caminos: el analógico, porque será el de la práctica profesional actual, y el tecnológico, porque representa el futuro de la disciplina.

Una de las fuentes describe la materia como optativa e inscripta en un ciclo de formación en tecnología jurídica del Colegio de Escribanos, dirigida a los operadores jurídicos ---notarios, abogados y martilleros--- que trabajan en un ecosistema en el que la firma digital, la prueba electrónica y los contratos electrónicos constituyen la forma cotidiana de dar fe, de contratar y de litigar. Como ejes tecnológicos se enuncian: blockchain, criptomonedas, inteligencia artificial, seguridad de la información, protección de datos personales, ciberseguridad y firma digital.
% TODO: verificar consistencia entre fuentes originales. Las fuentes atribuyen el curso a distintas instituciones (Universidad de Buenos Aires; Colegio de Escribanos) y rótulos de materia.

\begin{concept}{El derecho real como eje}
Los derechos reales se eligen como eje articulador porque el concepto de derecho real permite advertir que instituciones conocidas como fijas, inamovibles y absolutas (el dominio, por ejemplo) admiten numerosas variantes y matices. Esas variantes permiten dar respuestas a los clientes sin recurrir siempre a los contratos típicos. La organización patrimonial no recae exclusivamente sobre derechos reales, pero ellos son el hilo que atraviesa el recorrido, dado el \emph{numerus clausus} del CCyC.
\end{concept}

\section{Utilidad profesional y personal}

\begin{itemize}
    \item \textbf{Encuadrar antes de resolver.} Saber dónde encuadrar un problema antes de intentar resolverlo es la habilidad que distingue a quien ofrece soluciones reales de quien repite fórmulas (capítulo~\ref{cap:ejercicio}).
    \item \textbf{Herramientas concretas para la práctica diaria.} La sociedad en formación, la compra en gestión de negocios, el dominio revocable, el usufructo y el derecho de superficie no son abstracciones académicas: son las herramientas con las que se estructuran operaciones inmobiliarias, familiares y societarias reales, evitando pasivos, embargos e inhibiciones evitables.
    \item \textbf{Un criterio profesional, no solo técnico.} Saber litigar es necesario, pero también saber cuándo no litigar y cobrar por el conocimiento y no por el precio.
    \item \textbf{Asesorar en el mundo digital.} El mercado inmobiliario ya no se mueve solo con escrituras en papel, dinero fiat y fideicomisos tradicionales: aparecen clientes que desean pagar con criptomonedas, invertir en proyectos tokenizados o consultar si una plataforma es segura. El profesional que no sabe formular las preguntas correctas queda al margen de esa demanda y puede comprometer su responsabilidad si autoriza un acto sin verificar el origen de los fondos.
\end{itemize}

También a título personal: quien done, herede, compre o venda un inmueble, tenga bienes gananciales, forme parte de una sociedad familiar o quiera proteger su vivienda en la vejez sin perder autonomía económica puede encontrarse frente a las figuras que se estudian. Comprender qué se cede a cambio de los servicios digitales gratuitos, reconocer los patrones de una estafa piramidal presentada como proyecto de inversión y distinguir una firma electrónica de una firma digital antes de aceptar condiciones con un clic son herramientas de protección patrimonial que exceden lo profesional.

\section{El hilo conductor}

Los contenidos pueden leerse como las sucesivas respuestas a una misma pregunta: quién es dueño de qué, en qué momento y con qué límites, y cómo se transmite, financia, protege y prueba ese dominio. Una familia que accede a la vivienda, la financia, la transmite o la preserva para generaciones futuras dispone de las diversas figuras que el derecho real pone a su alcance; una empresa familiar se organiza mediante sociedades; y la tecnología modifica la forma de firmar, probar y, en última instancia, tokenizar esos derechos.

\begin{example}{Una casa que un vecino quiere comprar}
La parte tecnológica puede leerse como la historia de una casa que un vecino quiere comprar. Primero es necesario confiar en que el dinero con el que se la pagará vale algo. Luego esa confianza y esa información viajan por las redes que todos usan, con sus beneficios y riesgos. Cuando esas redes se vuelven descentralizadas, aparece una tecnología ---blockchain--- e incluso una inteligencia artificial que empieza a tomar decisiones por sí sola. Como toda herramienta poderosa puede usarse para estafar, por lo que el Estado y los profesionales deben aprender a prevenir el fraude y a regular el juego. Recién entonces, con las reglas claras, esa misma casa puede comprarse de una forma nueva: en cuotas digitales llamadas \textbf{tokens}, tal como antes se compraba en cuotas de un fideicomiso.
\end{example}

\section{Cómo está organizado este libro}

\begin{table}[H]
\centering
\caption{Estructura general del libro}
\renewcommand{\arraystretch}{1.2}
\begin{tabularx}{\textwidth}{>{\raggedright\arraybackslash}p{0.34\textwidth}X}
\toprule
\textbf{Parte} & \textbf{Pregunta que responde} \\
\midrule
I. Marco general & ¿Qué se estudia, para qué y con qué criterio? \\
II. Personas jurídicas y sociedades & ¿Quién puede ser titular y cómo se organiza y responde una sociedad? \\
III. El dominio y sus modalidades & ¿Qué es el dominio, cuáles son sus límites y cómo puede moldearse? \\
IV. Derechos reales sobre cosa ajena & ¿Cómo se desdobla el dominio: usufructo, uso, habitación y superficie? \\
V. Sociedades, condominio y sucesión & ¿Cómo se planifica el patrimonio familiar y empresarial? \\
VI. El documento digital & ¿Cómo se firman, prueban y contratan los actos en entorno digital? \\
VII. Blockchain y tokenización & ¿Cómo funciona blockchain y cómo se tokeniza un inmueble? \\
VIII. Riesgos, prevención y ejercicio profesional & ¿Qué riesgos existen y cómo se ejerce con criterio? \\
\bottomrule
\end{tabularx}
\end{table}

El orden responde al criterio de la escalera: primero los cimientos (personas, dominio y derechos reales), después las paredes (planificación patrimonial y marco digital) y al final el techo (blockchain, tokenización y ejercicio profesional). Cada bloque es un escalón que sostiene al siguiente, y no se puede comprender una etapa sin haber comprendido la anterior.

\section{Criterio de trabajo del curso}

El aprendizaje universitario es un proceso bidireccional: el nivel y la profundidad de cada curso se adaptan a lo que el grupo exige y está dispuesto a dar; la base del programa se mantiene, pero lo que se construye sobre ella depende del compromiso del grupo. La dinámica consiste en que los estudiantes lean, escuchen o investiguen los temas antes de cada encuentro, de modo que el tiempo se dedique a la discusión y aplicación práctica.

\begin{important}
\textbf{Regla de uso de la inteligencia artificial.} Los estudiantes pueden consultar herramientas como ChatGPT para resolver un caso, pero la respuesta obtenida debe pasar siempre por un filtro propio (¿actuaría de esa manera?, ¿sería lo correcto?). Lo que se entrega no es la respuesta de la inteligencia artificial, sino la conclusión personal construida a partir de ella.
\end{important}

Como técnica de estudio se propone, después de cada tema, escribir un breve resumen de cuatro palabras: escribir genera en la memoria algo diferente que solo decir o hablar. Los cuadros Cornell que cierran los capítulos de este libro cumplen una función análoga de recuperación y síntesis. La relación entre docentes y estudiantes, en esta etapa final de la carrera, no debe regirse por una jerarquía rígida sino por un vínculo de consenso y aprendizaje mutuo.

% ============================================================
% PARTE II — SOCIEDADES
% ============================================================
\part{Personas jurídicas y sociedades: libertad de organización y responsabilidad}
\label{parte:sociedades}

% ------------------------------------------------------------
\chapter{Un derecho societario en transformación}
\label{cap:contexto}

Este capítulo ofrece el marco desde el cual se estudian las sociedades en el resto del material. Primero se describe el contexto económico y profesional que impulsa el cambio; luego se ubica la reforma proyectada de la Ley General de Sociedades; y, por último, se presentan las ideas de fondo (persona jurídica, autonomía de la voluntad, jerarquía de valores del intérprete) que permiten leer los capítulos siguientes con un criterio común.

\section{El mundo del trabajo en transformación}

El derecho no puede permanecer ajeno a una transformación del mundo del trabajo que se venía anunciando desde hace años: la globalización se produjo, la pandemia se produjo, y ese cambio entró con toda su fuerza. La pandemia no originó la transformación, sino que aceleró un proceso ya en marcha, modificando de manera estructural la forma en que las personas trabajan y se vinculan profesionalmente.

\begin{example}{La nueva ciudadanía del trabajo remoto}
Dos encuentros personales ilustran el fenómeno: un argentino que se mudó a vivir a Australia y mantenía a la distancia sus costumbres y su hinchada de fútbol como si nunca se hubiera ido; y una persona cuyo trabajo, por ser enteramente en línea, podía ejercerse desde cualquier lugar del mundo, alternando temporadas en distintos países.

La conclusión extraída es que hoy «no importa dónde vivís», porque el trabajo se desprendió del territorio. Esa desterritorialización exige repensar también el derecho que regula la actividad económica y las profesiones.
\end{example}

De estos ejemplos se sigue una advertencia sobre la obsolescencia de ciertas expectativas laborales: profesiones y modos de trabajo que hoy parecen sólidos podrían dejar de existir en pocos años, por lo que conviene pensar la formación con una mirada flexible y de mediano plazo. Esta idea reaparece más adelante, cuando se estudia el rol del profesional del derecho frente a la economía digital (capítulo~\ref{cap:ejercicio}).

\section{La reforma proyectada de la Ley General de Sociedades}
\label{sec:reforma}

El Congreso tiene en agenda un proyecto de reforma integral de la Ley General de Sociedades. Según la valoración expuesta en el curso, comparte algunos de sus lineamientos y discrepa con otros, pero su mérito central es que no se limita a corregir la ley vigente: propone un cambio de paradigma. No copia el régimen anterior, sino que parte de una mirada nueva de cómo funciona hoy el mundo económico.

Para distinguir una reforma «cosmética» de una reforma estructural se recurrió a una metáfora: cambiar el vestido no cambia el cuerpo. Un vestido más juvenil sobre un cuerpo viejo no rejuvenece a quien lo lleva; del mismo modo, modificar aspectos puntuales de la ley sin alterar su estructura de fondo no transforma el sistema.

\begin{example}{El antecedente del proyecto de reforma de 2018}
Como responsable del Instituto de Derecho Comercial del colegio de escribanos, se impulsó en 2018 el debate del primer proyecto de reforma de la ley de sociedades, lo que generó fuertes resistencias internas dentro de la institución. Ese proyecto fue calificado como «muy bueno» desde el punto de vista doctrinario, pero limitado, porque modificaba aspectos puntuales sin transformar la estructura de fondo: seguía siendo el mismo «cuerpo» con un «vestido» distinto. El proyecto actual, en cambio, no reforma sobre lo existente, sino que parte de una base nueva.
\end{example}

Como antecedente de lo que implica un cambio normativo estructural se recordó la sanción del Código Civil y Comercial (CCyC): de un día para el otro dejó de poder aplicarse el código anterior, lo que obligó a todos los operadores jurídicos a un reentrenamiento intensivo, con un curso tras otro, para poder seguir trabajando. Aun cuando el nuevo proyecto de ley de sociedades no llegue a sancionarse, ya está operando un cambio de paradigma en la manera de mirar el derecho societario.

\begin{note}
El concepto de «objeto social» está en plena transformación. Bajo el régimen vigente, la responsabilidad se organiza en gran medida en torno al objeto social; el proyecto de reforma lo desplaza a un segundo plano y da centralidad a la actividad efectivamente desarrollada. Las consecuencias de este cambio para los actos ajenos al objeto se examinan en la sección~\ref{sec:extranos}.
\end{note}

\subsection{Derecho, tecnología e interdisciplina}

El vínculo creciente entre derecho y tecnología forma parte del mismo cambio de paradigma. Gran parte de las ciencias hoy están interrelacionadas con la inteligencia artificial, del mismo modo en que ya ocurre con la medicina. Esto no implica solamente el riesgo de que los profesionales del derecho sean reemplazados por sistemas de inteligencia artificial, sino también, y sobre todo, la posibilidad de usar esas herramientas como colaboradoras del ejercicio profesional.

Como ejemplo de un debate ya instalado en la doctrina se mencionó la discusión sobre si las personas jurídicas pueden integrar el directorio de una sociedad: aunque persisten posturas apegadas a conceptos antiguos, se sostuvo que en algún momento una persona jurídica podrá ser directora. En la misma línea, se anticipó que los próximos interrogantes del derecho societario girarán en torno a las DAO (organizaciones autónomas descentralizadas) y a las sociedades automatizadas, temas que hoy carecen de regulación específica.

\begin{example}{La inscripción registral y el trabajo interdisciplinario}
Una resolución del Registro de la Propiedad Inmueble (RPI) permite inscribir adjudicaciones por disolución de la sociedad conyugal instrumentadas por acta notarial, en las que solamente uno de los adjudicatarios solicita la adjudicación. Se la calificó como una solución práctica y ágil para un problema cotidiano de la práctica notarial.
\end{example}

\begin{important}
El derecho ya no se ejerce en compartimentos estancos: hoy trabaja también con ingenieros y con especialistas de las ciencias duras, entre ellos los de inteligencia artificial. Si la elaboración de una ley de sociedades queda en manos de un legislador sin formación en derecho societario, es probable que el resultado no capte los desarrollos doctrinarios y jurisprudenciales acumulados en la materia; de allí la importancia del trabajo interdisciplinario.
\end{important}

\section{Personas humanas y personas jurídicas}

El punto de partida conceptual es la distinción entre persona humana y persona jurídica, cuya relevancia práctica es directa: la mayoría de las contrapartes con las que un abogado, contador o escribano trabajará (ya sea para contratar, demandar o defender) serán personas jurídicas.

\begin{important}
Antes de contratar, demandar, impugnar una decisión o verificar una nulidad frente a una persona jurídica, el profesional debe verificar tres cuestiones:
\begin{itemize}
    \item si quien actúa en su representación tenía facultades suficientes;
    \item si esa representación era correcta; y
    \item si la toma de decisión (asamblea o reunión de socios) cumplió con los pasos exigidos según el objeto social.
\end{itemize}
\end{important}

\section{Autonomía de la voluntad, legítima y valores del intérprete}
\label{sec:autonomia}

El CCyC unificó los conceptos societarios, comerciales y civiles en un único cuerpo normativo. Como consecuencia, el derecho comercial como rama jurídica autónoma dejó de existir formalmente, aunque en la práctica los tribunales con competencia comercial se mantienen. Se valoró conservar la mirada comercial junto a la civil en el análisis de un mismo caso, porque cada perspectiva aporta matices distintos y complementarios.

\subsection{Libertades y restricciones}

La autonomía de la voluntad es el eje central de la mirada del derecho expuesta en el curso. El CCyC es asimétrico en el modo en que la garantiza: amplia libertad en materia de decisiones personalísimas (por ejemplo, sobre el propio cuerpo y la salud) y fuertes restricciones en materia patrimonial y sucesoria.

\begin{definition}{Porción legítima hereditaria}
El CCyC reserva a los herederos forzosos (descendientes, ascendientes y cónyuge) una porción del patrimonio de la que el titular no puede disponer libremente por testamento.

\textit{Fuente: Código Civil y Comercial de la Nación, régimen de la porción legítima (arts. 2444 y siguientes).}
\end{definition}

Frente a esta institución se expuso una postura personal crítica, autodefinida como «antilegitimaria»: el titular debería poder hacer con su patrimonio lo que quiera, y quien no tuvo hijos que lo cuidaran no debería verse obligado a inventar instrumentos jurídicos para dejar sus bienes a la persona que efectivamente lo cuidó. Se aclaró que se trata de una opinión personal, no de una posición unánime, sostenida incluso trabajando profesionalmente en consultoría de empresas familiares, donde la planificación patrimonial y sucesoria es moneda corriente. La tensión entre autonomía y legítima reaparece en el análisis de un fallo sobre usufructo vitalicio (sección~\ref{sec:legitima}).

\subsection{La pirámide de Kelsen y los valores del intérprete}

Para explicar por qué cada profesional puede llegar a soluciones distintas frente a un mismo caso se recurrió a la pirámide de Kelsen, no ya como esquema formal de jerarquía normativa, sino como metáfora del sistema de valores personales de quien interpreta el derecho: en la punta de la pirámide están los valores con los que cada uno trabaja. Para algunos lo máximo será la Constitución Nacional; para la posición expuesta en el curso, el derecho natural.

\begin{example}{El semáforo apagado}
Frente a un semáforo apagado, quien se detiene igual, sin que nadie se lo exija, actúa por convicción y no por temor a la sanción. Ese es el concepto de responsabilidad que, a juicio del material, todavía falta consolidar en el ejercicio societario: no cumplir la norma para evitar una sanción, sino cumplirla porque corresponde.
\end{example}

En el ámbito universitario y profesional coexisten distintas escuelas de interpretación del derecho societario, cada una con su propia jerarquía de valores. A título ilustrativo, y sin agotar el panorama doctrinario, se mencionaron las posiciones más conservadoras, identificadas con juristas como Ricardo Nissen, frente a lecturas más flexibles, identificadas con autores como Rafael Manóvil; ambas son igualmente válidas como puntos de partida.

\begin{important}
No existe una única forma «correcta» de mirar el derecho societario. El material expone una mirada propia sin desconocer la validez de otros criterios, para no sesgar la información que reciben quienes estudian.
\end{important}

Se citó también al jurista cordobés Efraín Hugo Richard (fallecido recientemente) como referencia doctrinaria central para pensar la relación entre libertad y responsabilidad societaria, resumida en la fórmula que organiza los capítulos siguientes:

\begin{concept}{Libertad y responsabilidad}
«A mayor libertad, mayor responsabilidad.» Cuanto más se flexibiliza la organización de una sociedad, más responsabilidad asume quien la organiza. Esta idea explica tanto la sociedad simple (sección~\ref{sec:simple}) como la S.A.S. (capítulo~\ref{cap:sas}).
\end{concept}

\cornell{Un derecho societario en transformación}{
\cq{¿Qué transformación del contexto impulsa el cambio?} & La globalización y la pandemia, que aceleró un proceso ya en marcha y desterritorializó el trabajo; ello exige repensar el derecho que regula la actividad económica. \\ \hline
\cq{¿Qué distingue a la reforma proyectada de una reforma «cosmética»?} & No corrige la ley vigente sino que parte de una base nueva (cambia el «cuerpo» y no solo el «vestido»); su efecto de cambio de paradigma ya opera aunque no llegue a sancionarse. \\ \hline
\cq{¿Qué debe verificar el profesional ante una persona jurídica?} & Facultades de quien actúa, corrección de la representación y cumplimiento de los pasos de la decisión societaria según el objeto social. \\ \hline
\cq{¿Cómo garantiza el CCyC la autonomía de la voluntad?} & De forma asimétrica: amplia libertad en decisiones personalísimas y fuertes restricciones en lo patrimonial y sucesorio (legítima). \\ \hline
\cq{¿Qué representa la pirámide de Kelsen en este enfoque?} & El sistema de valores personales de quien interpreta el derecho, que explica soluciones distintas ante un mismo caso. \\ \hline
\cq{¿Cuál es la fórmula que vincula libertad y responsabilidad?} & «A mayor libertad, mayor responsabilidad» (Efraín Hugo Richard). \\ \hline
}{El contexto de transformación acelerada del trabajo impulsa un cambio de paradigma en el derecho societario. Ese cambio se apoya en la autonomía de la voluntad, cuyo contrapeso es la responsabilidad de quien elige su forma de organización.}

% ------------------------------------------------------------
\chapter{Los tipos societarios y la crisis de la tipicidad}
\label{cap:tipos}

Este capítulo recorre el sistema de tipicidad de la Ley General de Sociedades N.º 19.550 (LGS), muestra sus límites cuando se aplica sin criterio a estructuras pequeñas y presenta las vías de flexibilización que precedieron a la S.A.S.: la sociedad simple y la sociedad anónima unipersonal.

\section{El principio de tipicidad}

\begin{formula}{Principio de tipicidad societaria}
La LGS exige, para constituir una sociedad regularmente encuadrada en su régimen, la adopción de uno de los tipos previstos por la propia ley. Toda organización que no adopte alguno de esos tipos queda fuera del sistema de tipicidad, con un régimen jurídico distinto (sociedad simple o residual).

\textit{Fuente: Ley General de Sociedades N.º 19.550, régimen de tipicidad.}
\end{formula}

Los tipos previstos, ordenados de mayor a menor complejidad estructural, son:

\begin{itemize}
    \item Sociedad Anónima (S.A.)
    \item Sociedad en Comandita por Acciones
    \item Sociedad de Responsabilidad Limitada (S.R.L.)
    \item Sociedad en Comandita Simple
    \item Sociedad Colectiva
    \item Sociedad de Capital e Industria
\end{itemize}

El criterio que recorre el resto del capítulo es que contar con un tipo previsto por la ley no equivale a que ese tipo sea el adecuado para cada caso. Elegir el tipo correcto, y no automáticamente el más conocido o prestigioso, es la tarea central del asesoramiento profesional.

\section{Sociedad Anónima y Sociedad de Responsabilidad Limitada}

La S.R.L. fue creada por el legislador de 1972 como un tipo intermedio, con un perfil más personalista que la S.A., ante la constatación de que muchos pequeños comerciantes no necesitaban la estructura corporativa de una sociedad anónima. En la práctica, sin embargo, suele ignorarse que es lo mismo ser gerente de una S.R.L. que presidente de una S.A., y se constituye una S.A. para un negocio pequeño con un ropaje jurídico que no le sirve.

\begin{example}{El kiosco constituido como S.A.}
Un comerciante constituye una sociedad anónima para un negocio familiar pequeño (``un kiosco'') y así incurre en el error de asesoramiento paradigmático. La S.A. exige reuniones de directorio cada tres meses, aun cuando los socios sean, por ejemplo, un matrimonio que se ve todos los días.

La consecuencia no es solo el costo y la carga administrativa de sostener el tipo: la falta de libros societarios actualizados y de actas de directorio puede convertirse en una herramienta para impugnar decisiones sociales o para habilitar reclamos de terceros, como el de quien trabajó sin estar registrado y puede valerse de la falta de libros en regla para reforzar su posición.
\end{example}

\begin{table}[H]
\centering
\caption{Sociedad Anónima frente a Sociedad de Responsabilidad Limitada}
\renewcommand{\arraystretch}{1.2}
\begin{tabularx}{\textwidth}{>{\raggedright\arraybackslash}p{0.22\textwidth}XX}
\toprule
\textbf{Aspecto} & \textbf{Sociedad Anónima (S.A.)} & \textbf{Sociedad de Resp. Limitada (S.R.L.)} \\
\midrule
Órgano de administración & Directorio, renovación máxima cada 3 años & Gerencia, por tiempo indeterminado \\
Reuniones obligatorias & Directorio cada 3 meses como máximo & Reunión de socios anual \\
Presentaciones ante IGJ & Mayor carga de presentaciones periódicas & Sin presentación de actas de gerencia \\
Perfil & Corporativo, pensado para grandes estructuras & Personalista, pensado para grupos reducidos \\
\bottomrule
\end{tabularx}
\end{table}

\begin{important}
El tipo societario no es un tema meramente formal, sino una decisión de asesoramiento profesional con consecuencias patrimoniales, tributarias y hasta laborales: hay que elegir la estructura acorde a la realidad económica del negocio, no a la que parece más prestigiosa.
\end{important}

\section{Tipos en desuso}

Tres tipos societarios están en franco desuso: la sociedad en comandita simple, la sociedad colectiva y la sociedad de capital e industria. En toda la trayectoria profesional relatada en el curso nunca se constituyó una sociedad de estos tipos desde cero; solo se trabajó con sociedades preexistentes para transformarlas en otras. Sin embargo, el perfil marcadamente personalista de la sociedad colectiva podría aprovecharse hoy para estructurar determinados contratos asociativos que no encuentran un molde adecuado en los tipos más utilizados.

\begin{example}{La sociedad de capital e industria y el startup actual}
La sociedad de capital e industria (un socio aporta el trabajo y el otro el capital) se equipara con la lógica de financiamiento de un startup actual: una persona aporta la idea y el trabajo, y otra el dinero. El problema es que el tipo, tal como está regulado, ya no refleja adecuadamente esa relación.

En un startup real, quien aporta la idea no puede quedarse con el 50\% del capital social porque no está aportando capital, pero tampoco quien pone el dinero puede quedarse con la idea. Forzar esa relación dentro de una sociedad de capital e industria obliga, en la práctica, a distorsionar en el instrumento constitutivo los aportes reales de cada parte, y expone además a quien aportó la idea a perder poder de decisión si más adelante se diluye su participación en una futura ronda de inversión.
\end{example}

\section{La sociedad simple (ex Sección IV)}
\label{sec:simple}

Con la sanción del CCyC las sociedades civiles, que hasta entonces convivían con las sociedades comerciales tipificadas, desaparecieron como categoría autónoma. Ganó protagonismo, en su lugar, una figura que ya existía en la Ley de Sociedades pero con otro perfil: la sociedad simple, regulada en la Sección IV. Era antes la sociedad de hecho irregular, no querida por el legislador de 1972, y hoy adquiere una preponderancia que antes no tenía.

\begin{definition}{Sociedad simple o residual}
Las sociedades que no se constituyen según alguno de los tipos previstos por la LGS, que omiten requisitos esenciales o que incumplen las formalidades exigidas, quedan regidas por el régimen de la Sección IV de la ley, conocido en la práctica como régimen de la «sociedad simple». A diferencia del sistema anterior a la reforma, este régimen ya no sanciona la irregularidad con la nulidad, sino que reconoce a estas sociedades personalidad jurídica plena, con reglas propias de responsabilidad y oponibilidad.

\textit{Fuente: Ley General de Sociedades N.º 19.550, Sección IV, arts. 21 a 26.}
\end{definition}

La ventaja central de este régimen es que permite a los socios pactar libremente su organización interna sin sujetarse a la rigidez de un tipo previsto, a costa de asumir una responsabilidad distinta frente a terceros. Es la aplicación directa de la fórmula «libertad con responsabilidad» (sección~\ref{sec:autonomia}): quien elige el «vestido» que quiere con la sociedad simple asume, a cambio, otra responsabilidad.

\section{La Sociedad Anónima Unipersonal (S.A.U.)}
\label{sec:sau}

La S.A.U. es resultado de un largo debate doctrinario: en un congreso de derecho societario se discutieron más de veinte proyectos distintos para regular la posibilidad de que una sola persona constituyera una sociedad, cada uno con su propio nombre. El proyecto original, tal como fue presentado en el debate parlamentario del CCyC, apuntaba a dar una herramienta ágil al pequeño emprendedor que no quería (o no podía) sumar un socio formal solo para cumplir con el requisito de pluralidad.

\begin{example}{El «socio de paja»}
Antes de la sociedad unipersonal, era habitual que un cliente que necesitaba constituir una sociedad pero no tenía con quién asociarse ofreciera la calidad de socio a su cónyuge, a un hermano con el que estaba distanciado o a un familiar de edad muy avanzada, solo para cumplir el requisito formal de pluralidad.

Ese «socio de paja», aunque tenga una participación mínima (por ejemplo, del 5\%), conserva derechos societarios como el de convocar a asamblea, lo cual puede generar conflictos no queridos por el socio mayoritario.
\end{example}

Cuando el proyecto pasó por el Poder Ejecutivo antes de su sanción definitiva se modificó su perfil original y se lo redujo a la actual Sociedad Anónima Unipersonal, con requisitos que desvirtuaron el objetivo inicial de asistir al pequeño emprendedor.

\begin{formula}{Requisitos de la Sociedad Anónima Unipersonal}
La constitución de una S.A.U. exige:
\begin{enumerate}
    \item un capital social mínimo de \$30.000.000, actualizado por el Decreto N.º 209/2024 (que modificó el art.~186 de la LGS, elevando el monto que regía desde 2012);
    \item la integración del 100\% del capital social en el mismo acto constitutivo, a diferencia de la sociedad anónima pluripersonal, donde la ley permite integrar un 25\% al inicio y el saldo dentro de los dos años; y
    \item sindicatura obligatoria, con el consecuente sometimiento al régimen de fiscalización estatal permanente del art.~299 de la ley.
\end{enumerate}
\textit{Fuente: Ley General de Sociedades N.º 19.550, arts. 186 y 299; Decreto N.º 209/2024.}
\end{formula}

Se subrayó especialmente la sindicatura obligatoria como el mayor obstáculo para el pequeño emprendedor, que no está en condiciones de afrontar ese costo adicional ni desea quedar sujeto a un régimen de control estatal reforzado.

\begin{important}
La S.A.U., pensada originalmente como una herramienta de simplificación para el pequeño emprendedor, terminó siendo utilizada casi exclusivamente por grandes grupos económicos y sociedades extranjeras que necesitan una sociedad vehículo sin socios externos. El pequeño comerciante, en la práctica, sigue necesitando otras alternativas (la sociedad simple o la S.A.S.) para evitar recurrir al «socio de paja».
\end{important}

\section{Suscripción e integración del capital social}
\label{sec:capital}

Dos conceptos que suelen confundirse en la práctica deben distinguirse con precisión.

\begin{definition}{Suscripción e integración}
\textbf{Suscripción:} monto de capital que el socio se compromete a aportar a la sociedad.

\textbf{Integración:} monto que el socio efectivamente aporta a la sociedad en un momento dado.
\end{definition}

La posibilidad de suscribir un capital e integrar solo una parte en el acto constitutivo es una excepción reservada exclusivamente a los aportes en dinero en efectivo. Quien aporta cualquier otro tipo de bien (dólares, un inmueble, etc.) debe suscribir e integrar la totalidad. Para los aportes dinerarios, la ley permite integrar un porcentaje inicial (25\% en la sociedad anónima pluripersonal) y completar el saldo dentro de un plazo de hasta dos años, según lo que determine el directorio.

\begin{example}{La venta del departamento}
Un cliente quiere constituir una sociedad con un capital social elevado, pero no dispone en el momento de la totalidad del dinero. La solución práctica es suscribir el total del capital previsto e integrar solo una parte (por ejemplo, un 25\%) al momento de la constitución, comprometiéndose a integrar el saldo dentro de los días siguientes, cuando cobre, por ejemplo, el producido de la venta de un inmueble. La ley da así un margen de planificación financiera, siempre que se trate de aportes dinerarios.
\end{example}

\cornell{Tipos societarios y crisis de la tipicidad}{
\cq{¿Qué exige el principio de tipicidad?} & Adoptar uno de los tipos previstos por la LGS; quien no lo hace queda bajo un régimen distinto (sociedad simple). \\ \hline
\cq{¿Por qué el «kiosco S.A.» es un error?} & Impone directorio trimestral y carga administrativa innecesaria; la falta de libros y actas puede usarse para impugnar decisiones o reforzar reclamos de terceros. \\ \hline
\cq{¿Qué es la sociedad simple hoy?} & El régimen residual de la Sección IV (arts.~21 a 26): personalidad jurídica plena, sin nulidad por irregularidad, con libertad de organización a cambio de otro régimen de responsabilidad. \\ \hline
\cq{¿Por qué la S.A.U. no sirvió al pequeño emprendedor?} & Exige capital mínimo elevado, integración del 100\% al constituir y sindicatura obligatoria (art.~299); quedó para sociedades vehículo. \\ \hline
\cq{¿Qué diferencia hay entre suscribir e integrar?} & Suscribir es comprometer el aporte; integrar es aportarlo efectivamente. Integración parcial solo es posible con aportes en dinero. \\ \hline
}{El sistema de tipicidad obliga a elegir un tipo, pero el tipo más prestigioso no siempre es el adecuado. La sociedad simple y la S.A.U. fueron dos respuestas a esa rigidez; la primera flexibiliza a cambio de otra responsabilidad y la segunda, pese a su objetivo original, quedó reservada a grandes estructuras.}

% ------------------------------------------------------------
\chapter{La sociedad por acciones simplificada y el valor de la forma}
\label{cap:sas}

La rigidez del sistema clásico y las limitaciones de la S.A.U. explican por qué la Sociedad por Acciones Simplificada (S.A.S.) se presenta como la alternativa más flexible y digital para el emprendedor real. Este capítulo describe su origen y sus rasgos, la compara con los tipos de mayor uso práctico y cierra con la reflexión sobre la relación entre forma y fondo.

\section{La Ley de Apoyo al Capital Emprendedor y el nacimiento de la S.A.S.}

La S.A.S. nace dentro de la Ley de Apoyo al Capital Emprendedor, con una técnica legislativa que consiste en incorporar figuras societarias novedosas a través de una ley específica, por fuera de la Ley General de Sociedades, cuando modificar esta última resulta políticamente más difícil.

\begin{definition}{Sociedad por Acciones Simplificada (S.A.S.)}
La Ley N.º 27.349, de Apoyo al Capital Emprendedor, creó la Sociedad por Acciones Simplificada como un nuevo tipo societario por fuera de la LGS, con constitución íntegramente digital, objeto amplio, mayor libertad para pactar la representación y la administración, y aplicación supletoria en primer lugar del CCyC y, en segundo lugar, de la LGS ante lagunas normativas.

\textit{Fuente: Ley N.º 27.349, de Apoyo al Capital Emprendedor.}
\end{definition}

Sus diferencias centrales frente a los tipos tradicionales son:

\begin{itemize}
    \item permite la unipersonalidad sin las exigencias de capital mínimo elevado ni sindicatura obligatoria que pesan sobre la S.A.U.;
    \item su constitución puede tramitarse en soporte íntegramente digital, sin instrumento en papel;
    \item el objeto social puede redactarse en forma amplia, sin la precisión y determinación que tradicionalmente exige la LGS; y
    \item permite diseñar formas de representación y administración más flexibles, incluyendo mandatos por tiempo indeterminado.
\end{itemize}

\begin{note}
Ante una laguna normativa en la S.A.S., el orden de prelación es inverso al de las sociedades típicas: corresponde recurrir primero al CCyC y, solo si este tampoco da respuesta, a la LGS.
\end{note}

La técnica de incorporar institutos nuevos por leyes especiales tiene antecedentes: la Ley N.º 24.441 introdujo el fideicomiso en el ordenamiento argentino sin necesidad de reformar el Código Civil entonces vigente.

\begin{table}[H]
\centering
\caption{S.A., S.R.L. y S.A.S.: comparación de exigencias centrales}
\renewcommand{\arraystretch}{1.2}
\begin{tabularx}{\textwidth}{>{\raggedright\arraybackslash}p{0.20\textwidth}XXX}
\toprule
\textbf{Aspecto} & \textbf{S.A.} & \textbf{S.R.L.} & \textbf{S.A.S.} \\
\midrule
Capital mínimo & \$30.000.000 (Dto. 209/2024) & Sin mínimo legal & Sin mínimo legal elevado \\
Constitución digital & No: requiere escritura o instrumento privado & No: requiere instrumento privado & Sí: íntegramente digital \\
Sindicatura obligatoria & Solo en casos del art.~299 LGS & No & No \\
Norma supletoria & LGS & LGS & Primero CCyC, luego LGS \\
\bottomrule
\end{tabularx}
\end{table}

\section{Resistencias doctrinarias y forma frente al fondo}
\label{sec:forma}

La S.A.S. generó resistencia en un sector de la doctrina societaria más conservadora, para el cual la flexibilidad del nuevo régimen ponía en riesgo la seguridad jurídica que aportaba el sistema clásico de tipicidad. Como anécdota se relató la reacción de colegas de mayor trayectoria ante una charla sobre esta figura: le señalaron que estaba «destruyendo la notarial», al restar protagonismo a la escritura pública en la constitución de sociedades.

Frente a esa resistencia se sostuvo que el soporte de un acto jurídico (papel, disco rígido o nube) no determina su validez ni su eficacia, sino el cumplimiento efectivo de la función que ese acto está llamado a cumplir.

\begin{concept}{La forma no hace el fondo}
Lo importante no es que la escritura pública sea el vehículo necesario para que el empresario opere, ni que el acto conste en papel, en un disco rígido o en la nube, sino la tarea que efectivamente se realiza.
\end{concept}

En esa misma línea se expresó una expectativa favorable respecto de un protocolo notarial digital: la certificación digital de una firma cumple la misma función de dar fe pública que el protocolo en papel, sin necesidad de sostener las cargas físicas tradicionales del oficio notarial. Este tema se retoma al estudiar la firma digital y el rol del escribano (capítulo~\ref{cap:firma}).

\cornell{La sociedad por acciones simplificada}{
\cq{¿Cómo y dónde nace la S.A.S.?} & Con la Ley N.º 27.349, de Apoyo al Capital Emprendedor, por fuera de la LGS, técnica usada cuando reformar esta última resulta políticamente difícil. \\ \hline
\cq{¿Qué rasgos la distinguen?} & Unipersonalidad sin capital mínimo elevado ni sindicatura obligatoria, constitución digital, objeto amplio y administración y representación flexibles. \\ \hline
\cq{¿Qué norma se aplica ante una laguna?} & Primero el CCyC y, subsidiariamente, la LGS (orden inverso al de los tipos clásicos). \\ \hline
\cq{¿Por qué generó resistencia?} & Un sector conservador vio en su flexibilidad un riesgo para la seguridad jurídica y para el protagonismo de la escritura pública. \\ \hline
\cq{¿Qué significa «la forma no hace el fondo»?} & Que la validez y eficacia del acto dependen de la función que cumple y no de su soporte (papel, disco rígido o nube). \\ \hline
}{La S.A.S. es la respuesta más flexible y digital a la rigidez de la tipicidad clásica. Su aceptación plantea la cuestión de la forma: para este enfoque, lo decisivo es la función cumplida por el acto y no su soporte.}

% ------------------------------------------------------------
\chapter{La sociedad en formación y la adquisición de inmuebles}
\label{cap:formacion}

El dominio se analiza en este capítulo desde otro ángulo: el de \emph{quién} puede ser titular de él. Para ello es necesario precisar el momento en que nace la sociedad como sujeto de derecho, qué puede hacer antes de su inscripción y quién responde por sus actos, con especial atención a la adquisición de inmuebles. El capítulo termina con el ciclo de vida final de la sociedad (liquidación, cancelación y reconducción), que conecta con la regularización de titularidades inmobiliarias (sección~\ref{sec:prescripcion}).

\section{Nacimiento de la sociedad y contrato plurilateral de organización}
\label{sec:nacimiento}

La sociedad nace con el \textbf{contrato constitutivo}. El ordenamiento argentino adopta la teoría del contrato plurilateral de organización, y no la teoría de la institución, de modo que el registro es meramente \emph{declarativo} y no constitutivo. Ese contrato da derecho al nacimiento de una persona jurídica diferente, que queda regularmente constituida con la inscripción en la Inspección General de Justicia (IGJ), en la Ciudad Autónoma de Buenos Aires, en la Dirección Provincial de Personas Jurídicas o en el órgano competente según cada provincia.

La IGJ y sus equivalentes provinciales dependen de cada provincia y no de la Nación, porque se trata de derechos de forma, no de fondo: todo derecho de forma que no haya sido delegado a la Nación pertenece a las provincias, conforme al reparto de competencias de la Constitución Nacional y del Código.

\begin{note}
Los comentarios expuestos en el curso sobre cómo debería reorganizarse la IGJ en el marco de la autonomía porteña constituyen una posición personal, expresamente presentada como tal, y no una descripción del derecho vigente. El régimen actual es que la IGJ de la Nación conserva competencia sobre las sociedades con sede en la Ciudad Autónoma de Buenos Aires.
\end{note}

\begin{definition}{Sociedad en formación}
Estado en el que se encuentra la sociedad en el transcurso que va desde la constitución (firma del contrato constitutivo) hasta la inscripción registral.
\end{definition}

\section{Capacidad y responsabilidad de la sociedad en formación}
\label{sec:capacidad}

Uno de los puntos técnicos más importantes es la distinción entre capacidad y responsabilidad. Dado que se parte del contrato plurilateral de organización, la sociedad en formación tiene plena capacidad jurídica desde su constitución: puede realizar todos los actos que quiera y tiene plena capacidad para adquirir derechos y contraer obligaciones. Lo único que varía es un problema de responsabilidad, no de capacidad: si esta recaerá sobre la sociedad, sobre la sociedad y quienes realizaron el acto, o directamente sobre quienes lo realizaron y no sobre la sociedad.

\begin{table}[H]
\centering
\caption{Capacidad y responsabilidad en la sociedad en formación}
\renewcommand{\arraystretch}{1.2}
\begin{tabularx}{\textwidth}{>{\raggedright\arraybackslash}p{0.24\textwidth}X}
\toprule
\textbf{Concepto} & \textbf{Alcance en la sociedad en formación} \\
\midrule
Capacidad & Plena desde la firma del contrato constitutivo: la sociedad puede adquirir derechos y contraer obligaciones. \\
Responsabilidad & Depende del acto realizado: puede recaer sobre la sociedad, sobre quien actuó, o sobre ambos (secciones~\ref{sec:adquisicion} y~\ref{sec:extranos}). \\
\bottomrule
\end{tabularx}
\end{table}

\section{Adquisición de inmuebles por sociedades en formación}
\label{sec:adquisicion}

El problema práctico central se presenta a través de un diálogo hipotético: un cliente anuncia que va a comprar un inmueble pero no quiere ponerlo a su nombre por temor a un embargo.

\subsection{La compra en gestión de negocios y sus riesgos}

Frente a ese planteo, el material critica una práctica notarial extendida: la llamada \textbf{compra en gestión de negocios}, en la que una persona compra «para sí y para una sociedad que oportunamente aceptará» (con dinero de la sociedad en formación). Se sostiene que es un error, porque trae consecuencias jurídicas que la gente no atiende.

\begin{example}{Juancito, Santi y los riesgos de la compra en gestión}
Dos personas con nombres ficticios compran un inmueble a título personal, con la promesa de que la sociedad lo aceptará más adelante. Los riesgos son:
\begin{itemize}
    \item mientras la sociedad no acepte dentro del período fiscal correspondiente, el comprador debe declarar el inmueble en su propio impuesto a las ganancias, porque forma parte de su patrimonio personal;
    \item si el comprador fallece antes de que la sociedad acepte, debe iniciarse su sucesión para poder completar la operación;
    \item si al comprador le trazan una inhibición general de bienes (por ejemplo, por no pagar una cuota alimentaria), debe levantarse antes de que la sociedad pueda aceptar; y
    \item si al inmueble le trazan un embargo (por ejemplo, por una deuda del colegio de sus hijos), debe levantarse antes de completar la operación.
\end{itemize}
\end{example}

\begin{example}{La sociedad no constituida al momento de la compra}
Una persona adquirió un inmueble en carácter de gestor de negocios para una sociedad que, oportunamente, aceptaría la compra. La sociedad ni siquiera había sido constituida al momento de la adquisición; estaba conformada por dos cuñadas. Falleció el marido de una de ellas sin que la sociedad hubiera aceptado, por lo que el inmueble ingresó en la sucesión. La viuda, con cuatro hijos menores, solicitó autorización judicial para vender. El mismo juez que autorizó la venta dispuso luego que no se inscribiera la escritura. El caso llegó a Cámara, que revocó la decisión de primera instancia.

El encuadre incorrecto del contrato (adquirir como gestor de una sociedad que en ese momento no existía ni disponía de los fondos) generó un conflicto que se prolongó en distintas instancias judiciales. Sobre la importancia del encuadre, véase la sección~\ref{sec:encuadre}.
\end{example}

\subsection{¿Cuándo sí conviene la compra en gestión?}

Existe un supuesto legítimo: cuando uno de los destinatarios de la compra está temporalmente ausente, pero su regreso e incorporación están asegurados en el corto plazo, lo que evita una escritura adicional posterior.

\begin{example}{Compra a favor de dos hijos, con uno ausente}
Un padre quiere comprar un departamento a favor de sus dos hijos, pero uno de ellos no pudo llegar (está en Mendoza o en España) y se sabe que regresa al mes. Compra entonces el hermano presente, para sí y para su hermano, que oportunamente aceptará.
\end{example}

\subsection{La solución preferida: encuadrar dentro del objeto social}

Para comprar un inmueble a nombre de una sociedad todavía no inscripta se propone un método de análisis en dos pasos:

\begin{enumerate}
    \item \textbf{Verificar el objeto social.} Determinar si la compra del inmueble está comprendida en el objeto social (por ejemplo, si su objeto es la construcción o la compraventa de inmuebles). En ese caso no hay problema: la sociedad en formación puede comprar directamente mientras se tramita la inscripción.
    \item \textbf{Autorizar expresamente el acto si excede el objeto social.} Incluir una autorización expresa en la parte del acta constitutiva donde hablan las partes, incluso para actos ajenos al giro ordinario (por ejemplo, alquilar un local o contratar personal), de modo que el acto quede imputado directamente a la sociedad y no a quien lo firmó.
\end{enumerate}

Una herramienta administrativa práctica es el \textbf{formulario 185}, trámite que permite obtener un número de CUIT para la sociedad en formación (y no solamente cuando la sociedad sale inscripta) mientras se completa la inscripción definitiva.

\subsection{El plazo de duración: desde la constitución o desde la inscripción}

La ley permite computar el plazo de duración de la sociedad desde la constitución o desde la inscripción. La preferencia expuesta en el curso es fijarlo \textbf{desde el contrato constitutivo}, por una razón práctica: permite lograr que todos los actos que realiza la sociedad sean imputados a la sociedad y no a quien los firmó.

\begin{important}
Si no se previó una autorización expresa en el acta constitutiva y el objeto social no cubre el acto realizado, quien firmó la escritura sin estar autorizado corre el riesgo de que el acto se le impute personalmente, ya que la sociedad aún no está definitivamente inscripta.
\end{important}

\section{Actos notoriamente extraños al objeto social}
\label{sec:extranos}

Se plantea un supuesto de conflicto: un director compra un inmueble para una sociedad cuyo objeto social nada tiene que ver con inmuebles (por ejemplo, una sociedad dedicada a la compra o confección de remeras). El procedimiento correcto consiste en esperar a que la sociedad quede regularmente constituida (inscripta) y, en ese momento, someter la operación a una \textbf{reunión de directorio}, donde la mayoría de los directores debe aprobar o rechazar la compra. Las consecuencias serán diferentes según que la mayoría apruebe o no. Esto muestra la importancia real de las reuniones de directorio, muchas veces subestimadas por los propios profesionales: no son un mero trámite formal.

\subsection{Responsabilidad según la aprobación del directorio y de la asamblea}

El circuito de responsabilidad se desarrolla paso a paso, y la responsabilidad se va sumando entre los distintos órganos societarios.

\begin{table}[H]
\centering
\caption{Responsabilidad por actos extraños al objeto social}
\renewcommand{\arraystretch}{1.2}
\begin{tabularx}{\textwidth}{>{\raggedright\arraybackslash}p{0.34\textwidth}X}
\toprule
\textbf{Supuesto} & \textbf{Quiénes resultan responsables} \\
\midrule
Acto dentro del objeto social & No hay responsabilidad especial: el acto es propio del giro ordinario y compromete normalmente a la sociedad. \\
Acto extraño al objeto, aprobado por directorio y por asamblea & Responden la sociedad, los directores que aprobaron el acto y quien lo realizó. \\
Acto extraño al objeto, no aprobado por la asamblea & Responden únicamente los directores que aprobaron el acto y quien lo realizó (no la sociedad). \\
\bottomrule
\end{tabularx}
\end{table}

\begin{note}
Se informó, aclarando que se trata de un proyecto de ley todavía no vigente, que en el eventual proyecto de reforma societaria la figura del acto «notoriamente extraño» al objeto social quedaría desdibujada, porque el objeto social pasaría a definirse de otra manera (sección~\ref{sec:reforma}). Mientras el proyecto no sea aprobado, sigue rigiendo el régimen explicado en esta sección.
\end{note}

\subsection{Actas volantes y libros societarios durante la formación}

Se plantea el problema de documentar decisiones cuando todavía no existen libros rubricados, por tratarse de una sociedad en formación con objeto no inmobiliario. La solución consiste en labrar un \textbf{acta de asamblea} (denominada coloquialmente «acta volante») con la unanimidad de los socios o con las mayorías exigidas por la ley, dejando constancia expresa de que, al no existir aún libros rubricados, el acto se transcribirá en el libro correspondiente cuando este sea rubricado.

\begin{important}
Para tener validez, un acta volante debe reunir los mismos requisitos que un acta ordinaria posterior a la inscripción: unanimidad o las mayorías exigidas por la ley. Si se exige la participación de todos o de una mayoría determinada y solo concurre una parte de los socios (por ejemplo, dos de ocho), el acto no puede realizarse.
\end{important}

El mismo criterio se aplica a la reunión de socios de una S.R.L., que, contrariamente a una creencia extendida entre los estudiantes, también está obligada a llevar libros rubricados.

\section{Liquidación, cancelación y reconducción societaria}
\label{sec:reconduccion}

Para determinar cuándo termina de existir una sociedad se utilizó una analogía con la vida y la muerte de las personas. La \textbf{liquidación} equivale al «estado de coma»: la sociedad todavía existe jurídicamente, pero solo para concluir las operaciones pendientes. La \textbf{cancelación} es la muerte definitiva. Entre ambos momentos la sociedad puede seguir viva durante un período prolongado: la liquidación puede durar décadas.

\begin{definition}{Reconducción societaria}
Posibilidad de «revivir» una sociedad en liquidación que todavía no fue cancelada, siempre que se cuente con la unanimidad exigida por la ley.
\end{definition}

Por ejemplo, si una liquidación comenzó hace tiempo y la sociedad todavía no fue cancelada, los herederos de quienes eran en su momento los accionistas pueden decidir reconducirla, con unanimidad y los demás requisitos legales, precisamente porque la sociedad sigue existiendo.

\begin{example}{La marca de los sombreros y la reconducción societaria}
Una sociedad fabricante de sombreros dejó de ser rentable porque el producto perdió demanda. Aunque el negocio original ya no tiene sentido, conserva un activo valioso: su nombre y su marca registrada. Alguien se acerca porque quiere usar esa marca para quien quiera fabricar, por ejemplo, cabezas de plástico. El titular no quiere transferir la marca, porque era de su familia.

En lugar de liquidar definitivamente la sociedad y perder el valor de la marca familiar, opta por reconducirla y celebrar un convenio de explotación conjunta con un tercero interesado en el nuevo uso de la marca, manteniendo viva la sociedad y aprovechando el activo intangible que ya tenía.
\end{example}

Esta sección se conecta con los derechos reales: una sociedad en liquidación (pero no cancelada) puede seguir siendo titular de dominio sobre inmuebles de su patrimonio, lo que condiciona las opciones de quien pretende regularizar la titularidad de esos bienes (sección~\ref{sec:prescripcion}).

\cornell{La sociedad en formación y la adquisición de inmuebles}{
\cq{¿Cuándo nace la sociedad y qué efecto tiene la inscripción?} & Con el contrato constitutivo (contrato plurilateral de organización); la inscripción registral es meramente declarativa. \\ \hline
\cq{¿Qué es la sociedad en formación?} & La que se encuentra entre la firma del contrato constitutivo y la inscripción registral. \\ \hline
\cq{¿Capacidad y responsabilidad son lo mismo?} & No. La sociedad en formación tiene plena capacidad; la responsabilidad puede recaer sobre la sociedad, sobre quien actuó o sobre ambos, según el acto. \\ \hline
\cq{¿Qué riesgos tiene la compra en gestión de negocios?} & Tributación personal del comprador, sucesión si fallece, levantamiento de inhibiciones o embargos, y conflictos judiciales si la sociedad ni siquiera existía. \\ \hline
\cq{¿Cómo se evitan esos riesgos?} & Verificando que el acto esté dentro del objeto social o autorizándolo expresamente en el acta constitutiva, y fijando el plazo de duración desde la constitución. \\ \hline
\cq{¿Quién responde por un acto extraño al objeto?} & Si lo aprobaron directorio y asamblea: la sociedad, los directores que aprobaron y quien lo realizó; si no lo aprobó la asamblea: solo los directores que aprobaron y quien lo realizó. \\ \hline
\cq{¿Qué es la reconducción?} & Revivir, con unanimidad, una sociedad en liquidación aún no cancelada. \\ \hline
}{La sociedad existe desde el contrato constitutivo y tiene plena capacidad aun antes de inscribirse; lo que cambia es quién responde. Un buen encuadre en el objeto social o una autorización expresa evita los riesgos de la compra en gestión, y la reconducción permite conservar activos de una sociedad en liquidación.}

% ============================================================
% PARTE III — EL DOMINIO
% ============================================================
\part{El dominio y sus modalidades}

% ------------------------------------------------------------
\chapter{El dominio: un derecho real absoluto con límites}
\label{cap:dominio}

El estudio del dominio es el punto de partida para comprender procesos más complejos de organización patrimonial, incluida, en última instancia, la tokenización de activos. Este capítulo presenta el dominio como el derecho real más fuerte del sistema y muestra luego que ese carácter absoluto está atravesado por excepciones y modalidades que lo matizan considerablemente.

\section{Por qué comenzar por el dominio}
\label{sec:introduccion}

Los contenidos se adaptan a los intereses del grupo, pero se orientan a un objetivo final común: la \textbf{organización patrimonial} y, en particular, la \textbf{tokenización de activos}. Para llegar a él es necesario recorrer temas notariales y de derecho civil, porque sin ellos no se comprenderían la importancia y el porqué de la tokenización. No se trata de un desvío del programa, sino de un tránsito necesario: firma digital, blockchain y tokenización forman parte del recorrido, pero antes es preciso tener fijos otros conceptos. Para comprender las ventajas de la tokenización hay que entender en profundidad qué es lo «no tokenizado».

El método responde a una lógica pedagógica específica: se estudia primero el dominio pleno para que, a partir de él, puedan comprenderse los derechos reales sobre cosa ajena, y esa progresión permite proyectar las nuevas figuras vinculadas a la tecnología. La estructura del curso responde, en simultáneo, a dos frentes: lo rígido (la estructura clásica de los derechos reales tal como está regulada) y lo tecnológico (las nuevas posibilidades que la tecnología ofrece para los mismos institutos). Como todavía no se sabe con certeza hacia dónde evolucionará la implementación tecnológica, es necesario enseñar ambos caminos: el analógico, porque será el de la práctica profesional actual, y el tecnológico, porque representa el futuro de la disciplina.

Una evaluación diagnóstica previa mostró carencias importantes en derecho societario, por lo que se refuerzan tanto el derecho societario (Parte~\ref{parte:sociedades}) como los derechos reales antes de avanzar.

Los derechos reales se eligen como eje articulador no porque la organización patrimonial recaiga exclusivamente sobre ellos (puede recaer, por ejemplo, sobre una fábrica en funcionamiento), sino porque el concepto de derecho real permite advertir que instituciones conocidas como fijas, inamovibles y absolutas admiten en la práctica numerosas variantes y matices. Esas variantes permiten dar respuestas a los clientes sin recurrir siempre a los contratos típicos.

\begin{note}
El eje que atraviesa todo el recorrido es el derecho real, dado el \emph{numerus clausus} establecido por el Código Civil y Comercial (CCyC): todo lo que no sea uno de los derechos reales taxativamente enumerados por la ley será, a lo sumo, un contrato nominado, pero nunca un derecho real.
\end{note}

\begin{important}
El dominio no se limita a los bienes inmuebles. La misma lógica se aplica al derecho societario, en particular a la posibilidad de constituir usufructo sobre cuotas sociales o acciones, algo que la ley de sociedades vigente permite en Argentina, aunque todavía existen algunas figuras que el derecho argentino no admite en ese terreno (capítulo~\ref{cap:empresafamiliar}).
\end{important}

\section{El dominio como derecho real absoluto}
\label{sec:dominio-absoluto}

\begin{definition}{Dominio}
Poder jurídico de estructura legal que otorga al titular las facultades de \textbf{uso}, \textbf{goce} y \textbf{disposición} sobre la cosa.
\end{definition}

Una característica estructural clave es que el \textbf{dominio pleno}, entendido como la titularidad del cien por ciento de esas facultades, no puede recaer sobre más de una persona a la vez. Cuando el derecho se reparte entre varios cotitulares, como en el \textbf{condominio}, ya no se está frente a un dominio absoluto en cabeza de un único titular sino frente a otra figura (véase el capítulo~\ref{cap:condominio}).

\subsection{¿Se pierde el dominio si no se usa?}

En principio, no. El dominio, como derecho real \textbf{perpetuo}, subsiste aunque el titular no ejerza activamente sus facultades: si no usa el bien, sigue siendo dueño. Sin embargo, la perpetuidad admite excepciones, la primera y más relevante de las cuales es la usucapión.

\begin{note}
Uso, goce y disposición no son garantía de dominio eterno. La tensión central se resume en la expresión «nada es absoluto, nada es tan perfecto»: el dominio se presenta primero como el derecho real más fuerte del sistema, para luego revelarse sometido a numerosas excepciones legales y a modalidades (como el dominio revocable) que lo matizan.
\end{note}

\section{Excepciones a la perpetuidad del dominio}
\label{sec:excepciones}

La \textbf{usucapión} (prescripción adquisitiva) es la primera gran excepción. Quien no ejerce actos posesorios sobre la cosa puede ver desplazada su posición frente a quien sí los ejerce con ánimo de dueño durante el plazo legal: la usucapión da preferencia a quien estuvo en el uso, el goce y la posesión, y puede desplazar al titular registral.

Un punto técnico central es que la \textbf{calidad jurídica} en que se ocupa la cosa es determinante. Un locatario, por más que ocupe el inmueble durante décadas, nunca podrá usucapir, porque su posesión reconoce el derecho ajeno y le falta el ánimo de dueño exigido por la ley.

\begin{example}{El locatario que nunca puede usucapir}
Quien celebra veinte contratos de locación durante cuarenta años en forma ininterrumpida nunca logrará la prescripción adquisitiva, porque en cada contrato reconoce que el otro tiene la calidad de dueño.
\end{example}

Esa distinción, la calidad jurídica en que se posee, es lo que en la práctica permite «defender al que está poseyendo» o «defender al que es dueño del inmueble», según el caso.

\subsection{Causales de pérdida del dominio perfecto}

A partir de una consigna abierta (identificar las excepciones a la perpetuidad) se construyó el siguiente listado de causales, cada una con sus matices.

\begin{table}[H]
\centering
\caption{Causales de pérdida del dominio perfecto}
\renewcommand{\arraystretch}{1.2}
\begin{tabularx}{\textwidth}{>{\raggedright\arraybackslash}p{0.27\textwidth}>{\raggedright\arraybackslash}X}
\toprule
\textbf{Causal} & \textbf{Explicación} \\
\midrule
Usucapión (prescripción adquisitiva) & Pérdida por falta de ejercicio de actos posesorios, sumada al ejercicio de la posesión con ánimo de dueño por un tercero durante el plazo legal. \\
Expropiación & Puede ser \textbf{directa} (el Estado, la Nación, la provincia o el municipio expropian el bien) o \textbf{inversa} (el propio titular solicita que se lo expropie por determinada causa). \\
Destrucción total de la cosa & Solo la destrucción \emph{total} hace perder el dominio; si es parcial, el dominio subsiste sobre lo que resta. \\
Abandono (cosas muebles) & Aplica principalmente a cosas muebles; en inmuebles, el abandono por sí solo no produce la pérdida del dominio si nadie más ocupa la cosa y ejerce actos posesorios. \\
Ejecución forzosa & Derivada, por ejemplo, de una sentencia judicial que ordena el remate del bien para cobrar una deuda. \\
Confiscación & Asimilable, en inmuebles, a una forma de expropiación; sobre cosas muebles opera de manera más directa. \\
Accesión & Se da cuando una cosa se une a otra y quedan unidas; no siempre implica pérdida total del dominio, muchas veces deriva en un condominio. \\
\bottomrule
\end{tabularx}
\end{table}

\subsection{El abandono de inmuebles y los actos posesorios}

En inmuebles suele perderse el \emph{ejercicio} de la posesión (por ejemplo, dejando de pagar impuestos, uno de los actos posesorios típicos) sin que ello derive automáticamente en la prescripción adquisitiva a favor de un tercero, si nadie ocupó efectivamente el inmueble. Puede haber abandono sin que nadie tome la cosa: no todo dejar de poseer conduce a la prescripción.

\begin{example}{Terrenos comprados por boleto y nunca escriturados}
Una persona compraba terrenos (en la provincia de Córdoba, entre otros lugares) mediante una libreta de pagos, sin que la familia llegara a conocer con exactitud dónde estaban ubicados. Esos terrenos no tenían dominio a nombre del comprador porque no habían sido escriturados: el titular de dominio los había vendido por boleto de compraventa pero nunca había hecho la transferencia.

En estos casos el vendedor original (muchas veces una sociedad ya disuelta, sin actividad ni liquidador) conserva formalmente la titularidad registral, aunque el uso y goce reales estén en cabeza de otra persona. La posibilidad de usucapir dependerá de que haya existido posesión continuada con ánimo de dueño durante el plazo legal.
\end{example}

\section{Encuadrar antes de resolver: casos de pérdida y falta de dominio}
\label{sec:casos-practicos}

Dos aportes espontáneos de una estudiante (la donación con cargo y el caso de las bóvedas del cementerio de la Recoleta) permitieron remarcar la importancia de encuadrar correctamente cada problema jurídico.

\begin{example}{Las bóvedas del cementerio de la Recoleta}
Caso de conocimiento público: el gobierno de la Ciudad de Buenos Aires intentó reasignar bóvedas del cementerio de la Recoleta, y una heredera se opuso alegando que allí se encontraban sepultados sus familiares. Según lo relatado, el gobierno habría argumentado la falta de pago de la tasa correspondiente durante varios años.
\end{example}

Una advertencia metodológica central es \textbf{no mezclar categorías jurídicas distintas}. Una donación con cargo, un régimen especial de dominio (como el de los cementerios privados, sujeto a un régimen diferente del dominio público y del dominio privado ordinario) y la usucapión son instituciones distintas, y confundirlas lleva a errores de encuadre que perjudican al cliente.

\begin{important}
Cuando una persona trae un problema, lo primordial es saber dónde encuadrarlo. Primero se encuadra el problema; recién después se busca la solución. Al tratarse de derechos reales, no corresponde mezclar el análisis con instituciones ajenas (por ejemplo, la donación con cargo).
\end{important}

También se mencionaron los \textbf{conjuntos inmobiliarios} (countries) como ejemplo de regímenes especiales de dominio, regulados de forma diferente al dominio ordinario y a la propiedad horizontal: «dominio» no es una categoría única y homogénea, sino que admite regímenes especiales según el tipo de bien.

\section{Accesión: del dominio al condominio}
\label{sec:accesion}

La \textbf{accesión} es la causal de pérdida del dominio menos «absoluta». Consiste en la unión de una cosa con otra, de modo que quedan juntas. No toda accesión provoca la pérdida de dominio: muchas veces genera un \textbf{condominio} entre ambos titulares.

Después de repasar usucapión, expropiación, destrucción total, abandono, ejecución forzosa, confiscación y accesión, queda claro que el dominio, presentado como el derecho real «más duro» y «más absoluto», está atravesado por numerosas excepciones legales: aquel derecho que se decía perpetuo y absoluto lleva, en definitiva, un signo de interrogación.

\section{Formas del dominio: perfecto e imperfecto}
\label{sec:formas}

El dominio puede ser \textbf{perfecto} o \textbf{imperfecto}. Las causales de pérdida estudiadas corresponden al dominio perfecto. El dominio imperfecto, en cambio, no es pleno, y admite dos categorías:

\begin{itemize}
    \item \textbf{Dominio revocable}: el sujeto a una condición o a un plazo resolutorio (capítulo~\ref{cap:revocable}).
    \item \textbf{Dominio desmembrado}: aquel en el que alguna de las facultades que integran el dominio (uso, goce y disposición) se transfiere a otra persona.
\end{itemize}

\begin{definition}{Dominio desmembrado}
El dominio se integra por tres facultades: uso, goce y disposición. Existe desmembramiento cuando alguna de esas facultades se traslada a otro sujeto, lo cual puede ocurrir mediante la constitución de otro derecho real, aunque no siempre el desmembramiento implica necesariamente esa constitución.
\end{definition}

El dominio desmembrado conduce a los derechos reales sobre cosa ajena, cuyo estudio comienza en la Parte~\ref{parte:cosaajena} con el usufructo: si el dominio se integra por uso, goce y disposición, y el usufructo otorga solamente uso y goce, la disposición permanece en cabeza del titular de dominio, a quien se denomina \textbf{nudo propietario}.

\section{Regularizar una titularidad: usucapión y alternativas}
\label{sec:prescripcion}

Si el titular registral de un inmueble es una sociedad en liquidación pero todavía no cancelada (sección~\ref{sec:reconduccion}), el inmueble sigue integrando su patrimonio. Frente a esta situación hay dos caminos posibles: pedir la \textbf{liquidación judicial}, si la sociedad no tiene liquidador, o evaluar la vía de la \textbf{usucapión}. Como primer paso frente a un cliente con una posesión sin título perfecto, se recomienda indagar si existe una vía administrativa o societaria más rápida antes de embarcarse en un juicio de usucapión (más largo) o en una acción de reivindicación (más difícil).

\begin{important}
Antes de iniciar un juicio de usucapión conviene verificar si la sociedad titular todavía está vigente o si existe alguien, como un liquidador, que pueda firmar la escritura. Si esa vía existe, es mucho más fácil que enviar al cliente a usucapir: no se debe mandar directamente a la gente a iniciar un juicio de usucapión si hay otras posibilidades.
\end{important}

\begin{example}{Un juicio de usucapión que duró más de treinta años}
Caso real de la práctica profesional que ilustra los costos económicos, emocionales y de tiempo de un juicio de usucapión mal evaluado desde el inicio. Se trataba de una prescripción adquisitiva sobre un inmueble bajo propiedad horizontal, iniciada en 1983 y concluida recién en el año en que se dictó la clase. El inmueble fue escriturado al terminar el juicio y de inmediato vendido en 50.000 dólares.

La demora se debió, en gran parte, a decisiones procesales del juzgado: en lugar de avanzar de manera más expeditiva ante la falta de comparecencia de los demandados (en rebeldía), el juez tramitó individualmente los expedientes sucesorios de los ocho hermanos titulares originales, varios de los cuales fallecieron durante el proceso, generando nuevas sucesiones de sus propios herederos (ocho hijos, otros siete hijos). A ello se sumaron los costos de intereses, escrituras, juicio y abogados, con dos abogados fallecidos durante el proceso.
\end{example}

Este caso introduce una postura sobre la resolución de conflictos que se desarrolla en el capítulo~\ref{cap:ejercicio}: conviene evaluar alternativas antes de judicializar los temas.

\cornell{El dominio y sus límites}{
\cq{¿Qué es el dominio?} & Poder jurídico de estructura legal con facultades de uso, goce y disposición; el dominio pleno no puede recaer sobre más de una persona a la vez. \\ \hline
\cq{¿Por qué se estudia primero el dominio?} & Para comprender luego los derechos reales sobre cosa ajena y, en última instancia, la tokenización: hay que entender qué es lo «no tokenizado». \\ \hline
\cq{¿Qué implica el \emph{numerus clausus}?} & Lo que no sea un derecho real enumerado por la ley será, a lo sumo, un contrato nominado, nunca un derecho real. \\ \hline
\cq{¿Se pierde el dominio por no usarlo?} & En principio no (es perpetuo), salvo excepciones como la usucapión. \\ \hline
\cq{¿Por qué un locatario nunca puede usucapir?} & Porque reconoce el dominio ajeno y carece de ánimo de dueño: la calidad jurídica de la posesión es determinante. \\ \hline
\cq{¿Qué causales de pérdida del dominio perfecto se estudiaron?} & Usucapión, expropiación (directa o inversa), destrucción total, abandono, ejecución forzosa, confiscación y accesión. \\ \hline
\cq{¿Qué es dominio imperfecto?} & El no pleno: revocable (condición o plazo resolutorio) o desmembrado (una facultad se traslada a otro). \\ \hline
\cq{¿Qué conviene hacer antes de iniciar una usucapión?} & Verificar si la sociedad titular está vigente o si hay un liquidador que pueda escriturar. \\ \hline
}{El dominio se presenta como el derecho real más fuerte, pero admite excepciones y modalidades. Antes de resolver conviene encuadrar el problema y explorar vías más rápidas que el juicio; el dominio imperfecto (revocable o desmembrado) abre el camino a herramientas prácticas.}

% ------------------------------------------------------------
\chapter{El dominio revocable}
\label{cap:revocable}

El dominio revocable es la principal manifestación del dominio imperfecto y una de las herramientas más relevantes para la práctica notarial. Este capítulo examina su régimen legal, la diferencia entre plazos ciertos y condiciones inciertas, y dos aplicaciones concretas: el derecho de reversión en las donaciones y la garantía del cumplimiento de obligaciones sin ejecutar una hipoteca.

\section{Marco normativo}
\label{sec:marco-revocable}

Interesa observar las modalidades que puede tener el dominio y cómo puede usarse el dominio imperfecto. La estructura básica del instituto es que el dominio revocable es un dominio sometido a una condición, que puede ser suspensiva o resolutoria; puede tratarse de un tiempo fijo que transcurre o de un hecho puntual que ocurre.

\begin{definition}{Dominio revocable (art.~1965 CCyC)}
Dominio sometido a condición o plazo resolutorios, a cuyo cumplimiento el dueño debe restituir la cosa a quien se la transmitió.
\end{definition}

\begin{formula}{Artículo 1965 CCyC --- Dominio revocable}
\textit{Dominio revocable es el sometido a condición o plazo resolutorios, a cuyo cumplimiento el dueño debe restituir la cosa a quien se la transmitió. La condición o el plazo deben ser impuestos por disposición voluntaria expresa o por la ley. Las condiciones resolutorias impuestas al dominio se deben tener por no escritas si el efecto es sujetar la revocación a la sola voluntad del transmitente, salvo pacto en contrario. Las condiciones resolutorias impuestas al dominio se deben tener limitadas al término de diez años, aunque no pueda realizarse el hecho previsto dentro de aquel plazo, o este sea mayor o incierto. Si los diez años transcurren sin haberse producido la resolución, el dominio debe quedar definitivamente establecido. El plazo se computa desde la fecha del título constitutivo del dominio imperfecto.}

\smallskip
{\small Código Civil y Comercial de la Nación, Ley 26.994, Libro Cuarto --- Derechos Reales, Título III --- Dominio, Capítulo 5 (Dominio revocable).}
\end{formula}

El contenido normativo esencial es: condición o plazo resolutorios, límite de diez años y consolidación definitiva del dominio transcurrido ese plazo.

\section{Plazos ciertos y condiciones inciertas}
\label{sec:plazos}

Una de las advertencias técnicas más importantes es la diferencia entre un \textbf{plazo cierto} y una \textbf{condición incierta}.

\begin{example}{Dar la casa «cuando se reciba de abogado»}
Se propone dar una casa a alguien «cuando se reciba de abogado». La condición es posible, pero resulta \emph{incierta}. El problema concreto surge al combinarla con el límite legal de diez años del art.~1965: si la persona se recibe este año, el año próximo o dentro de veinticinco años, a los diez años el dominio queda consolidado. Si se recibió en el año diez y un día, se perdió la posibilidad de dar la casa en el tiempo que se quería.
\end{example}

\begin{formula}{Redacción de cláusulas resolutorias}
Conviene redactar cláusulas resolutorias con \textbf{plazos ciertos} y evitar condiciones basadas en hechos futuros de ocurrencia incierta.
\end{formula}

Fórmulas como «para cuando mi hija se case» o «para cuando mi hijo tenga un hijo» son riesgosas, porque no se sabe si el hecho ocurrirá alguna vez. En cambio, «para cuando mi hijo tenga la mayoría de edad» constituye un plazo cierto, porque en algún momento llegará a ella; y «para cuando mi hija cumpla 15 años» también, porque ya se sabe cuándo los cumplirá.

\begin{table}[H]
\centering
\caption{Plazo cierto y condición incierta en las cláusulas resolutorias}
\renewcommand{\arraystretch}{1.2}
\begin{tabularx}{\textwidth}{>{\raggedright\arraybackslash}p{0.17\textwidth}>{\raggedright\arraybackslash}p{0.31\textwidth}X}
\toprule
\textbf{Tipo de cláusula} & \textbf{Ejemplo} & \textbf{Riesgo jurídico} \\
\midrule
Plazo cierto & «Cuando mi hijo cumpla la mayoría de edad» / «cuando mi hija cumpla 15 años» & Bajo: el hecho ocurrirá necesariamente y en una fecha determinable de antemano. \\
Condición incierta & «Cuando se reciba de abogado» / «cuando se case» / «cuando tenga un hijo» & Alto: el hecho puede no ocurrir nunca y, aun si ocurre, puede hacerlo después de vencido el plazo máximo legal de 10 años. \\
\bottomrule
\end{tabularx}
\end{table}

\begin{important}
Aun dentro de un plazo cierto, conviene fijar \textbf{hitos intermedios verificables} (por ejemplo, presentación de planos, aprobación de observaciones), para reducir el margen de discusión sobre si la condición se cumplió. Se deben poner causales ciertas, que puedan verificarse y no den lugar a interpretación: si se pone simplemente «iniciar la obra», pueden objetar que la obra no se inició porque no aprobaron el plano.
\end{important}

\section{Aplicaciones prácticas del dominio revocable}
\label{sec:aplicaciones}

\subsection{Derecho de reversión y derecho de acrecer en las donaciones}
\label{sec:reversion}

El derecho de acrecer entre cónyuges muchas veces se aplica mal: cuando alguien es titular único de un inmueble, el cónyuge no acrece; solo acrece el titular de ese derecho real y, después, el cónyuge del titular de ese derecho.

\begin{example}{Derecho de reversión en una donación a dos hijos}
Un señor dona a favor de sus dos hijos y uno de ellos fallece previamente. Se da entonces el derecho de reversión: el inmueble vuelve al donante.

Una precisión que muchas veces se pasa por alto es que el derecho de reversión, al derivar de un dominio revocable, está sujeto al límite de diez años del art.~1965. Transcurrido ese plazo, la condición pierde eficacia, aunque el hecho al que estaba sujeta finalmente ocurra. Mucha gente piensa que el derecho de reversión es \emph{sine die}, pero el plazo legal lo limita.
\end{example}

\begin{definition}{Derecho de reversión}
Cláusula que puede incorporarse a una donación y que permite al donante recuperar el bien donado si el donatario fallece antes que él, siempre que ese fallecimiento ocurra dentro del plazo de diez años contados desde la donación. Vencido ese plazo, la reversión deja de ser exigible.
\end{definition}

\begin{note}
\textbf{Derecho de acrecer y derecho de reversión no son lo mismo.} En las donaciones a varios beneficiarios, el \textbf{derecho de acrecer} opera cuando uno de los donatarios no puede o no quiere recibir su parte, permitiendo que los restantes «acrezcan» su porción. El \textbf{derecho de reversión} opera cuando fallece uno de los donatarios antes de vencido el plazo legal, y hace que el bien vuelva al donante originario. Ambos deben pactarse expresamente y están sujetos al límite de diez años del dominio revocable cuando corresponda.
\end{note}

\begin{warning}
Es riesgoso pactar solamente el derecho de reversión sin reservarse además el usufructo del bien: si el donante fallece antes de que se cumplan los diez años, pierde toda posibilidad de usar la cosa durante ese lapso, ya que se desprendió tanto de la nuda propiedad como del uso. Véase la sección~\ref{sec:donacion}.
\end{warning}

\subsection{Garantizar el cumplimiento de una obligación sin ejecutar una hipoteca}

La aplicación considerada más interesante consiste en usar el dominio revocable como garantía de que el comprador de un inmueble inicie efectivamente una obra determinada dentro de un plazo pactado, sin constituir y eventualmente ejecutar una hipoteca. La pregunta que motiva el mecanismo es qué ocurre si se transmite el dominio de un inmueble a un tercero para la construcción de un edificio y este no comienza la obra.

\begin{example}{Venta a una desarrolladora sujeta a condición resolutoria}
En la compraventa de un inmueble a una desarrolladora se pacta que, si dentro de un plazo determinado (por ejemplo, tres años) no se inicia la obra proyectada, el dominio se revoca automáticamente y vuelve al vendedor, sin necesidad de tramitar ninguna acción judicial adicional. No hay que ejecutar la hipoteca y se asegura que el comprador tenga que empezar la obra sí o sí, porque si no la empieza no conserva el dominio.
\end{example}

La principal ventaja operativa frente a una garantía hipotecaria tradicional es que, producido el incumplimiento de la condición, el dominio se revoca de pleno derecho, sin acción judicial de ejecución.

\subsection{Instrumentación registral y ejecución de la revocación}

En el Registro de la Propiedad Inmueble, la anotación consigna directamente la \textbf{causal de dominio revocable}, del mismo modo que se anota el derecho de reversión. Para acreditar en la práctica el cumplimiento de la condición resolutoria, una solución instrumental habitual es otorgar, en la misma escritura que constituye el dominio revocable, un \textbf{poder especial irrevocable} que habilite al beneficiario de la reversión a otorgar por sí solo la escritura de revocación llegado el caso, sin nueva intervención del titular revocado.

\begin{example}{Cláusula de dominio revocable para asegurar el inicio de obra}
Elementos del caso, útiles como guía de redacción en la práctica notarial:
\begin{itemize}
    \item transmisión del dominio de un inmueble a un tercero, sujeta a condición resolutoria: el inicio efectivo de una obra determinada;
    \item plazo cierto para el cumplimiento de la condición (en el ejemplo, tres años);
    \item hitos intermedios verificables dentro de ese plazo (por ejemplo, presentación de planos dentro de determinados meses, respuesta a observaciones dentro de otro plazo);
    \item otorgamiento, en el mismo acto, de un poder especial irrevocable a favor del transmitente original, que le permita instrumentar por sí solo la escritura de revocación en caso de incumplimiento; y
    \item anotación registral de la causal de dominio revocable, de forma análoga a como se anota el derecho de reversión.
\end{itemize}
\end{example}

\cornell{El dominio revocable}{
\cq{¿Qué es el dominio revocable?} & Dominio sometido a condición o plazo resolutorios, a cuyo cumplimiento el dueño debe restituir la cosa a quien se la transmitió (art.~1965 CCyC). \\ \hline
\cq{¿Qué límite temporal tiene?} & Diez años desde el título constitutivo; transcurridos sin resolverse, el dominio queda definitivamente establecido. \\ \hline
\cq{¿Por qué preferir plazos ciertos?} & Las condiciones inciertas pueden no ocurrir o ocurrir después de los diez años; los hitos verificables reducen la discusión sobre su cumplimiento. \\ \hline
\cq{¿Diferencia entre acrecer y reversión?} & Acrecer: los demás donatarios amplían su parte si uno no puede o no quiere recibirla. Reversión: el bien vuelve al donante si fallece el donatario dentro del plazo. \\ \hline
\cq{¿Cómo garantiza el dominio revocable una obra sin hipoteca?} & Se condiciona el dominio al inicio de la obra en plazo cierto; si no se inicia, se revoca de pleno derecho, con poder especial irrevocable para instrumentarlo y anotación registral. \\ \hline
}{El dominio revocable permite moldear el dominio con condiciones o plazos, siempre dentro del límite legal de diez años. Su eficacia práctica depende de redactar plazos ciertos y de instrumentar la revocación con poder especial irrevocable y anotación registral.}

% ============================================================
% PARTE IV — DERECHOS REALES SOBRE COSA AJENA
% ============================================================
\part{Derechos reales sobre cosa ajena}
\label{parte:cosaajena}

% ------------------------------------------------------------
\chapter{El usufructo: concepto, constitución y régimen}
\label{cap:usufructo}

El usufructo es, junto con la hipoteca y la servidumbre, uno de los derechos reales de mayor aplicación práctica y, sin embargo, uno de los más subutilizados por los profesionales del derecho. Dominar su lógica (quién tiene el uso y el goce, quién conserva la disposición, cómo se constituye, cómo se transmite y cómo se extingue) permite resolver problemas que exceden el aula: planificar la sucesión de un patrimonio familiar sin perder el control de los bienes en vida, estructurar la sucesión de una empresa familiar sin que el fundador pierda el manejo de las decisiones estratégicas, diseñar operaciones inmobiliarias que beneficien tanto a quien vende como a quien compra, y asesorar sobre las consecuencias impositivas de cada decisión.

El usufructo se entrelaza con el derecho de familia (bienes gananciales, asentimiento conyugal), con el derecho societario (usufructo de acciones, derechos políticos y económicos), con el derecho sucesorio (planificación patrimonial, donaciones con reserva de usufructo) y con el derecho tributario (impuesto a la transmisión gratuita de bienes). También es útil a título personal: quien done, herede, compre o venda un inmueble, tenga bienes gananciales, integre una sociedad familiar o quiera proteger su vivienda en la vejez sin perder autonomía económica puede encontrarse frente a estas figuras.

Este capítulo presenta el concepto y el régimen general; el siguiente aborda su uso en la planificación patrimonial y sus consecuencias fiscales.

\section{Concepto: nuda propiedad y reserva de usufructo}

El derecho real que permite desdoblar el dominio en uso y goce, por un lado, y disposición, por otro, es el usufructo, que otorga a su titular el uso y el goce. Se trata de un derecho real que suele asociarse únicamente al esquema familiar clásico, sin advertir que admite otros usos de gran interés práctico.

\begin{definition}{Reserva de usufructo}
Cuando el único titular de un inmueble transmite la nuda propiedad y se reserva para sí el usufructo, se configura la \textbf{reserva de usufructo}: la titularidad de dominio (la propiedad «nuda», vacía de uso) se transmite a otra persona, mientras que el uso y goce permanece, por reserva, en cabeza de quien transmite.
\end{definition}

\begin{example}{El lápiz}
Si una persona tiene un lápiz y le da a otra la nuda propiedad, se queda para sí con el uso del lápiz. El dueño del lápiz pasa a ser el receptor de la nuda propiedad; quien transmite conserva el uso.
\end{example}

Una familia con un inmueble y una empresa familiar ilustra el conjunto: los padres pueden donar la nuda propiedad de la casa a los hijos y reservarse el derecho de seguir viviendo en ella; esa misma lógica puede trasladarse a la empresa, donando las acciones y reservándose el usufructo (capítulo~\ref{cap:empresafamiliar}); y el día en que la familia decide vender el inmueble o cancelar el usufructo, cada camino tiene un costo impositivo distinto (sección~\ref{sec:fiscal}).

\section{Bienes gananciales y titularidad registral}

Una pregunta central: ¿puede una persona casada, titular registral de un bien ganancial, reservarse el usufructo para sí y además constituirlo a favor de su cónyuge? La respuesta es negativa en cuanto a la reserva a favor del cónyuge: puede reservarse para sí, pero no reservar al cónyuge, porque el cónyuge es titular de dominio.

\begin{important}
Aunque un bien sea ganancial, la \textbf{titularidad registral} corresponde exclusivamente a quien figura como titular en la escritura. La gananciabilidad solo produce efectos patrimoniales (participación del 50\%) en dos momentos: la disolución de la sociedad conyugal o el fallecimiento. Mientras esas situaciones no se produzcan, la administración del bien corresponde en forma exclusiva al titular registral.
\end{important}

Si un matrimonio adquiere un bien durante la vigencia de la sociedad conyugal pero la escritura se otorga solo a nombre de uno de los cónyuges, no es posible subsanar esa situación mediante un contrato entre cónyuges para cambiar la titularidad: quien registró el bien a su nombre conserva la administración durante todo el matrimonio. Que un bien ganancial es «de los dos» es cierto desde lo patrimonial, pero no desde la realidad fáctica, porque la administración la tiene solamente el titular y solo recaerá sobre los dos cuando se parta.

\section{Constitución de usufructo entre cónyuges}

Descartada la reserva directa a favor del cónyuge, existen dos alternativas para lograr un resultado equivalente.

La primera es el cambio del régimen patrimonial matrimonial: puede optarse por mutar de comunidad de gananciales a separación de bienes, lo que habilitaría una adjudicación del 50\% a cada cónyuge. Se relativizó su utilidad práctica, dado que la mayoría de las personas no tiene interés en modificar todo su régimen patrimonial solo para otorgar el usufructo al cónyuge.

\begin{formula}{Doble operación simultánea entre cónyuges}
La segunda vía consiste en una doble operación simultánea: el titular registral transmite el dominio a la otra persona y, en el mismo acto, la otra persona le constituye el usufructo a quien transmitió. Aplicado a una cuota parte, el titular puede reservarse el usufructo sobre su 50\% y, en el mismo acto, el nuevo nudo propietario constituye el usufructo sobre el otro 50\% a favor de su cónyuge.

El resultado es que ambos cónyuges terminan siendo cousufructuarios en las proporciones acordadas, aun cuando uno solo conserve la nuda propiedad plena. Esta estructura resulta útil en operaciones con montos elevados de impuesto de sellos, donde la ingeniería de la operación puede incidir en la carga fiscal.
\end{formula}

\section{El asentimiento conyugal}
\label{sec:asentimiento}

Un supuesto complementario es la constitución de un usufructo a favor del cónyuge del nudo propietario, cuando quien dona no tiene cónyuge al donar pero lo adquiere después (por ejemplo, una donación a un hijo casado, sobre la que luego se pretende constituir usufructo a favor de la nuera o el yerno). Si el hijo (nudo propietario) está casado y recibió el bien por donación (bien propio), la constitución de un derecho real sobre él requiere, en principio, la conformidad del cónyuge, porque se está disponiendo del bien.

\begin{important}
\textbf{Criterio.} Cuando se constituye, transfiere o modifica un derecho real se requiere el asentimiento conyugal, con la excepción de los bienes propios que no constituyan la sede del hogar conyugal. Si el bien propio es la sede del hogar conyugal, el asentimiento se exige igualmente, para proteger la vivienda familiar.
\end{important}

En el propio acto de donación, si el bien resulta propio para el donatario por provenir de una liberalidad, suele requerirse igualmente que el cónyuge del donatario ratifique o reconozca ese carácter propio.

\begin{example}{Observación registral sobre una partición sucesoria}
Una escritura instrumentaba una adjudicación por partición derivada de una sucesión, sin que compareciera el cónyuge del adjudicatario a prestar el asentimiento. El escribano no lo exigió por entender que no se trataba de la sede del hogar conyugal y que el carácter propio del bien era indudable, por provenir de una sucesión. El acto fue observado por el organismo registral, que exigió el asentimiento. Se manifestó desacuerdo con esa observación: no era sede del hogar conyugal y era claro que se trataba de un bien propio.

Como criterio profesional, no se pide el asentimiento cuando el bien viene de bienes propios y es claro; pero si hay dudas sobre dónde viven (por ejemplo, si coincide el domicilio del acto con el del inmueble que se enajena), sí se lo pide.
\end{example}

La exigencia del asentimiento cede, así, frente a la certeza sobre el carácter propio del bien y sobre el hecho de que no constituye la sede del hogar conyugal; ante la duda sobre este último punto corresponde requerirlo, como resguardo de la vivienda familiar.

\section{El usufructo oneroso}
\label{sec:oneroso}

El usufructo puede adquirirse por reserva, por constitución y también por compra. Es allí donde se advierte la riqueza de este derecho real, subutilizado en la práctica.

\begin{example}{El propietario de 85 años y el inquilino de toda la vida}
Un propietario de 85 años que alquila desde hace diez años un local o departamento al mismo inquilino le ofrece venderle la nuda propiedad: para cuando el propietario fallezca, el inmueble queda para el inquilino. El propietario se reserva el usufructo, sigue cobrando el alquiler y le da al inquilino la tranquilidad de que el día en que fallezca el inmueble será suyo. El precio se reduce respecto del valor pleno del inmueble.

La quita responde a que el precio de la nuda propiedad se reduce en función de la expectativa de vida del usufructuario, de modo análogo a una renta vitalicia, pero cobrada en un pago único en lugar de cuotas periódicas.
\end{example}

\begin{example}{La inversión de un gran inversor sobre un terreno}
Un nudo propietario recibe la oferta de un inversor de gran escala interesado en desarrollar una inversión importante sobre un terreno. En lugar de un contrato de locación por diez años ---que no lo resguardaría frente al fallecimiento del propietario o a un eventual embargo---, el inversor le compra el usufructo del terreno a título oneroso. Así blinda su inversión frente a contingencias sucesorias o de embargos que podrían afectar a un simple contrato de locación.
\end{example}

\begin{formula}{Plazo del usufructo oneroso}
El usufructo adquirido a título oneroso puede pactarse hasta el fallecimiento del usufructuario o con un plazo cierto y determinado. Si el usufructuario es una persona humana, el usufructo se extingue con su fallecimiento, aun cuando se hubiera pactado un plazo fijo mayor.

En consecuencia, si se prevé adquirir un usufructo, conviene que el adquirente sea una persona jurídica con un plazo de duración mayor al tiempo pactado del usufructo: la extinción se producirá entonces por cancelación de esa sociedad y no por fallecimiento.
\end{formula}

También existe un mercado real para la compraventa exclusiva de la nuda propiedad: quien compra una nuda propiedad con un usufructuario de noventa años a un precio muy reducido hace una apuesta financiera razonable sobre la duración esperada del usufructo. Esta práctica se señaló como habitual en España, comparándosela con la hipoteca inversa.

\subsection{Cesión del usufructo y renta vitalicia}

El usufructuario puede ceder su derecho a un tercero. Quien adquiere el usufructo cedido queda sujeto al mismo límite temporal del usufructuario original: el usufructo cedido se extingue con el fallecimiento del usufructuario originario, no con el del cesionario.

\begin{definition}{Renta vitalicia}
Contrato aleatorio cuyo resultado económico depende de un hecho incierto: la duración de la vida del beneficiario. Implica una transacción en la que una parte transmite la propiedad de un bien a cambio de recibir una renta periódica durante toda su vida.
\end{definition}

\begin{important}
La renta vitalicia puede ser un excelente o un pésimo negocio según variables que no pueden conocerse de antemano: nadie sabe con certeza cuánto vivirá el rentista. Además, quien recibe la renta puede ver comprometida su capacidad de pago si cambian las circunstancias externas (por ejemplo, si aumentan los costos de un geriátrico que se comprometió a pagar). Toda decisión patrimonial sujeta a variables externas exige cautela y revisión constante.
\end{important}

\section{Extinción del usufructo y derecho de acrecer}
\label{sec:acrecer}

El usufructo se extingue con el fallecimiento del usufructuario: no hay sucesión, no puede testarse y no se incorporan los herederos. El punto se complejiza cuando hay pluralidad de usufructuarios, según exista o no derecho de acrecer entre ellos.

\begin{definition}{Derecho de acrecer entre cousufructuarios}
Si varios usufructuarios \textbf{no} tienen entre sí derecho de acrecer, el usufructo se extingue parcialmente a medida que cada uno fallece, y esa porción consolida el dominio pleno en cabeza del nudo propietario. Si \textbf{sí} se pactó derecho de acrecer, el o los usufructuarios sobrevivientes se quedan con la totalidad del usufructo.
\end{definition}

\begin{warning}
Es frecuente omitir, por error, el derecho de acrecer al constituir un usufructo a favor de varios cousufructuarios (por ejemplo, ambos cónyuges donantes que se reservan el usufructo). Si no se incluye y fallece uno, el nudo propietario adquiere el dominio pleno sobre la parte ideal correspondiente. Los cousufructuarios no necesariamente son cónyuges: pueden ser convivientes o personas sin vínculo afectivo, por lo que su inclusión debe decidirse expresamente en cada caso.
\end{warning}

\begin{example}{Acrecimiento proporcional entre tres usufructuarios}
Tres personas son cousufructuarias de un mismo bien: la primera titular del 50\% del usufructo y las otras dos del 25\% cada una. Al fallecer la primera, si existe derecho de acrecer, su 50\% no se reparte en partes iguales entre las restantes, sino a prorrata de la proporción que cada una ya tenía sobre el 50\% remanente. Así, si la segunda tenía el 30\% de ese 50\% restante y la tercera el 20\%, cada una acrece en esa misma proporción sobre el porcentaje liberado.
\end{example}

\cornell{El usufructo: concepto, constitución y régimen}{
\cq{¿Qué es la reserva de usufructo?} & La transmisión de la nuda propiedad por quien se reserva el usufructo: el dominio pasa a otro, pero el uso y goce queda en cabeza del transmitente. \\ \hline
\cq{¿Puede reservarse usufructo a favor del cónyuge sobre un bien ganancial?} & No: el titular solo puede reservarse para sí; el cónyuge no es titular registral. La gananciabilidad solo opera al disolverse la sociedad conyugal o al fallecer. \\ \hline
\cq{¿Cómo se logra el usufructo de ambos cónyuges?} & Cambiando al régimen de separación de bienes o mediante una doble operación simultánea: transmisión del dominio y constitución del usufructo en el mismo acto. \\ \hline
\cq{¿Cuándo se requiere asentimiento conyugal?} & Al constituir, transferir o modificar un derecho real, salvo bienes propios que no sean sede del hogar conyugal; ante duda sobre la sede, se pide. \\ \hline
\cq{¿Por qué el inversor prefiere comprar el usufructo a alquilar?} & Porque el usufructo es un derecho real que protege la inversión frente a fallecimiento o embargos que afectarían a una locación. \\ \hline
\cq{¿Cómo se extingue el usufructo de una persona humana?} & Con su fallecimiento, aun con plazo mayor; conviene un adquirente persona jurídica si se busca plazo largo. \\ \hline
\cq{¿Qué ocurre con varios cousufructuarios?} & Sin acrecer, cada porción consolida en el nudo propietario al fallecer su titular; con acrecer, los sobrevivientes acrecen a prorrata. \\ \hline
\cq{¿Qué riesgo tiene la renta vitalicia?} & Es aleatoria: depende de cuánto viva el rentista y de las circunstancias futuras. \\ \hline
}{El usufructo desdobla el dominio: el nudo propietario conserva la disposición y el usufructuario el uso y goce. Su constitución exige atender a la titularidad registral, al asentimiento conyugal y al derecho de acrecer, y su extinción con la muerte del usufructuario condiciona todo diseño de largo plazo.}

% ------------------------------------------------------------
\chapter{El usufructo en la planificación patrimonial y sus aspectos fiscales}
\label{cap:usufructoplan}

Una vez presentado el régimen general, este capítulo examina cómo se utiliza el usufructo para ordenar el patrimonio familiar, qué riesgos tienen las soluciones tradicionales (en especial la donación) y cuáles son las consecuencias fiscales de constituirlo, cancelarlo y valuarlo.

\section{Donación con reserva de usufructo}
\label{sec:donacion}

\begin{example}{El padre que dona la nuda propiedad y se reserva el usufructo}
Un padre transmite la nuda propiedad del inmueble a sus hijos, pero desea continuar viviendo en él. Se reserva el usufructo, de modo que la nuda propiedad pasa a los hijos mientras el padre puede seguir percibiendo los frutos del bien.
\end{example}

\begin{important}
Si el donante no se reserva el usufructo y el donatario fallece antes que él, el inmueble ingresa en la sucesión del donatario y el donante puede quedar sin ningún derecho sobre el bien. Sobre la combinación con el derecho de reversión, véase la sección~\ref{sec:reversion}.
\end{important}

\subsection{Riesgos de la donación como planificación sucesoria}

La práctica frecuente de donar bienes en vida para evitar la sucesión merece un análisis crítico. La donación puede ser una herramienta razonable cuando la persona tiene una edad avanzada o una situación de salud que hace previsible el desenlace, pero no como recurso preventivo general y anticipado. Entre sus riesgos:

\begin{itemize}
    \item \textbf{Exposición a contingencias del donatario:} una vez que el bien ingresa al patrimonio del hijo donatario queda expuesto a las contingencias propias de ese patrimonio. Por ejemplo, si el hijo atraviesa un divorcio, el bien puede verse afectado por una medida de no innovar sobre bienes propios, aun cuando el donante hubiera conservado el usufructo.
    \item \textbf{Limitación del acceso al crédito:} el hijo, al ser titular registral de un inmueble (aunque no lo use, porque el usufructo sigue en cabeza del donante), puede ver limitado su acceso al crédito bancario, porque a los ojos de la entidad financiera ya es propietario de una vivienda.
    \item \textbf{Falta de flexibilidad:} la donación es un acto difícil de revocar y carece de flexibilidad frente a cambios en la situación familiar o económica.
\end{itemize}

Antes de recurrir automáticamente a la donación, al usufructo o a otras figuras tradicionales, conviene evaluar si existen herramientas mejor adaptadas a la situación de cada familia.

\section{Transmisión onerosa del usufructo}

Cuando el titular de un usufructo gratuito desea que los nudos propietarios vendan el inmueble sin renunciar a su participación económica, puede transmitir en forma onerosa su propio usufructo de manera simultánea a la venta de la nuda propiedad.

\begin{example}{Venta simultánea de la nuda propiedad y del usufructo}
Un padre de 70 años se había reservado el usufructo gratuito de un inmueble al ceder a sus hijas la nuda propiedad. Ante la posibilidad de que ellas quieran vender, la estrategia consiste en que las hijas vendan la nuda propiedad y el padre transmita en forma onerosa su usufructo en el mismo acto: el padre recibe una parte del precio y las hijas la otra.

Alternativamente, el padre puede transferir en forma onerosa el usufructo del inmueble original y adquirir, también en forma onerosa, el usufructo sobre un nuevo inmueble adquirido por las hijas, para seguir contando con un lugar donde vivir mientras estas disponen del inmueble anterior.
\end{example}

\section{Usufructo vitalicio, legítima y colación: análisis de un fallo}
\label{sec:legitima}

\begin{example}{Usufructo vitalicio consentido por los hermanos}
Una mujer con cuatro hijos constituyó a favor del más pequeño un usufructo vitalicio sobre una propiedad, reservándose ella la nuda propiedad. Todos los hermanos firmaron consintiendo el acto. Transcurridos cuatro años la mujer falleció. Los hermanos pretendieron entonces dejar sin efecto el usufructo del hermano menor, invocando la colación y la afectación de la porción legítima. Ese hermano quedó luego incapacitado y actúa mediante representante.
\end{example}

El fallo de primera instancia, confirmado por la Cámara, rechazó el pedido por dos razones:

\begin{enumerate}
    \item \textbf{Teoría de los actos propios:} los hermanos habían consentido el acto, por lo que no podían impugnarlo después.
    \item \textbf{Validez del usufructo vitalicio:} había sido otorgado en forma vitalicia y válida, por lo que no podía revocarse unilateralmente; la nuda propiedad quedó en cabeza de los hermanos.
\end{enumerate}

Se señala que, más allá de la postura personal expresada sobre el caso, este ilustra una manera en que puede vulnerarse una norma de orden público del Código: la protección de la porción legítima de los herederos forzosos (sección~\ref{sec:autonomia}).

\begin{example}{Usufructo constituido a favor de cónyuge joven}
Un hombre de 90 años, con hijos de 60, contrae matrimonio con una mujer de 30 y le constituye el usufructo sobre su patrimonio. Los hijos conservan la nuda propiedad, pero es probable que la cónyuge los sobreviva por un extenso período, de modo que quedarían con una nuda propiedad vacía de contenido económico durante ese lapso.
\end{example}

\begin{important}
El CCyC no distingue en función de la edad de las partes ni de probabilidades de supervivencia al regular la constitución del usufructo. Sin perjuicio de ello, cuando se constituye un usufructo sobre la totalidad del patrimonio a favor de una persona con expectativa de vida considerablemente mayor a la de los herederos forzosos, puede vulnerarse en los hechos la protección de la porción legítima, en tanto los herederos quedan con una nuda propiedad inaprovechable durante un período prolongado.
\end{important}

\section{Cancelación del usufructo y aspectos impositivos}
\label{sec:fiscal}

Con el usufructuario con vida, el usufructo se cancela mediante su renuncia, lo que consolida el dominio pleno en cabeza del nudo propietario.

\begin{important}
\textbf{Referencia normativa: impuesto a la transmisión gratuita de bienes (Provincia de Buenos Aires).} La Provincia de Buenos Aires es, dentro del país, la única jurisdicción que mantiene vigente el impuesto a la transmisión gratuita de bienes, aplicable a los actos a título gratuito. Las causales que determinan su procedencia son la ubicación del bien en la Provincia o el domicilio del causante o del sucesor en dicha jurisdicción.
\end{important}

\begin{example}{Renuncia al usufructo: Ciudad de Buenos Aires frente a Provincia de Buenos Aires}
Una madre usufructuaria y su hijo nudo propietario deciden vender un inmueble.

Si el inmueble está en la Ciudad de Buenos Aires, la renuncia al usufructo no genera preocupación tributaria adicional, porque esa jurisdicción no aplica el impuesto a la transmisión gratuita de bienes: la madre renuncia gratuitamente, el hijo queda como titular del dominio pleno y puede vender libremente.

En la Provincia de Buenos Aires, en cambio, la renuncia gratuita configura, a criterio fiscal, un acto a título gratuito que beneficia patrimonialmente al nudo propietario, y tributa el impuesto a la transmisión gratuita de bienes.
\end{example}

\begin{formula}{Alternativa a la renuncia gratuita en Provincia de Buenos Aires}
En lugar de que la usufructuaria renuncie gratuitamente y luego el hijo venda el dominio pleno, puede estructurarse una doble venta a título oneroso al mismo comprador: el nudo propietario vende la nuda propiedad y el usufructuario vende, en el mismo acto, el usufructo. El efecto económico final es equivalente (el comprador termina con el dominio pleno) pero sin el acto gratuito que dispara el impuesto provincial.

% TODO: verificar consistencia entre fuentes originales. La fuente indica alícuotas de 1\% (usufructo) y 2,7\% (nuda propiedad) "a los fines de los honorarios notariales" en la Ciudad de Buenos Aires, y en otro pasaje una alícuota de sellos de 2,7\%.
En la Ciudad de Buenos Aires, según lo expuesto en el curso y a los fines de los honorarios notariales, la venta del usufructo tributa el 1\% y la venta de la nuda propiedad el 2,7\%.
\end{formula}

\begin{warning}
El uso de la diferencia de alícuotas entre la venta del usufructo y la de la nuda propiedad debe manejarse con cuidado: recurrir a esta estructura únicamente por motivos impositivos, para pagar menos impuesto, requiere especial prudencia profesional.
\end{warning}

\section{Valuación del usufructo a efectos impositivos}

Cuando el usufructo y la nuda propiedad se venden por separado al mismo comprador surge el problema de cómo distribuir el precio total entre ambos.

\begin{example}{Subvaluación indebida del usufructo}
Algunos profesionales, buscando reducir el impuesto de sellos (más bajo sobre la cuota del usufructo), asignan al usufructo un valor artificialmente bajo: por ejemplo, vender el usufructo en diez mil y la nuda propiedad en cincuenta mil, sin relación con sus valores reales. Esta práctica no debe aceptarse.

El fundamento es técnico: el valor del usufructo debe guardar proporción con la expectativa de vida ---y, por lo tanto, de uso y goce efectivo--- del usufructuario; no puede fijarse arbitrariamente. Quien resulta perjudicado es, fundamentalmente, el comprador, que termina pagando por la nuda propiedad un sellado superior al que correspondería si la proporción reflejara razonablemente la expectativa de vida del usufructuario.
\end{example}

\begin{formula}{Proporción máxima orientativa del usufructo}
Como criterio orientativo y sujeto a la edad del usufructuario, el valor asignado al usufructo no debería superar, como máximo, entre el 25\% y el 30\% del valor total de la operación. La proporción se calcula sobre el valor total de la venta y no sobre una porción aislada.
\end{formula}

\begin{exercise}{Distribución del precio en una operación con usufructo}
Sobre una operación de \$200.000, con una alícuota de sellos del 2,7\%, se planteó cómo distribuir el precio entre el usufructo y la nuda propiedad. Frente a la propuesta de calcular el 30\% sobre la porción del usufructo aisladamente, se corrigió el criterio: la proporción se calcula sobre el total de la operación, porque ambas ventas integran una misma operación económica y deben evaluarse conjuntamente para determinar si la distribución resulta razonable.
\end{exercise}

En síntesis, el manejo impositivo puede generar situaciones que afectan al usufructuario y al nudo propietario, pero sobre todo al comprador, que suele ser el más perjudicado.

\cornell{El usufructo en la planificación patrimonial y sus aspectos fiscales}{
\cq{¿Qué riesgo corre el donante que no se reserva el usufructo?} & Si el donatario fallece antes, el bien ingresa en su sucesión y el donante puede quedar sin derecho alguno. \\ \hline
\cq{¿Por qué la donación no siempre conviene?} & Expone el bien a contingencias del donatario (p.~ej. divorcio), limita su acceso al crédito y es poco flexible. \\ \hline
\cq{¿Cómo se resolvió el caso del usufructo vitalicio consentido?} & Se rechazó la impugnación por la teoría de los actos propios y por la validez del usufructo vitalicio. \\ \hline
\cq{¿Cómo puede afectarse la legítima con un usufructo?} & Un usufructo total a favor de alguien con mayor expectativa de vida deja a los herederos con una nuda propiedad inaprovechable. \\ \hline
\cq{¿Qué diferencia fiscal hay entre CABA y PBA al renunciar al usufructo?} & En PBA la renuncia gratuita tributa el impuesto a la transmisión gratuita; en CABA no. \\ \hline
\cq{¿Qué alternativa evita el acto gratuito?} & Doble venta onerosa simultánea (nuda propiedad y usufructo) al mismo comprador. \\ \hline
\cq{¿Qué proporción orientativa se asigna al usufructo?} & Como máximo entre el 25\% y el 30\% del valor total de la operación, según la edad del usufructuario. \\ \hline
}{El usufructo permite ordenar el patrimonio en vida, pero exige evaluar riesgos y costos: la donación puede exponer y limitar al donatario, el usufructo total puede afectar la legítima y la cancelación o la venta tienen efectos impositivos distintos según la jurisdicción. La valuación debe ser razonable y guardar relación con la expectativa de vida.}

% ------------------------------------------------------------
\chapter{Los derechos de uso y de habitación}
\label{cap:usohab}

Aunque parecen figuras poco utilizadas en la práctica, bien empleadas resultan sumamente útiles, sobre todo para garantizar vivienda a personas vulnerables con una solidez jurídica muy superior a la de los acuerdos personales. Este capítulo distingue el uso y la habitación entre sí y respecto del usufructo, y muestra sus aplicaciones y límites.

\section{Diferencias entre el uso y la habitación}
\label{sec:usohab}

\begin{table}[H]
\centering
\caption{Derecho de uso y derecho de habitación}
\small
\renewcommand{\arraystretch}{1.2}
\begin{tabularx}{\textwidth}{>{\raggedright\arraybackslash}p{0.20\textwidth}>{\raggedright\arraybackslash}X>{\raggedright\arraybackslash}X}
\toprule
\textbf{Criterio} & \textbf{Derecho de uso} & \textbf{Derecho de habitación} \\
\midrule
Titularidad & Puede constituirse a favor de personas humanas. & Solo puede constituirse a favor de personas humanas. \\
Objeto & Puede recaer sobre cosas muebles e inmuebles (por ejemplo, un tractor). & Solo puede recaer sobre un inmueble. \\
Amplitud & Derecho más amplio: además de usar, permite gozar de los frutos que la cosa produce. & Derecho más reducido, limitado a morar en el inmueble. \\
Vigencia & Puede ser vitalicio o pactarse por un plazo determinado. & Puede ser vitalicio o temporal. \\
Figura personal equivalente & Puede constituirse también como derecho personal (por ejemplo, mediante un contrato de comodato). & Solo se admite como derecho real; no como derecho personal equivalente. \\
\bottomrule
\end{tabularx}
\end{table}

\begin{note}
% TODO: verificar consistencia entre fuentes originales. Una fuente indica que el derecho de uso permite gozar de los frutos de la cosa; otra indica que el uso es solo para uso personal del titular y su familia, no puede arrendarse y no es transferible.
Las fuentes describen el alcance del derecho de uso con matices que conviene contrastar: una lo presenta como un derecho más amplio que la habitación, que incluye el goce de los frutos; otra lo describe, respecto de un inmueble, como limitado al uso personal del titular y su familia, sin posibilidad de arrendarlo ni de transferirlo. Ambas descripciones se conservan en este capítulo.
\end{note}

\subsection{Solidez jurídica de los derechos reales}

Un derecho real de habitación ya inscripto se sostiene por sí mismo, salvo que se configure alguna causal de extinción, sin necesidad de recurrir a los tribunales para hacerlo valer. Un derecho meramente personal (como un comodato), en cambio, depende de la autonomía de la voluntad contractual y puede ser discutido por terceros, obligando a litigar para hacerlo respetar.

\section{Uso, habitación y usufructo}

\begin{table}[H]
\centering
\caption{Habitación, uso y usufructo}
\renewcommand{\arraystretch}{1.2}
\begin{tabularx}{\textwidth}{>{\raggedright\arraybackslash}X c c c}
\toprule
\textbf{Aspecto} & \textbf{Habitación} & \textbf{Uso} & \textbf{Usufructo} \\
\midrule
¿Puede alquilar? & No & No (solo para sí) & Sí \\
¿Puede vivir con pareja? & No & Sí (familia) & Sí \\
¿Es transferible? & No & No & Sí \\
Robustez & Menor & Media & Mayor \\
\bottomrule
\end{tabularx}
\end{table}

Según lo expuesto, el derecho de habitación se orienta a asegurar techo a una persona vulnerable; el derecho de uso, al uso personal con la familia; y el usufructo, a la explotación económica del bien junto con su goce como vivienda. Respecto del uso frente al usufructo, el inmueble sujeto a uso se destina al uso personal del titular y su familia, no es transferible y su robustez es menor; el usufructo puede alquilarse y transferirse onerosamente y permite habitar al titular y a quien este decida.

\begin{example}{Derecho de uso para uso familiar estacional}
Una persona posee una vivienda en Mar del Plata que utiliza únicamente en temporada. Su hermano trabaja en esa ciudad durante el año, por lo que le otorga un derecho de uso para que pueda instalarse en esos períodos. El titular puede usar la vivienda para necesidades de su familia durante ese tiempo, pero no puede alquilarla ni disponer de ella.
\end{example}

\begin{important}
Corresponde evaluar con precisión el objetivo perseguido antes de constituir un usufructo cuando en realidad se pretende otorgar un derecho más limitado, como el de habitación: el usufructo concede facultades considerablemente mayores a las que en ocasiones se desea otorgar.
\end{important}

\section{Preferencia de la habitación sobre el usufructo}

En ciertos casos suele preferirse el derecho de habitación al usufructo. Si a la persona que cuidó de un familiar se le constituye un usufructo, esta podría alquilar la vivienda a un tercero sin habitarla ella misma, desvirtuando la finalidad original de asegurarle un lugar donde vivir. El derecho de habitación, en cambio, está más directamente orientado a garantizar la morada efectiva del beneficiario, sin perjudicarlo ni perjudicar a los demás herederos.

\begin{example}{Derecho de habitación como gesto de gratitud}
Una familia, en reconocimiento a la mujer que había cuidado a su madre durante años y que tenía prácticamente su misma edad, le constituyó un derecho real de habitación sobre la vivienda familiar hace aproximadamente diez años. La hija de la persona cuidada no deseaba que quien había asistido a su madre quedara sin hogar. Bajo esta figura, la beneficiaria solo puede habitar el inmueble ella sola, sin sumar a otros integrantes.

Con el tiempo la madre falleció y, más adelante, la propia beneficiaria desarrolló Alzheimer, quedando imposibilitada de seguir viviendo sola en la casa. La familia decidió entonces no ejercer más el derecho de habitación en especie, sino alquilar la vivienda y destinar lo producido a solventar los gastos de un geriátrico para la beneficiaria.

El derecho real de habitación no necesariamente se traduce en la ocupación efectiva del inmueble: puede reconvertirse, según las circunstancias, en una fuente de sustento económico equivalente, sin perder su función de gratitud y protección hacia quien prestó años de cuidado.
\end{example}

\begin{example}{Alojamiento de un alumno de intercambio}
El titular de un departamento que no utiliza lo destina a un alumno que llega a estudiar durante un año desde el extranjero. Se le constituye un derecho de habitación para que resida en el inmueble durante ese período; bajo esta figura solo puede vivir el beneficiario, sin poder traer a otras personas ni alquilar el bien. A diferencia del comodato, donde el comodatario podría alojar a un tercero, el derecho de habitación es personalísimo, por lo que resulta la figura más precisa cuando el objetivo es que habite exclusivamente el beneficiario.
\end{example}

\section{Límites prácticos del derecho de uso: la cochera y el uso temporal}

Cuando se pretende constituir un derecho de uso sobre una unidad complementaria, como una cochera, corresponde verificar si esa unidad admite el ejercicio de un derecho real sobre una parte indivisa cuando el beneficiario no es titular del edificio.

\begin{note}
Para situaciones de uso temporal y acotadas en el tiempo ---por ejemplo, cuatro meses al año--- no resulta conveniente constituir e inscribir un derecho real de uso en el registro. Una figura más adecuada es el contrato de comodato, más flexible, que no requiere inscripción registral y permite recuperar el uso del inmueble si el comodante cambia de opinión (sección~\ref{sec:comodato}).
\end{note}

\cornell{Los derechos de uso y de habitación}{
\cq{¿Qué diferencia al uso de la habitación?} & El uso puede recaer sobre muebles e inmuebles y es más amplio; la habitación solo recae sobre inmuebles, se limita a morar y solo existe como derecho real. Ambos son a favor de personas humanas y pueden ser vitalicios o temporales. \\ \hline
\cq{¿Qué ventaja tiene un derecho real inscripto frente a uno personal?} & Se sostiene por sí mismo sin litigar; un comodato puede ser discutido por terceros. \\ \hline
\cq{¿Puede el titular de habitación alquilar o alojar a otros?} & No: es personalísimo y solo habita el beneficiario. \\ \hline
\cq{¿Por qué preferir habitación al usufructo para quien cuidó a un familiar?} & El usufructo permitiría alquilar sin habitar; la habitación garantiza la morada efectiva y puede reconvertirse en sustento económico. \\ \hline
\cq{¿Cuándo no conviene inscribir un derecho de uso?} & En usos temporales y acotados (p.~ej. cuatro meses al año): es preferible el comodato. \\ \hline
}{El uso y la habitación protegen la vivienda de personas concretas con más solidez que un acuerdo personal. La habitación, personalísima y limitada a morar, es la figura precisa cuando solo debe habitar el beneficiario; el usufructo concede facultades mayores que a veces no se desean otorgar.}

% ------------------------------------------------------------
\chapter{El derecho real de superficie: concepto y estructura}
\label{cap:superficie}

El derecho real de superficie permite generar ingresos, financiar construcciones y aprovechar un mismo terreno en varios niveles sin que el propietario se desprenda del dominio. Este capítulo presenta su origen, su definición legal y su estructura; el siguiente desarrolla sus aplicaciones, el financiamiento, la extinción y el registro.

\section{Un caso introductorio: el inmueble desocupado de la familia}

\begin{example}{El inmueble desocupado de la familia}
Una familia es titular de varios bienes, entre ellos un inmueble totalmente desocupado. Esa misma familia necesita liquidez para reinvertirla en su empresa familiar. El interrogante es: \textbf{¿qué herramienta jurídica resulta recomendable antes de recurrir a la venta del bien?}
\end{example}

La primera reacción ante un inmueble ocioso y una necesidad de liquidez suele ser la venta, pero no siempre es lo más adecuado, sobre todo cuando el bien tiene valor de legado o relevancia patrimonial. La distinción entre quien solo ofrece la compraventa como salida y quien acompaña a una familia a conservar su patrimonio explorando otras alternativas es una diferencia fundamental en el ejercicio profesional. Entre las alternativas estudiadas se encuentran el usufructo, la locación y el derecho real de superficie, que permite generar ingresos sin desprenderse del dominio.

\begin{important}
El primer aprendizaje metodológico es \textbf{encuadrar correctamente el problema} antes de ofrecer una solución. Si no se identifican con precisión las necesidades, restricciones y objetivos del cliente, es fácil quedarse en detalles secundarios y perder de vista lo que específicamente se requiere. Cuando alguien trae un problema, la primera tarea es cómo se lo encuadra (sección~\ref{sec:encuadre}).
\end{important}

\section{Origen histórico}

El derecho real de superficie fue incorporado al ordenamiento argentino con la sanción del CCyC. La figura ya contaba con un antecedente: una \textbf{ley de forestación}, sancionada cuando aún no se hablaba del derecho de superficie como categoría general, pero que respondía a la necesidad práctica de forestar, sembrar y cosechar en terrenos ajenos sin transferir la propiedad.

\begin{example}{Por qué la forestación necesitaba esta herramienta}
La vida de un árbol no se mide en días sino en años: en muchos casos se necesitan entre 15 y 25 años para que una forestación dé sus frutos. Esa inversión no puede pensarse a corto plazo. Esta necesidad, sumada a las recomendaciones de organismos internacionales sobre el uso sostenible del medio ambiente, explica el origen de la figura: se requería una herramienta legal que otorgara al forestador la seguridad de usar, gozar y disponer de lo plantado durante todos esos años, sin ser dueño del terreno.
\end{example}

\section{Concepto legal: artículo 2114 del CCyC}
\label{sec:art2114}

Para comprender el derecho de superficie es necesario partir de la lógica clásica del dominio: quien es dueño de un terreno lo es del suelo, del subsuelo y del espacio aéreo por encima de él (el vuelo). El derecho de superficie rompe con ese paradigma, al permitir separar temporalmente esas tres capas y otorgarlas a distintos titulares.

\begin{formula}{Artículo 2114 CCyC --- Concepto}
\textit{``El derecho de superficie es un derecho real temporario, que se constituye sobre un inmueble ajeno, que otorga a su titular la facultad de uso, goce y disposición material y jurídica del derecho de plantar, forestar o construir, o sobre lo plantado, forestado o construido en el terreno, el vuelo o el subsuelo, según las modalidades de su ejercicio y plazo de duración establecidos en el título suficiente para su constitución, dentro de lo previsto en el presente Título y las leyes especiales.''}
\end{formula}

\begin{definition}{Derecho real de superficie}
Derecho real \textbf{temporario} que se constituye sobre un \textbf{inmueble ajeno} y otorga a su titular (el \textbf{superficiario}) la facultad de uso, goce y disposición material y jurídica para plantar, forestar o construir, o sobre lo ya plantado, forestado o construido, en el suelo, el vuelo o el subsuelo, según las modalidades y el plazo pactados en el título constitutivo. El superficiario puede enajenarlo a terceros, transmitiendo su posición como derecho superficiario.
\end{definition}

Los dos elementos centrales son:

\begin{itemize}
    \item \textbf{Temporalidad:} a diferencia del dominio, que es perpetuo, tiene un plazo determinado.
    \item \textbf{Ajeneidad:} se constituye sobre un inmueble ajeno, mientras que el dominio recae sobre un inmueble propio.
\end{itemize}

\subsection{Diferencia con el dominio pleno y con el usufructo}

En el dominio pleno, el propietario es simultáneamente dueño del suelo, el subsuelo y el vuelo. Con el derecho de superficie puede conservar el suelo y otorgar a otra persona el subsuelo o el vuelo en carácter de \textbf{dominio superficiario}: desmembrado y temporal.

Frente al usufructo, el superficiario tiene facultades de \textbf{uso, goce y disposición} sobre lo que construye, algo que el usufructuario no posee respecto de la cosa en sí. Además, el usufructo se extingue con la muerte del usufructuario y no se transmite a los herederos, mientras que el derecho de superficie no se extingue por el fallecimiento del superficiario y se transmite a sus herederos (sección~\ref{sec:extincion-sup}).

\section{Objeto del derecho: plantar, forestar y construir}

El objeto comprende:

\begin{itemize}
    \item \textbf{Plantar, forestar o construir} sobre el suelo, el vuelo o el subsuelo ajenos.
    \item \textbf{Adquirir un derecho ya existente} sobre algo previamente plantado, forestado o construido: no es obligatorio partir de cero.
\end{itemize}

El derecho puede recaer sobre cualquiera de las tres capas del inmueble, según las modalidades pactadas en el título constitutivo.

\begin{example}{El estacionamiento bajo un edificio}
Un titular de terreno se reserva el subsuelo porque allí funciona un estacionamiento. El titular de ese garaje es dueño de él en carácter de \textbf{dominio superficiario}: desmembrado y temporal. Vencido el plazo, se producen las consecuencias reguladas en el art.~2125 CCyC (sección~\ref{sec:efectos-extincion}).
\end{example}

\section{Pluralidad de superficiarios}

Un mismo inmueble puede estar afectado a \textbf{distintos derechos de superficie simultáneos}: el suelo, el subsuelo y el vuelo pueden entregarse a distintos superficiarios, cada uno con su propia habilitación y su propio plazo. Esta superposición potencia el aprovechamiento económico del bien, como se ilustra en el capítulo siguiente.

\cornell{Derecho de superficie: concepto y estructura}{
\cq{¿Qué es el derecho de superficie?} & Derecho real temporario sobre inmueble ajeno que faculta a plantar, forestar o construir, o a usar, gozar y disponer de lo ya plantado, forestado o construido (art.~2114 CCyC). \\ \hline
\cq{¿Cuáles son sus notas centrales?} & Temporalidad (plazo determinado) y ajeneidad (recae sobre inmueble ajeno). \\ \hline
\cq{¿Qué separa respecto del dominio pleno?} & Suelo, subsuelo y vuelo, que pueden entregarse a distintos titulares como dominio superficiario (desmembrado y temporal). \\ \hline
\cq{¿En qué se diferencia del usufructo?} & El superficiario dispone de lo construido y su derecho se transmite a los herederos; el usufructo se extingue con la muerte del usufructuario. \\ \hline
\cq{¿Por qué surgió?} & Antecedente en la ley de forestación: inversiones de 15 a 25 años sin ser dueño del terreno. \\ \hline
}{El derecho de superficie rompe la unidad clásica de suelo, subsuelo y vuelo, permitiendo a un tercero construir o plantar por un plazo determinado sin que el propietario pierda el dominio. Su rasgo distintivo frente al usufructo es la facultad de disposición y su transmisión a los herederos.}

% ------------------------------------------------------------
\chapter{Aplicaciones, financiamiento y extinción del derecho de superficie}
\label{cap:superficieaplic}

Este capítulo muestra cómo se emplea el derecho de superficie en casos concretos, cómo permite financiar vivienda, cómo se extingue y liquida, cómo se registra y qué efectos económicos y sociales genera.

\section{Casos prácticos de aplicación}

\subsection{Múltiples derechos de superficie sobre un mismo terreno}

\begin{example}{El terreno frente al predio ferial}
La titular de un terreno ubicado frente a un predio ferial, adquirido en subasta pública, no tiene intención de invertir en construir porque toda su inversión se agotó en la compra. Aparecen distintos interesados:
\begin{enumerate}
    \item Uno propone construir un garaje o estacionamiento en el subsuelo: se le otorga un \textbf{derecho de superficie por quince años sobre el subsuelo}.
    \item Un segundo quiere desarrollar un centro cultural en la planta: se le otorga un \textbf{derecho de superficie por veinte años}.
    \item Un tercero propone construir, por encima del centro cultural (desde el segundo piso hacia arriba), un centro de convenciones: se le cede por \textbf{veinte años la porción del derecho de superficie correspondiente al vuelo}.
\end{enumerate}
Cada superficiario invierte en su propia obra y la explota durante el plazo convenido. La titular cobra retribución de los tres. El dominio nunca se transfiere: al vencimiento de cada plazo, los bienes construidos revierten a su patrimonio.
\end{example}

Este mecanismo permite \textbf{aprovechar un mismo terreno en varios niveles y con varios destinos simultáneamente} sin que la propietaria pierda la titularidad ni deba vender.

\subsection{El convenio con el gobierno de la ciudad: el cargo como herramienta de protección}

En el contexto de un convenio entre el gobierno de la ciudad y el colegio profesional para la escrituración de un barrio construido para personas provenientes de un asentamiento informal, algunos predios estaban destinados a asignarse a asociaciones civiles para su funcionamiento. La cuestión: ¿conviene transferir el dominio pleno a la asociación o existe una herramienta más adecuada para proteger el destino social del bien? El riesgo de la donación en dominio pleno es que, si la asociación deja de realizar los actos para los cuales se le asignó el inmueble, o si lo vende, el objetivo del gobierno se pierde definitivamente.

\begin{example}{La escritura con cargo}
La solución propuesta fue constituir un \textbf{derecho de superficie con cargo}: el derecho se extingue si no se cumple el objetivo para el cual fue otorgado, es decir, si la asociación deja de desarrollar la actividad comprometida. El inmueble puede entonces reasignarse a otra asociación con interés en esa actividad, sin que el predio quede desocupado ni se use en beneficio de intereses ajenos al objetivo original.
\end{example}

Este mecanismo protege la \textbf{finalidad social del bien público} más allá del acto formal de la escrituración.

\subsection{El legado familiar y la alternativa a la donación}

Cuando una familia posee un terreno con \textbf{valor de legado} pero ya no puede mantenerlo económicamente, la alternativa a la donación o venta en dominio pleno es constituir un \textbf{derecho de superficie con cargo}.

\begin{example}{La donación del abuelo}
Un titular había comprado un terreno para destinarlo, junto con terrenos linderos, al desarrollo de un club social. Todos los propietarios donaron sus terrenos en dominio pleno. El club nunca se desarrolló como estaba previsto y ese terreno terminó vendido para otra obra de bien público, ajena al proyecto que había motivado la donación.

Si en lugar de donar el dominio pleno se hubiera constituido un derecho de superficie con cargo, el terreno habría podido volver a los donantes (o a sus herederos) al no cumplirse el objetivo pactado, preservando el legado familiar.
\end{example}

\subsection{La revalorización de zona: el caso del hipermercado}

Cuando una gran cadena comercial se instala en una zona genera a su alrededor un polo de actividad: galerías, locales y otros comercios. Pero puede no querer comprar el terreno, porque no tiene certeza sobre el rendimiento de la ubicación.

\begin{example}{Por qué el hipermercado prefiere un derecho de superficie}
Mediante un derecho de superficie por veinte años, la cadena construye el hipermercado, las galerías y el estacionamiento, y explota el negocio durante ese plazo. Al vencimiento puede renovar el derecho, asociarse con el titular del terreno o pagarle un alquiler por continuar.

Para el titular resulta doblemente beneficioso: cobra la retribución pactada y, si posee inmuebles en la zona, su valor aumenta por el solo hecho de tener ese polo comercial cerca.
\end{example}

La ventaja frente al usufructo es que el titular \textbf{no pierde el dominio}: lo cede temporalmente y recupera el bien con las mejoras del superficiario.

\subsection{Superficie sin ser dueño del suelo: agronegocio y bonos verdes}

En el ámbito agropecuario, empresas agroindustriales pueden desarrollar su actividad sin ser propietarias de los terrenos, aprovechando el derecho de superficie: pueden trasladarse de una zona climática a otra y cosechar a gran escala sin capital para adquirir los terrenos. La alternativa tradicional, el arrendamiento, tiene desventajas para ambas partes: el arrendatario no tiene un derecho real inscribible, y el arrendador puede ver limitada su capacidad de renegociación.

El derecho de superficie también se vincula con el mercado de \textbf{bonos verdes}: las empresas que emiten dióxido de carbono pueden compensar ese impacto promoviendo bosques a través de la compra de bonos verdes. Esta práctica, muy extendida en Europa, podría aplicarse en Argentina mediante derechos de superficie orientados a la forestación, con beneficios para los parques nacionales y el desarrollo sostenible.

\begin{example}{La cochera bajo la iglesia}
Una iglesia situada sobre una loma no deseaba vender el terreno bajo ningún concepto, dado que sobre él se encuentra el templo. Para aprovechar el subsuelo sin desprenderse del dominio se constituyeron distintos derechos de superficie que permitieron desarrollar una cochera debajo del edificio religioso, sin afectar su titularidad sobre el inmueble.
\end{example}

\section{Financiamiento de vivienda mediante el derecho de superficie}
\label{sec:financiamiento}

El derecho de superficie es una herramienta especialmente potente para el desarrollo y financiamiento inmobiliario.

\begin{example}{Financiamiento de vivienda mediante derecho de superficie}
La propietaria de un terreno constituye a favor de otra persona (el superficiario) un derecho de superficie por un plazo de 30 a 50 años. El superficiario construye veinte unidades habitacionales y puede vender a terceros el derecho de superficie sobre ellas en forma financiada: los compradores entran con una entrega inicial y pagan cuotas mensuales durante, por ejemplo, diez años.

Con lo producido, el superficiario recupera lo invertido mucho antes de que se cumpla el plazo total de su propio derecho. Se genera así un derecho a la vivienda alternativo: compradores que no pueden acceder al dominio pleno adquieren un derecho de superficie financiado, con una seguridad jurídica muy superior a la de un simple alquiler.
\end{example}

\subsection{La hipoteca sobre el derecho de superficie}

La hipoteca que garantiza la venta financiada de una unidad superficiaria se constituye sobre el propio derecho de superficie del comprador, nunca sobre el dominio del terreno, que sigue perteneciendo a su dueña original.

\begin{important}
Quien se asegura mediante hipoteca es el dueño del derecho de superficie, no la propietaria del terreno, quien no participa de esa garantía. En el ejemplo de las veinte unidades, cada comprador adquiere una sobre la base de un derecho de superficie, y esa adquisición se financia con una hipoteca sobre ese mismo derecho ---nunca sobre el dominio del terreno.
\end{important}

El superficiario originario puede incorporar en ese derecho subcontratado cláusulas para proteger su posición, oponibles dentro del propio derecho:

\begin{itemize}
    \item que el comprador no pueda alquilar la unidad mientras no cancele la hipoteca;
    \item que no pueda disponer de ella hasta el pago total del precio; y
    \item que no pueda vender a personas que el vendedor no acepte.
\end{itemize}

\subsection{Subcontratación, plazos y reversión}

Un punto de confusión frecuente es qué ocurre cuando el derecho originario (por ejemplo, de 50 años) se subcontrata por un plazo menor (por ejemplo, 20 años) con los compradores finales.

\begin{formula}{Regla de extinción y reversión en la subcontratación}
Cuando el derecho de superficie subcontratado se extingue (después de 20 años), no revierte a la propietaria original del terreno, sino al superficiario intermedio que lo había constituido, quien conserva su propio derecho por el plazo restante. El principio general de toda la figura es que \textbf{nadie puede adquirir más derechos de los que quien transmite estaba en condiciones de transmitir}.
\end{formula}

Los adquirentes de las unidades no son simples locatarios: durante el plazo de su derecho pueden venderlo, alquilarlo o hipotecarlo, con una amplitud que un locatario nunca tiene. Extinguido finalmente el derecho originario (cumplidos los 30 a 50 años, salvo pacto en contrario), la propiedad vuelve a su titular original. En el propio contrato de constitución pueden pactarse alternativas para ese momento:

\begin{itemize}
    \item que el superficiario pague una suma de dinero a cambio de que se le transmita el dominio pleno;
    \item que se prorrogue el derecho de superficie por un plazo adicional;
    \item que se reconozca al superficiario una preferencia u opción de compra sobre el inmueble; o
    \item que ambas partes queden en condominio sobre todo o parte del bien.
\end{itemize}

\begin{warning}
Se advierte sobre un riesgo frecuente del mercado inmobiliario: la compra de emprendimientos sin planos aprobados ni derecho de superficie efectivamente constituido, con la sola expectativa de que en el futuro se adjudique una unidad funcional. Al aprobarse los planos puede ser necesario modificar el proyecto o ceder una porción del terreno, lo que afecta a quienes compraron sobre un proyecto todavía no firme (sección~\ref{sec:noconsolidados}).
\end{warning}

\section{Extinción del derecho de superficie}
\label{sec:extincion-sup}

La extinción está regulada en el art.~2124 del CCyC.

\begin{formula}{Artículo 2124 CCyC --- Extinción}
\textit{``El derecho de construir, plantar o forestar se extingue por renuncia expresa, vencimiento del plazo, cumplimiento de una condición resolutoria, por consolidación y por el no uso durante diez años, para el derecho a construir, y de cinco años, para el derecho a plantar o forestar.''}
\end{formula}

Las causales son, en consecuencia:

\begin{itemize}
    \item \textbf{renuncia expresa} del superficiario;
    \item \textbf{vencimiento del plazo} convenido en el título constitutivo;
    \item \textbf{cumplimiento de una condición resolutoria} pactada;
    \item \textbf{consolidación}: reunión en una misma persona de la titularidad del suelo y del derecho de superficie; y
    \item \textbf{no uso}: diez años para el derecho a construir; cinco años para el derecho a plantar o forestar.
\end{itemize}

Las partes pueden pactar en la escritura \textbf{causales adicionales de extinción}, por ejemplo limitando el tipo de actividad del superficiario. El derecho de superficie \textbf{no se extingue por el fallecimiento del superficiario}: se transmite a sus herederos, a diferencia del usufructo.

\begin{important}
Si bien no se puede prohibir directamente la transmisión del derecho de superficie, sí es posible \textbf{limitarla}: por ejemplo, pactando que el derecho sea cedido a terceros solo en tanto estos no desarrollen la misma actividad que el titular del suelo, para evitar que la cesión instale a un competidor directo.
\end{important}

\subsection{Efectos de la extinción}
\label{sec:efectos-extincion}

Los efectos están regulados en el art.~2125 CCyC.

\begin{formula}{Artículo 2125 CCyC --- Efectos de la extinción}
\textit{``Al momento de la extinción del derecho de superficie por el cumplimiento del plazo convencional o legal, el propietario del suelo hace suyo lo construido, plantado o forestado, libre de los derechos reales o personales impuestos por el superficiario. Si el derecho de superficie se extingue antes del cumplimiento del plazo legal o convencional, los derechos reales constituidos sobre la superficie o sobre el suelo continúan gravando separadamente las dos partes, como si no hubiese habido extinción, hasta el transcurso del plazo del derecho de superficie. Subsisten también los derechos personales durante el tiempo establecido.''}
\end{formula}

En síntesis: si la extinción se produce por vencimiento del plazo, el propietario del suelo hace suyo lo construido, plantado o forestado \textbf{libre de las cargas} impuestas por el superficiario; si se produce antes del plazo pactado, los derechos reales y personales constituidos sobre la superficie continúan vigentes hasta el vencimiento del plazo original, como si no hubiera habido extinción.

\subsection{Modalidades de pago y liquidación final}

La retribución al titular del suelo puede pactarse de distintas maneras:

\begin{enumerate}
    \item \textbf{pago total al contado} al suscribir el derecho de superficie;
    \item \textbf{pago en cuotas o plazos}, a medida que se perciben los ingresos de la explotación (por ejemplo, alquileres temporarios); o
    \item \textbf{cesión de un porcentaje de las unidades construidas} a favor del titular del suelo (por ejemplo, el diez por ciento de los departamentos).
\end{enumerate}

Al extinguirse, el superficiario \textbf{no está obligado a demoler} lo construido ni a devolver el terreno vacío: lo entrega junto con la construcción. También puede pactarse de antemano cómo se reparten las unidades al finalizar el plazo.

\begin{example}{Liquidación al vencimiento}
Si se construyeron veinte departamentos, puede pactarse que al vencer el plazo el superficiario conserve ocho y el titular del suelo reciba doce, extinguiéndose el derecho de superficie y transmitiéndose la propiedad de las unidades según lo convenido.
\end{example}

\section{Aspectos registrales}

La inscripción se realiza mediante \textbf{escritura pública}, que se inscribe en el Registro de la Propiedad dentro de la \textbf{misma matrícula del inmueble}, en la parte correspondiente al dominio, del mismo modo en que se inscribe un usufructo. La titularidad del dominio permanece a nombre del propietario del suelo; el derecho de superficie se anota en la matrícula como gravamen, con la misma minuta con que se anota un dominio.

La \textbf{cancelación registral} al finalizar el plazo es análoga a la de una hipoteca: se cancela al cumplirse su objeto y debe contemplar lo pactado en la escritura, incluyendo eventuales pactos complementarios dentro del mismo instrumento.

\section{Ventajas frente a la locación y el usufructo}

El principal argumento a favor del derecho de superficie frente a la locación es jurídico: es un \textbf{derecho real}, oponible \emph{erga omnes}, mientras que la locación es un derecho personal (contractual).

\begin{example}{El depósito de camionetas}
Una empresa propone instalar un depósito de camionetas en un terreno familiar heredado con valor sentimental. La titular no quiere vender, pero la empresa necesita invertir sumas importantes en seguridad y tecnología. Si se celebra una simple locación, la inversión queda a merced de la duración del contrato y de la voluntad de las partes al vencimiento. El derecho de superficie, en cambio, otorga un \textbf{derecho real inscribible, con plazo cierto}, sobre el cual puede invertirse con seguridad jurídica. Al finalizar el plazo, la titular podría incluso vender ese derecho de superficie a un tercero.
\end{example}

\begin{table}[H]
\centering
\caption{Comparación entre el derecho de superficie y la locación}
\renewcommand{\arraystretch}{1.2}
\begin{tabularx}{\textwidth}{>{\raggedright\arraybackslash}p{0.26\textwidth}>{\raggedright\arraybackslash}X>{\raggedright\arraybackslash}X}
\toprule
\textbf{Criterio} & \textbf{Derecho de superficie} & \textbf{Locación (alquiler)} \\
\midrule
Naturaleza jurídica & Derecho real, oponible \emph{erga omnes} & Derecho personal (contrato) \\
Inscripción registral & Sí, en la matrícula del inmueble & No se inscribe como derecho real \\
Facultad de disposición & El superficiario puede disponer del derecho & El locatario no puede disponer de la cosa \\
Modificar la cosa & Sí, es la esencia del derecho (construir, plantar, forestar) & En principio no, salvo autorización expresa \\
Embargabilidad & El derecho es embargable & El contrato en sí no es un bien embargable del mismo modo \\
Seguridad para quien invierte & Alta: derecho real con plazo cierto & Baja: depende de la duración y renovación del contrato \\
Transmisión por muerte del titular & Se transmite a los herederos & Depende de lo pactado y de la ley aplicable al contrato \\
\bottomrule
\end{tabularx}
\end{table}

Frente al usufructo, cuando el superficiario es una \textbf{persona jurídica} existe además una ventaja operativa: en lugar de transmitir el derecho mediante una nueva escritura puede cederse la participación societaria (acciones o cuotas), simplificando la operación.

\begin{warning}
La facilidad de transferir el derecho mediante cesión de acciones conlleva un riesgo: si se contrató con una empresa por su reputación específica y luego esa empresa cede sus acciones a otra persona sin la misma solvencia o reputación, el titular del terreno puede quedar vinculado, de hecho, con una contraparte distinta de la que eligió. Debe considerarse al redactar la escritura constitutiva.
\end{warning}

\section{El efecto multiplicador económico y social}

La construcción no se agota en la obra: genera una \textbf{cadena de actividades conexas} que derrama beneficios hacia otros sectores. La obra da trabajo al arquitecto, a los obreros y a los gremios; esos trabajadores, a su vez, dinamizan el comercio local. Fomentar el derecho de superficie genera, en consecuencia, un efecto multiplicador que excede el negocio inmobiliario puntual.

\begin{example}{El local con dos pisos para coworking}
Un comerciante es dueño de un local y recibe la propuesta de que alguien construya dos pisos por encima para destinarlos a coworking. Teme comprometer el destino futuro de su esquina, que podría valorizarse. La solución es el derecho de superficie: sigue explotando su local en planta baja, cede por un plazo el derecho de construir y explotar el coworking en los pisos superiores, y al vencimiento recupera la disponibilidad plena del inmueble, incluidas las mejoras. Quien invierte tiene seguridad jurídica ---puede incluso vender su derecho durante el plazo--- y el propietario no compromete el destino futuro de su propiedad.
\end{example}

La misma lógica puede aplicarse al \textbf{problema habitacional}: en lugar de transferir el dominio pleno de una vivienda, pueden otorgarse derechos de uso temporal por plazos determinados ---por ejemplo, diez años--- a jóvenes que comienzan su vida económica. Transcurrido ese plazo estarían en mejores condiciones de acceder a un crédito y a una vivienda propia, liberando la unidad para otra persona joven.

\begin{example}{El modelo de vivienda comunitaria en los kibutz de Israel}
En los kibutz de Israel existen unidades habitacionales de dos, tres y cuatro ambientes que se asignan según la etapa de vida de cada familia: una pareja recién casada comienza en una unidad pequeña; cuando llegan los hijos se muda a una mayor; cuando esos hijos crecen y se van, vuelve a una más pequeña, liberando la mayor para otra familia con hijos. La propiedad no es individual: la vivienda se comparte entre la comunidad, de modo que el espacio habitacional se aprovecha de manera constante.
\end{example}

La resistencia cultural a estas herramientas no se debe a la falta de instrumentos legales, sino a la \textbf{falta de costumbre y de creatividad} de los clientes y de los propios profesionales que los asesoran: el derecho es una herramienta de creación cuando se lo ejerce con imaginación.

\cornell{Aplicaciones, financiamiento y extinción de la superficie}{
\cq{¿Cómo aprovecha un terreno con varios superficiarios?} & Se otorgan derechos separados sobre subsuelo, planta y vuelo, cada uno con su plazo; el dominio no se transfiere y los bienes revierten al vencimiento. \\ \hline
\cq{¿Para qué sirve el derecho de superficie con cargo?} & Para proteger la finalidad del bien: si no se cumple el objetivo, el derecho se extingue y el inmueble puede reasignarse o volver a la familia. \\ \hline
\cq{¿Cómo financia vivienda?} & El superficiario construye y vende financiado el derecho de superficie sobre las unidades, con hipoteca sobre ese derecho y nunca sobre el terreno. \\ \hline
\cq{¿A quién revierte un derecho subcontratado que se extingue?} & Al superficiario intermedio, no al dueño del terreno: nadie transmite más derechos de los que tiene. \\ \hline
\cq{¿Cuáles son las causales de extinción (art.~2124)?} & Renuncia expresa, vencimiento del plazo, condición resolutoria, consolidación y no uso (10 años para construir; 5 para plantar o forestar). \\ \hline
\cq{¿Qué efectos tiene la extinción (art.~2125)?} & Por vencimiento, el dueño del suelo hace suyo lo construido libre de cargas; si es anticipada, los derechos constituidos continúan hasta el plazo original. \\ \hline
\cq{¿Por qué es mejor que la locación?} & Es un derecho real inscribible, con plazo cierto, disponible y transmisible a herederos, lo que da seguridad a quien invierte. \\ \hline
}{El derecho de superficie multiplica los usos de un terreno, permite financiar vivienda y proteger destinos sociales o familiares, y se extingue por causales legales y pactadas con efectos regulados por los arts.~2124 y 2125. Su ventaja sobre la locación es su naturaleza real; su riesgo, la posible cesión a contrapartes no queridas.}

% ------------------------------------------------------------
\chapter{Elegir la figura adecuada: comparación de los derechos sobre cosa ajena}
\label{cap:comparacion}

Después de estudiar cada derecho por separado, este capítulo los compara y presenta el criterio de fondo que atraviesa la Parte~\ref{parte:cosaajena}: antes de aplicar automáticamente una figura tradicional, conviene analizar la situación concreta de cada persona o familia y diseñar la combinación de derechos y cláusulas que mejor proteja sus intereses reales.

\section{Cuadro comparativo}

\begin{table}[H]
\centering
\caption{Derechos reales sobre inmueble ajeno}
\small
\renewcommand{\arraystretch}{1.2}
\begin{tabularx}{\textwidth}{>{\raggedright\arraybackslash}p{0.14\textwidth}>{\raggedright\arraybackslash}X>{\raggedright\arraybackslash}X>{\raggedright\arraybackslash}X}
\toprule
\textbf{Derecho real} & \textbf{Principales facultades} & \textbf{Limitaciones} & \textbf{Uso típico} \\
\midrule
Usufructo & Usar, gozar, percibir frutos, ceder o alquilar & No puede deteriorar ni alterar el destino & Donación con reserva de renta; protección del donante \\
Uso & Usar el inmueble para sí y su familia & No puede arrendarlo; no transferible & Uso familiar temporal o estacional \\
Habitación & Habitar personalmente el inmueble & No puede albergar a terceros distintos de sí mismo & Protección de vivienda de persona vulnerable \\
Superficie & Construir, plantar y explotar sobre suelo ajeno & Plazo máximo legal; al vencer, revierte & Proyectos de construcción o explotación a largo plazo \\
\bottomrule
\end{tabularx}
\end{table}

El derecho de superficie es el más robusto de los derechos reales sobre cosa ajena: permite construir o plantar sobre un inmueble ajeno durante un plazo máximo fijado por el CCyC y, al vencer, lo construido revierte al propietario del terreno sin indemnización, salvo pacto en contrario. Una ventaja adicional frente al usufructo es que el superficiario puede, a su vez, alquilar o subarrendar lo construido.

\begin{example}{Salón de fiestas sobre terreno ajeno}
Una constructora necesita un predio para levantar un salón de fiestas. El titular de un terreno le otorga un derecho de superficie por 40 a 50 años para construir y explotar el salón; al finalizar el plazo, lo construido vuelve al titular del terreno.
\end{example}

\begin{example}{Construcción de departamentos por una empresa sobre terreno ajeno}
Una empresa desea construir departamentos sobre un terreno que no le pertenece. El propietario no está dispuesto a vender pero acepta que la empresa construya y sea titular de lo edificado durante 50 años. Al vencer el plazo, si no se renueva, lo construido vuelve al titular del terreno.
\end{example}

\section{Comodato frente a derechos reales}
\label{sec:comodato}

Encasillar una situación dentro de un derecho real o de un tipo contractual cuando ello no es necesario puede perjudicar al cliente. Para un uso acotado a pocos meses, por ejemplo, no resulta conveniente constituir e inscribir un derecho real de uso.

\begin{table}[H]
\centering
\caption{Comodato, derecho de uso y locación temporal}
\small
\renewcommand{\arraystretch}{1.2}
\begin{tabularx}{\textwidth}{>{\raggedright\arraybackslash}p{0.16\textwidth}>{\raggedright\arraybackslash}p{0.17\textwidth}c>{\raggedright\arraybackslash}X>{\raggedright\arraybackslash}X}
\toprule
\textbf{Figura} & \textbf{Tipo} & \textbf{Inscripción} & \textbf{Duración típica} & \textbf{Recomendado para} \\
\midrule
Derecho de uso & Derecho real & Sí (registro) & Mediano/largo plazo & Uso permanente o estable \\
Comodato & Contrato personal & No & Corto/mediano plazo & Préstamo temporal sin costo \\
Locación temporal & Contrato personal & No (salvo más de 3 años) & Meses o años & Uso oneroso con plazo determinado \\
\bottomrule
\end{tabularx}
\end{table}

\begin{note}
El comodato es un contrato de préstamo de uso, distinto de los derechos reales analizados. Distinguir con precisión cuándo corresponde constituir un derecho real y cuándo basta un vínculo contractual personal es central para un asesoramiento notarial adecuado.
\end{note}

\section{Aplicar lo aprendido: consigna integradora}

Como cierre de este bloque se propone un ejercicio breve y aplicado: plantear situaciones concretas ---no hipótesis abstractas--- en las que se aplicaría cada uno de los derechos reales estudiados (uso, habitación, usufructo, superficie, entre otros), incorporando las cláusulas pertinentes en cada caso.

\begin{exercise}{Trabajo práctico integrador}
Formule la respuesta en primera persona y con brevedad, por ejemplo: «Yo constituiría un derecho real de uso porque es conveniente para la explotación de tal cosa; a ella le daría un usufructo por tal motivo; a él le daría el derecho de habitación por tal situación; acá constituiría un derecho real de superficie sobre el subsuelo por tal circunstancia».

Se espera brevedad y aplicabilidad práctica, no extensión: un ejercicio breve aplicado a casos concretos y propios de cada estudiante, y no un trabajo extenso de veinte o cuarenta páginas. Su finalidad principal no es acreditar un puntaje, sino verificar si los contenidos fueron comprendidos y pueden aplicarse a situaciones reales.
\end{exercise}

\cornell{Comparación de los derechos sobre cosa ajena}{
\cq{¿Qué derecho real es el más robusto sobre cosa ajena?} & La superficie: permite construir o plantar por un plazo máximo legal y lo construido revierte sin indemnización, salvo pacto. \\ \hline
\cq{¿Qué distingue a usufructo, uso y habitación?} & El usufructo puede alquilarse y transferirse onerosamente; el uso es para el titular y su familia, sin arrendar ni transferir; la habitación es personalísima y limitada a morar. \\ \hline
\cq{¿Cuándo conviene el comodato?} & En usos temporales sin costo: no requiere inscripción, es flexible y permite recuperar el uso. \\ \hline
\cq{¿Qué riesgo hay en encasillar sin necesidad?} & Constituir un derecho real donde bastaba un vínculo personal puede perjudicar al cliente. \\ \hline
}{Cada figura responde a una necesidad distinta: la superficie a construir y explotar, el usufructo a explotar y gozar, el uso a la vivienda familiar temporal, la habitación a proteger a una persona y el comodato a préstamos breves. La elección debe partir del objetivo concreto del cliente.}

% ============================================================
% PARTE V — SOCIEDADES, CONDOMINIO Y SUCESIÓN
% ============================================================
\part{Sociedades, condominio y sucesión: la planificación patrimonial}
\label{parte:planificacion}

% ------------------------------------------------------------
\chapter{Usufructo de acciones y cláusulas societarias en la empresa familiar}
\label{cap:empresafamiliar}

Los derechos reales y las sociedades se cruzan cuando el bien no es un inmueble sino una participación societaria. Este capítulo muestra cómo la lógica del usufructo se traslada a las acciones y cuotas sociales, cómo se planifica la sucesión de una empresa familiar y cómo las cláusulas societarias de arrastre y acompañamiento pueden resolver bloqueos entre socios y, más adelante, entre condóminos.

\section{Usufructo de acciones y cuotas sociales}
\label{sec:usufructoacciones}

En lo estructural, el usufructo de acciones sigue reglas equivalentes a las del usufructo inmobiliario: si la acción es un bien ganancial corresponde el asentimiento conyugal en los mismos supuestos y con las mismas excepciones (sección~\ref{sec:asentimiento}); y el titular puede reservarse el usufructo o constituirlo a favor de un tercero, sin distinciones respecto del régimen aplicable al dominio inmobiliario.

\begin{definition}{Derechos políticos y derechos económicos}
La acción de una sociedad comprende dos categorías de derechos: los \textbf{derechos económicos} (dividendos, cuota de liquidación) y los \textbf{derechos políticos} (voto en las asambleas). Históricamente, la doctrina asociaba el usufructo de acciones exclusivamente a los derechos económicos, dejando los derechos políticos en cabeza del nudo propietario.
\end{definition}

\begin{important}
\textbf{Referencia normativa: Ley de Sociedades.} Se hace referencia a la Ley 19.550 de Sociedades Comerciales y a los proyectos de reforma societaria en debate, que regulan de manera distinta la posición del usufructuario minoritario frente a las decisiones políticas de la sociedad.
\end{important}

\begin{example}{Donación de acciones sin reserva de derechos políticos}
Un padre dona a su hijo las acciones de la empresa familiar, con fines de planificación sucesoria, reservándose únicamente los derechos económicos (dividendos y cuota de liquidación), sin advertir que cede también los derechos políticos.

Como consecuencia, el hijo ---titular de los derechos políticos--- puede votar en una asamblea la disolución y liquidación de la sociedad, con lo cual el padre queda con derecho a una cuota de liquidación pero pierde la posibilidad de seguir cobrando dividendos, que era lo que pretendía conservar.

Otros supuestos en que el nudo propietario, ejerciendo en soledad los derechos políticos, puede licuar o extinguir la expectativa económica del usufructuario: un aumento de capital que diluye su participación (por ejemplo, del 50\% original a un 10\%), una fusión societaria frente a la cual el usufructuario solo puede ejercer el derecho de receso, o el cambio del objeto social hacia una actividad distinta.
\end{example}

La conclusión doctrinaria expuesta es que, tratándose de derechos políticos, la nuda propiedad no es una caja vacía de contenido: quien se reservó el uso y el goce transmite igualmente al nudo propietario algo valioso.

\begin{note}
Como alternativa superadora al esquema de todo o nada, es posible graduar en el tiempo la cesión de derechos económicos y políticos según la maduración del sucesor dentro de la empresa: se dona la acción reservándose el usufructo de los derechos políticos y económicos y, con los años, se cede un porcentaje creciente de los derechos económicos reservados. La justificación es formativa: permite que quien recibe la nuda propiedad se prepare para ejercer en el futuro los derechos políticos en la asamblea.
\end{note}

Sobre la comparecencia en la asamblea con usufructo sobre las acciones, la regla depende de a quién se reservaron los derechos políticos: si el usufructuario los conserva, comparece y vota, y el nudo propietario no comparece y solo percibe los derechos económicos que le correspondan. Cuando existe comparecencia conjunta en condominio de derechos rige la unificación de la voluntad de voto: una acción, un voto.

\begin{definition}{Fideicomiso accionario}
Estructura mediante la cual las acciones ---gravadas con usufructo y nuda propiedad--- se aportan a un fideicomiso, cuyo fiduciario ejerce los derechos políticos conforme a las pautas fijadas en el instrumento de constitución, en lugar de que cada cousufructuario o nudo propietario comparezca individualmente a la asamblea.
\end{definition}

\section{Planificación sucesoria de la empresa familiar}

El usufructo de acciones es una herramienta de planificación sucesoria empresarial que permite al fundador transmitir la propiedad de las acciones a la siguiente generación ---con el beneficio de ordenar el patrimonio en vida y evitar conflictos futuros entre herederos--- sin resignar necesariamente el control económico ni político del negocio mientras viva y lo considere necesario. La combinación de reserva de usufructo, derechos políticos condicionados a determinados actos societarios (fusión, aumento de capital, cambio de objeto social) y cesión progresiva de derechos económicos conforma un esquema más rico y flexible que la donación lisa y llana de las acciones.

\section{Cláusulas de arrastre y de acompañamiento}
\label{sec:dragtag}

Un problema típico del derecho societario es la situación de los socios minoritarios frente a las decisiones de los mayoritarios. Para resolverlo existen dos cláusulas incorporadas expresamente al ordenamiento a partir de la Ley de Sociedades por Acciones Simplificadas (Ley de S.A.S.):

\begin{definition}{Derecho de arrastre (\emph{drag-along})}
Cláusula que obliga a los socios minoritarios a vender también su participación cuando el socio principal decide vender la suya, de modo que el comprador pueda adquirir el 100\% del capital y los minoritarios no queden atrapados en una sociedad cuyo control cambió de manos sin ellos.
\end{definition}

\begin{definition}{Derecho de acompañamiento (\emph{tag-along})}
Cláusula que otorga a los minoritarios la posibilidad ---no la obligación--- de vender su participación en las mismas condiciones y al mismo valor que se le pagó al socio principal, si así lo desean. Si no quieren vender, no están obligados a hacerlo.
\end{definition}

Estas cláusulas son especialmente convenientes en el diseño de un startup: quien tuvo la idea original y conserva, por ejemplo, un 30\% o un 40\% del capital puede verse sin ningún poder real de decisión si el socio mayoritario decide vender y no cuenta con una cláusula de acompañamiento que le permita salir en las mismas condiciones.

\cornell{Usufructo de acciones y cláusulas societarias}{
\cq{¿Qué derechos comprende una acción?} & Económicos (dividendos, cuota de liquidación) y políticos (voto en asambleas). \\ \hline
\cq{¿Qué riesgo hay al donar acciones reservándose solo los derechos económicos?} & Se ceden también los políticos: el hijo puede votar la disolución, aumentar el capital diluyendo al padre, fusionar o cambiar el objeto. \\ \hline
\cq{¿Cómo se gradúa la cesión?} & Reservándose el usufructo de derechos políticos y económicos y cediendo progresivamente los económicos según la maduración del sucesor. \\ \hline
\cq{¿Quién vota con usufructo de acciones?} & Quien conserva los derechos políticos; con comparecencia conjunta, una acción, un voto unificado. \\ \hline
\cq{¿Qué es el fideicomiso accionario?} & Aporte de las acciones a un fideicomiso cuyo fiduciario ejerce los derechos políticos según el instrumento. \\ \hline
\cq{¿Qué diferencia hay entre arrastre y acompañamiento?} & Arrastre: los minoritarios deben vender con el mayoritario. Acompañamiento: pueden, si quieren, vender en las mismas condiciones. \\ \hline
}{El usufructo de acciones permite planificar la sucesión de la empresa familiar conservando control, pero exige distinguir derechos económicos y políticos, y graduar su cesión. Las cláusulas de arrastre y acompañamiento protegen a mayoritarios y minoritarios frente a un cambio de control.}

% ------------------------------------------------------------
\chapter{Condominio, comunidad hereditaria y herederos en la sociedad}
\label{cap:condominio}

El condominio nace cuando el dominio se reparte entre varios cotitulares; la comunidad hereditaria funciona en la práctica como un condominio, pero no lo es en sentido técnico. Este capítulo trata los conflictos que surgen entre copropietarios, entre herederos y frente a la sociedad, y traslada al condominio la lógica de las cláusulas estudiadas en el capítulo anterior.

\section{Las cláusulas de arrastre y acompañamiento aplicadas al condominio}

La lógica del arrastre y del acompañamiento puede trasladarse al condominio. En un condominio compuesto por tres personas con el 20\%, el 40\% y el 40\% del bien, basta con que una sola no quiera vender para que nadie esté dispuesto a comprar una propiedad con apenas el 99\% de sus partes libres.

Por eso puede proponerse que los condóminos firmen un pacto que replique esa lógica: si existe un comprador interesado en la totalidad del inmueble, todos los condóminos quedan obligados a vender; y si la mayoría decide vender, la minoría puede sumarse a la venta en las mismas condiciones.

\begin{important}
Los pactos entre condóminos son oponibles entre las partes que los suscriben, pero en principio no lo son frente a terceros. Sin embargo, la ley admite que un pacto sea oponible a un tercero que haya tomado conocimiento efectivo de su existencia, de modo que basta exhibirle el pacto para que quede alcanzado por él.
\end{important}

Para instrumentar la obligación de vender suele recurrirse a un poder especial e irrevocable otorgado al firmar el pacto, que permite ejecutar la venta llegado el caso sin depender de la voluntad posterior del condómino remiso. Sin ese poder, el condómino podría simplemente negarse a firmar cuando llegara el momento.

\section{División de condominio y comunidad hereditaria}

Dos situaciones producen efectos prácticos similares pero tienen distinta naturaleza jurídica: el condominio propiamente dicho, nacido de un derecho real, y la comunidad hereditaria, que a los efectos prácticos funciona como un condominio pero no lo es desde el punto de vista técnico.

\begin{example}{Herencia indivisa y venta}
Dieciséis herederos, de los cuales quince están de acuerdo en vender y uno se niega: la vía es la \textbf{división judicial de la herencia}, ya que no existe otra forma de forzar la partición si no se logra un acuerdo extrajudicial, y la cuestión queda directamente judicializada.

Si en cambio se trata de un condominio de derecho real y no de una comunidad hereditaria, la vía es la \textbf{acción de división de condominio}, también judicial si no hay acuerdo.
\end{example}

\section{Riesgos de comprar derechos no consolidados}
\label{sec:noconsolidados}

Un riesgo frecuente en el mercado actual es la compra de emprendimientos que todavía no cuentan con planos aprobados ni con un derecho de superficie efectivamente constituido, sino apenas con la expectativa de que en el futuro se adjudique una unidad funcional.

\begin{important}
Al aprobarse finalmente los planos puede ser necesario modificar el proyecto original o ceder una porción del terreno, lo que afecta directamente a quienes compraron sobre un proyecto todavía no firme jurídicamente.
\end{important}

\begin{example}{El departamento en desarrollo sin derecho constituido}
Quien compra debe preguntarse: ¿a quién le estoy comprando? ¿Qué es lo que estoy comprando exactamente? Muchas veces lo que se adquiere no es un derecho real consolidado sino apenas la cesión de un boleto o de una expectativa contractual. El riesgo no depende únicamente del nivel de sofisticación del comprador: muchas personas que pagan sumas millonarias cuentan con asesores y escribanos y aun así compran derechos que no están firmes.
\end{example}

\section{Incorporación de herederos y quórum en asambleas}
\label{sec:herederos}

La relación entre condominio, comunidad hereditaria y sociedades se torna compleja cuando fallece un socio o accionista y sus herederos todavía no fueron declarados. Mientras no exista declaratoria de herederos ni partición no hay certeza definitiva sobre quién representa al causante en la sociedad, y la propia partición podría adjudicar esa participación a un heredero distinto del que venía ejerciendo los derechos del fallecido.

\begin{example}{Exclusión de herederos en asamblea}
La familia de un socio fallecido asiste, al día siguiente de velarlo, a una asamblea de la que los herederos son excluidos con el argumento de que todavía no existe declaratoria de herederos que los legitime. Estas situaciones pueden derivar en decisiones sociales que perjudican los intereses del causante y de sus sucesores.
\end{example}

Respecto de las asambleas extraordinarias es necesario distinguir entre quórum y mayoría:

\begin{definition}{Quórum en asambleas}
Porcentaje de capital con derecho a voto que debe estar presente para que la asamblea pueda sesionar válidamente. En primera convocatoria se exige un quórum del 60\% de las acciones con derecho a voto; en segunda convocatoria ese quórum se reduce a un porcentaje sensiblemente menor.
\end{definition}

\begin{definition}{Mayoría en asambleas}
Porcentaje de votos necesario, sobre los presentes, para aprobar una decisión.
\end{definition}

\begin{important}
El estatuto puede elevar los porcentajes mínimos legales de quórum, pero nunca hasta el punto de exigir una unanimidad encubierta.
\end{important}

Existe debate doctrinario sobre el momento en que se considera incorporado un heredero a la sociedad. Hay tres posturas: (1) al fallecimiento del causante; (2) a la declaratoria de herederos y designación de administrador de la sucesión; y (3) a la partición definitiva. Es una cuestión difícil de resolver en la práctica y está siendo objeto de tratamiento en el proyecto de reforma de la ley de sociedades.

\cornell{Condominio, comunidad hereditaria y herederos en la sociedad}{
\cq{¿Por qué aplicar arrastre y acompañamiento a un condominio?} & Porque basta que un condómino no quiera vender para que nadie compre el inmueble; el pacto obliga a vender ante un comprador de la totalidad o permite sumarse a la mayoría. \\ \hline
\cq{¿Cómo se ejecuta la obligación de vender?} & Con un poder especial e irrevocable otorgado al firmar el pacto. Los pactos obligan a las partes y a terceros que los conozcan efectivamente. \\ \hline
\cq{¿Cómo se divide un condominio o una herencia con un disidente?} & Condominio de derecho real: acción de división de condominio. Comunidad hereditaria: división judicial de la herencia. \\ \hline
\cq{¿Qué se compra al adquirir derechos no consolidados?} & A veces solo la cesión de un boleto o expectativa sobre un proyecto sin planos aprobados. \\ \hline
\cq{¿Qué diferencia hay entre quórum y mayoría?} & Quórum: capital presente necesario para sesionar. Mayoría: votos necesarios, sobre los presentes, para aprobar. \\ \hline
\cq{¿Desde cuándo se incorpora el heredero a la sociedad?} & Hay tres posturas: desde el fallecimiento, desde la declaratoria de herederos o desde la partición; está en debate en el proyecto de reforma. \\ \hline
}{El condominio y la comunidad hereditaria requieren vías distintas de división, y los pactos inspirados en el arrastre y el acompañamiento evitan bloqueos. En las sociedades, la falta de certeza sobre los herederos plantea conflictos de participación en asambleas que la doctrina resuelve con tres posturas.}

% ============================================================
% PARTE VI — MARCO JURÍDICO DE LA ECONOMÍA DIGITAL
% ============================================================
\part{El documento digital: gobierno electrónico, internet, firma y prueba}
\label{parte:digital}

% ------------------------------------------------------------
\chapter{Gobierno electrónico y marco normativo digital}
\label{cap:gobierno}

La segunda gran orientación del curso, junto con la tradicional del derecho inmobiliario, es la tecnológica. Está dirigida a los operadores jurídicos ---notarios, abogados y martilleros--- que trabajan hoy en un ecosistema en el que la firma digital, la prueba electrónica y los contratos electrónicos son la forma cotidiana de dar fe, de contratar y de litigar. Este capítulo parte del origen del marco normativo digital argentino.

\section{El criterio de la escalera}

La tecnología debe pensarse como una escalera y no como un salto, con la misma lógica con que se construye una casa: primero los cimientos, después las paredes y al final el techo. No se puede hablar de blockchain sin haber entendido antes qué es una firma digital, y no se puede hablar de firma digital sin entender por qué el Estado empezó a digitalizarse en el año 2000. Cada bloque temático es un escalón que sostiene al siguiente. Es también la actitud necesaria para no quedar desactualizado en una profesión que cambia más rápido que la ley que la regula.

\begin{concept}{Aprendizaje escalonado}
Si no se comprende la firma digital no puede avanzarse hacia la tecnología blockchain, y sin esta no es posible comprender la actualización tecnológica del derecho. Los capítulos de esta Parte y de la siguiente siguen ese orden.
\end{concept}

\section{El gobierno electrónico}

Para comprender el origen del marco normativo digital argentino es necesario partir del gobierno electrónico, su punto de partida histórico.

\begin{definition}{Gobierno electrónico}
Materia propia del derecho administrativo. Consiste en la relación entre el Estado y sus administrados, que deja de ser presencial y pasa a darse a través de los canales digitales.
\end{definition}

\subsection{El contexto del año 2000}

El fenómeno se ubica en el año 2000. Ya se venía consolidando en el mundo un debate en torno al gobierno electrónico, vinculado por algunas corrientes doctrinarias a fenómenos de \textbf{control social}, dado que la digitalización del Estado implica un seguimiento y un perfilamiento del ciudadano. Sin embargo, en ese año aún no existían las herramientas de inteligencia artificial generativa, el uso masivo de datos biométricos ni el desarrollo actual de la internet de las cosas, por lo que la privacidad, aunque ya expuesta, no enfrentaba el nivel de exposición actual.

\subsection{Adopción de la firma digital}

\begin{example}{Los colegios profesionales y la firma digital}
En los primeros intentos de implementación de la firma digital, algunos colegios profesionales pequeños evaluaron adoptarla, pero los costos resultaban imposibles y se optó por no avanzar. Esa decisión fue calificada como un error.

Más allá de la existencia de la tecnología, la ciudadanía no adoptó masivamente la firma digital porque la administración pública continuó vinculándose con el ciudadano de manera presencial, de mostrador, y porque faltaba un cambio cultural previo para que la implementación tecnológica prosperara.
\end{example}

\subsection{El factor catalítico y el aceleracionismo}

Se planteó cuál fue el factor catalítico que hizo que este proceso explotara. Se introdujo la referencia a una corriente filosófica, la de los \textbf{aceleracionistas}, que analiza cómo los factores disruptivos ---como las irrupciones tecnológicas--- generan sacudidas sociales que llevan a la sociedad, como si fuese un corral de ovejas, hacia determinado lado.

\section{Las leyes 25.506 y 25.326}

Ambas leyes fundacionales del derecho digital argentino fueron sancionadas en el año 2000.

\begin{itemize}
    \item \textbf{Ley 25.506, de Firma Digital.} Regula el régimen de la firma digital y su equiparación jurídica con la firma manuscrita en las condiciones que la propia ley establece. Es la primera norma sobre la cual se trabaja.
    \item \textbf{Ley 25.326, de Protección de Datos Personales.} Regula el tratamiento de datos personales asentados en archivos, registros, bancos de datos u otros medios técnicos de tratamiento de datos.
\end{itemize}

\begin{important}
Las leyes 25.506 y 25.326 nacen del mismo impulso histórico: la construcción del gobierno electrónico. Mientras la primera habilita jurídicamente la actuación digital del Estado y de los particulares mediante la firma digital, la segunda busca contener los riesgos que esa misma digitalización genera sobre la privacidad del ciudadano. Son dos caras de un mismo proceso.
\end{important}

\section{Decreto 182/2019}

El decreto reglamentario vinculado a la infraestructura de firma digital es el \textbf{Decreto 182/2019}, que tiene una modificación posterior en el año 2024. Su lectura, junto con la de la Ley de Firma Digital y los arts.~286, 287, 288, 317 y 319 del CCyC, fue indicada como material de estudio para abordar el capítulo~\ref{cap:firma}.

\cornell{Gobierno electrónico y marco normativo digital}{
\cq{¿Qué es el gobierno electrónico?} & Materia del derecho administrativo: la relación entre el Estado y los administrados deja de ser presencial y pasa a los canales digitales. \\ \hline
\cq{¿Qué leyes surgen en el año 2000 y qué regulan?} & Ley 25.506 (firma digital) y Ley 25.326 (protección de datos personales): dos caras del mismo proceso de digitalización. \\ \hline
\cq{¿Por qué no se adoptó masivamente la firma digital al principio?} & Costos para colegios pequeños, administración de mostrador y falta de cambio cultural previo. \\ \hline
\cq{¿Qué es el criterio de la escalera?} & Aprender la tecnología por etapas: gobierno electrónico, firma digital, blockchain. \\ \hline
}{El marco normativo digital argentino nace con el gobierno electrónico del año 2000: la firma digital habilita la actuación digital y la protección de datos contiene sus riesgos. Comprender ese origen es el primer escalón para estudiar firma, prueba, blockchain y tokenización.}

% ------------------------------------------------------------
\chapter{De ARPANET a la Web 2.0: internet, datos y plataformas}
\label{cap:web}

La evolución de internet se presenta en tres etapas: Web 1.0, Web 2.0 y Web 3.0. En todas ellas se habla siempre de comunicaciones electrónicas. Este capítulo desarrolla las dos primeras y sus consecuencias sociales, institucionales y probatorias; la Web 3.0, tecnología de valor, se estudia en la Parte~\ref{parte:blockchain}.

\section{El circuito de la innovación tecnológica}

\begin{important}
Toda tecnología y todo descubrimiento tecnológico surge primero en el ámbito militar, luego se traslada a la academia y recién después al ámbito comercial. Ese es el circuito que siguen las innovaciones tecnológicas.
\end{important}

\section{ARPANET y la Web 1.0}

El origen del ecosistema digital se sitúa en la década de 1960 con \textbf{ARPANET}, la primera red de comunicación a distancia, antecesora del consorcio que daría origen a la World Wide Web. En su funcionamiento básico la información se organiza en \textbf{paquetes de datos} que no viajan por un único canal: se dividen y se distribuyen por múltiples vías de la red y se reúnen en el proveedor de servicios de internet del destinatario. Este esquema reemplazó progresivamente a los primeros modelos de redes centralizadas, controladas desde un único centro de cómputos, y dio paso a una lógica de descentralización de la información (véase la sección~\ref{sec:redes}). Esta infraestructura se apoya en una base material, muchas veces invisible: los cables submarinos que atraviesan los océanos.

Aplicando el circuito de la innovación, la primera etapa corresponde al momento en que las comunicaciones electrónicas circulan entre universidades: los hitos son la primera comunicación exitosa a distancia entre casas de estudio y el primer correo electrónico. Como hito generacional previo a las redes sociales se mencionan los primeros sistemas de mensajería instantánea (ICQ, MSN Messenger), antecedentes directos de WhatsApp.

\section{La Web 2.0 y las empresas GAFAM}

La Web 2.0 es la web de las redes sociales (Facebook, YouTube, WhatsApp, entre otras). En esta etapa aparecen las grandes empresas tecnológicas, llamadas «tigres tecnológicos» o \textbf{GAFAM}: Google, Amazon, Facebook, Apple y Microsoft.

\begin{definition}{Prosumidor}
Usuario de la Web 2.0 que \textbf{produce contenido} y, a la vez, lo \textbf{consume}, a diferencia del espectador pasivo de la Web 1.0.
\end{definition}

\subsection{El activo principal de las plataformas: los datos}

El activo principal de estas empresas es la información, pero la respuesta es demasiado genérica: se trata de la información y los datos de los propios usuarios. El producto gratis es el propio usuario. La lógica de las plataformas puede describirse con la imagen de un corral: «entren, entren\ldots{} y cierro, y ahora soy dueño de todo esto». Al aceptar los términos y condiciones, el usuario entrega todo el contenido generado durante años, sin retener derechos de propiedad sobre él.

\begin{important}
Las grandes plataformas construyeron un modelo de negocio basado en la obtención de datos de los usuarios a cambio de servicios aparentemente gratuitos. Las personas cedieron sus datos de manera masiva y voluntaria porque, a cambio, accedían a un mundo de conexión antes impensado: reencuentros con compañeros de estudio, vínculos recuperados después de años y nuevas formas de relación social y patrimonial.
\end{important}

\begin{concept}{Información es poder}
El valor, en la economía digital, reside en los datos y en quién tiene acceso y control sobre ellos. Se recuerda como ilustración la máquina Enigma de la Segunda Guerra Mundial: una vez descifrado el código alemán, los servicios británicos debieron en ciertas ocasiones no actuar sobre la información interceptada, incluso a costa de vidas humanas, para no revelar que el código había sido quebrado.
\end{concept}

\begin{example}{Consecuencias sociales y legales del uso de las redes}
El uso masivo de redes produjo situaciones que impactaron en lo jurídico: reencuentros en Facebook entre exparejas de la adolescencia que derivaron en la disolución de matrimonios constituidos, y modificaciones ---aumentos o reducciones--- de cuotas alimentarias fundadas en publicaciones que evidenciaban la situación económica real de alguna de las partes (véase el capítulo~\ref{cap:prueba}).
\end{example}

Como lectura recomendada sobre los aspectos socio-técnicos de estos fenómenos se sugiere la obra de la antropóloga argentina radicada en Brasil Paula Sibilia, quien desarrolla el concepto de «extimidad» para describir la exposición voluntaria de la intimidad en las plataformas digitales.

\begin{example}{Las ideas disruptivas y el caso de Steve Jobs}
Las ideas disruptivas suelen ser rechazadas en un primer momento. Como analogía se recuerda la biografía de Steve Jobs: sus propuestas iniciales fueron resistidas dentro de su propia empresa, lo que lo llevó a apartarse y fundar otra compañía, para luego regresar y transformar definitivamente la industria. Ideas que en 2009 parecían extravagantes hoy resultan, con el tiempo, acertadas.
\end{example}

\section{La respuesta del notariado: Comisión de Innovación y Tecnología}

Como respuesta institucional se recuerda la experiencia de la Comisión de Innovación y Tecnología del colegio profesional, impulsada por seis colegas inquietos por la irrupción de la tecnología en la práctica y presidida entre 2020 y 2023 por una de sus impulsoras. Su primer proyecto fue un \textbf{glosario} para traducir al español la abundante terminología técnica que ingresaba en inglés: el equipo revisaba a diario entre diez y veinte noticias del ámbito digital y extraía de ellas los términos nuevos, hasta reunir cerca de trescientas palabras. El material, elaborado en 2021, quedó publicado en el sitio web de la institución.

El segundo desarrollo fue un \textbf{libro}, publicado en 2023, surgido de un formato de exposición interna llamado «café notarial», en el que cada comisión presenta a los colegas sus proyectos anuales. Sus autores no perciben derechos económicos: toda la tirada fue donada a la biblioteca de la institución. Está en preparación un nuevo libro sobre derecho informático, cuya finalización se posterga porque la materia se actualiza más rápido de lo que puede cerrarse cada capítulo.

\begin{important}
\textbf{Actualización normativa reciente.} Se citó una resolución del gobierno de la Ciudad que obliga a los administradores de consorcio a llevar los libros de actas y de registro en \textbf{formato digital}, en lugar del libro en papel. Al momento de la clase no se había podido acceder al texto completo de la resolución. Se la cita como ejemplo de la velocidad con que la normativa se adapta a la digitalización de la práctica profesional.
\end{important}

\cornell{Internet, datos y plataformas}{
\cq{¿Qué es ARPANET y qué introdujo?} & La primera red de comunicación a distancia (década de 1960); introdujo el envío por paquetes de datos por múltiples vías, en reemplazo de redes centralizadas. \\ \hline
\cq{¿Cómo se diferencian Web 1.0, 2.0 y 3.0?} & 1.0: estática y unidireccional. 2.0: redes sociales; el usuario produce y consume (prosumidor). 3.0: tecnología de valor, transferencia de valor entre pares. \\ \hline
\cq{¿Cuál es el activo de las plataformas GAFAM?} & Los datos de los propios usuarios, cedidos a cambio de servicios «gratuitos»: el usuario es el producto. \\ \hline
\cq{¿Qué enseña el caso Enigma?} & Que información es poder: controlar la información tiene un precio estratégico. \\ \hline
\cq{¿Qué hizo la Comisión de Innovación y Tecnología?} & Elaboró un glosario de cerca de 300 términos (2021) y un libro (2023) surgido del «café notarial». \\ \hline
}{La evolución de internet pasa del usuario espectador al prosumidor, cuyos datos son el activo de las plataformas; esa lógica de control de la información prepara el terreno de la Web 3.0. El notariado respondió con formación y producción propia.}

% ------------------------------------------------------------
\chapter{Firma electrónica, firma digital y el rol del escribano}
\label{cap:firma}

La distinción entre firma electrónica y firma digital determina la validez y la fuerza probatoria de un documento: puede ser la diferencia entre un contrato válido y un instrumento que nadie puede hacer valer judicialmente. Este capítulo estudia esa distinción y sus consecuencias en la carga de la prueba, el valor agregado de la certificación notarial, la infraestructura y las plataformas disponibles, los riesgos de la identidad digital y los límites jurisdiccionales de la actuación notarial.

Para quien ejerce la función notarial, la distinción orienta el asesoramiento sobre qué tipo de firma corresponde según el riesgo del acto; reconocer mecanismos de fraude de identidad mediante \emph{deepfakes} o suplantación en videollamadas es una competencia cada vez más determinante; e incorporar firma digital, certificación remota y validación de documentos constituye una ventaja competitiva real.

\section{Punto de partida}

El déficit de formación específica en firma digital es un problema estructural: en general los programas de acceso a la función notarial no incluyen el estudio de la Ley 25.506, ni siquiera en distritos con mayor desarrollo institucional, como la Ciudad de Buenos Aires. Esto favorece que el tema se conozca de manera superficial y se confunda con el régimen de la firma electrónica.

\begin{important}
Si el acceso a la función notarial nunca exigió siquiera una lectura de la Ley de Firma Digital, el propio ingreso a la profesión queda desactualizado frente al mundo de los documentos electrónicos que ya circula en la práctica diaria.
\end{important}

\section{Firma electrónica y firma digital}
\label{sec:firmas}

La Ley 25.506 reconoce validez a la firma electrónica, aunque con un régimen probatorio distinto al de la firma digital.

\begin{formula}{Artículo 5, Ley 25.506 --- Firma electrónica}
Se entiende por firma electrónica al conjunto de datos electrónicos integrados, ligados o asociados de manera lógica a otros datos electrónicos, utilizado por el signatario como su medio de identificación, que carezca de alguno de los requisitos legales para ser considerada firma digital. En caso de ser desconocida la firma electrónica corresponde a quien la invoca acreditar su validez.
\end{formula}

\begin{formula}{Artículo 2, Ley 25.506 --- Firma digital}
Se entiende por firma digital al resultado de aplicar a un documento digital un procedimiento matemático que requiere información de exclusivo conocimiento del firmante, encontrándose ésta bajo su absoluto control, y que además sea susceptible de verificación por terceros.
\end{formula}

Lo determinante para que una firma sea «digital» y no simplemente «electrónica» no es la sofisticación del procedimiento matemático o criptográfico, sino quién está detrás de la infraestructura que emite el certificado.

\begin{important}
Si la firma fue otorgada por la ONTI (Oficina Nacional de Tecnologías de la Información) o por un certificador licenciado por ella ---como los Colegios de Escribanos que prestan ese servicio--- es firma digital. Todo lo demás, sin importar cuán sofisticada sea la tecnología subyacente, es firma electrónica.
\end{important}

\subsection{El caso del pase de un futbolista}

Para poner en tensión la distinción se plantea un caso real: un contrato de pase de un futbolista por una cifra muy significativa, firmado a través de una plataforma de firma electrónica ---no digital---, por un monto ejemplificado entre cien mil y trescientos mil dólares. ¿El contrato es válido? ¿Se considera firmado? ¿Adolece de algún vicio de forma?

\begin{example}{El pase del futbolista}
Un club y un jugador firman un contrato de pase por una suma muy importante utilizando firma electrónica. Al mes siguiente, el club no deposita el sueldo pactado. El jugador reclama; el club responde que no existe «ningún contrato firmado». El conflicto se traslada a una pregunta técnica: ¿ese documento está firmado en sentido jurídico, y quién tiene que probarlo?
\end{example}

Una postura posible, sostenida como ejercicio de debate, es que al no cumplir con los requisitos de la firma digital el instrumento debe considerarse un instrumento particular no firmado, aun cuando exprese una manifestación de voluntad. Sobre esa base se profundiza en los arts.~7 y 8 de la Ley 25.506, que instituyen la presunción de autoría y la presunción de integridad, ambas exclusivas de la firma digital.

\begin{formula}{Artículos 7 y 8, Ley 25.506}
\textbf{Artículo 7 (presunción de autoría):} se presume que un documento digital firmado digitalmente pertenece al titular del certificado digital que permite verificar dicha firma.

\textbf{Artículo 8 (presunción de integridad):} si el resultado de la verificación de una firma digital es válido, se presume que el documento digital no ha sido modificado desde el momento de su firma.
\end{formula}

La clave no está solo en si el documento es válido en abstracto, sino en sobre quién recae la carga de la prueba: en la firma electrónica, quien afirma haber firmado debe probarlo; en la firma digital opera una presunción legal a favor de quien firmó, y la contraparte debe probar que no lo hizo.

\begin{important}
La diferencia no es meramente técnica: define quién tiene que salir a probar algo en un conflicto judicial. Firmar electrónicamente traslada la carga de la prueba a quien firmó; firmar digitalmente la traslada a quien pretende desconocer la firma.
\end{important}

\subsection{Tensión entre la Ley de Firma Digital y el Código Civil y Comercial}

El estudio de la tensión entre la Ley de Firma Digital y la redacción del CCyC de 2015 respecto de la distinción entre firma electrónica y firma digital parte de una pregunta: ¿la firma electrónica es, jurídicamente, firma? Los arts.~286, 287 y 288 CCyC regulan la forma y la firma; los arts.~317 y 319 se vinculan con la valoración de la prueba electrónica por parte del juez. Como actividad de estudio se propuso elaborar dos columnas comparativas entre ambos cuerpos normativos, identificando la tensión conceptual que emerge de ellos.

\subsection{Tres niveles}

Puede construirse una estructura de tres niveles: en la base, la \textbf{firma electrónica}, reconocida por la ley y a la que no puede negársele validez sin negar buena parte del mundo digital cotidiano (por ejemplo, una comunicación de WhatsApp, que lleva inserta una firma electrónica en tanto vincula al usuario con su número telefónico); en un nivel intermedio, la \textbf{firma digital}, con sus presunciones legales de autoría e integridad; y en el nivel superior, la \textbf{certificación notarial de la firma digital}, que aporta certeza sobre la identidad de quien firmó.

\section{Presunciones, hash y riesgos de la identidad digital}

\subsection{La función hash}

La presunción de integridad se sostiene técnicamente en la función hash. Al firmar digitalmente un documento, el sistema aplica un algoritmo matemático que produce una cadena de caracteres única asociada a ese contenido exacto. Si el documento se modifica, aunque sea mínimamente, la función hash arroja un resultado distinto, lo que revela la alteración. La validación consiste en verificar que el hash generado por la clave pública coincide con el generado por la clave privada del firmante.

\begin{example}{Cómo funciona el «tilde verde»}
La clave privada genera un código numérico de 64 caracteres (letras y números). Cuando el sistema verifica la firma con la clave pública correspondiente, si el resultado coincide exactamente con ese mismo código, en el mismo orden, el validador muestra un tilde verde: el documento no fue alterado y existen la presunción de autoría y de integridad que otorga la ley.
\end{example}

\begin{important}
La presunción legal de autoría e integridad es una herramienta jurídica poderosa, pero no sustituye a la certeza sobre la identidad de la persona física que está detrás de la firma. Presunción y certeza son cosas distintas.
\end{important}

\subsection{Valor probatorio de las comunicaciones electrónicas}

Se ha demostrado en capacitaciones técnicas que es posible manipular el hash de una conversación de WhatsApp guardada en una computadora, lo que relativiza su valor probatorio fuera de un contexto pericial adecuado: da fe del procedimiento, pero no necesariamente del contenido de los mensajes, porque podría tratarse de una simulación. Los estados de WhatsApp presentan una complejidad adicional: son volátiles y desaparecen a las veinticuatro horas (véase el capítulo~\ref{cap:prueba}).

\subsection{Identidad digital: token, deepfakes y zero trust}

La existencia de una firma digital válida tampoco garantiza que el titular esté vivo y en condiciones de consentir al firmar: da una presunción, no una certeza, lo que habilita usos y abusos del token por terceros ---incluidos herederos--- una vez fallecido el titular. A ello se suman los fraudes de identidad mediante inteligencia artificial.

\begin{example}{Fraude por videollamada con identidades sintéticas}
En un caso ampliamente difundido, un empleado autorizó una transferencia de fondos muy importante tras una videollamada en la que creyó ver y escuchar a todo el directorio de su empresa. En realidad, todos los participantes ---salvo él--- eran entidades sintéticas generadas por inteligencia artificial. El trabajo remoto desde cualquier lugar del mundo también facilita que empresas crean estar contratando personas físicas cuando en verdad están dando acceso a espías industriales.
\end{example}

Corresponde aplicar el principio de \emph{zero trust} (cero confianza): no dar nunca por sentada la identidad de un interlocutor solo porque «se lo ve» a través de una pantalla, y requerir validación adicional cuando no se trate de un interlocutor conocido.

\subsection{La amenaza de la computación cuántica}

La computación cuántica amenaza a futuro los algoritmos criptográficos que hoy sostienen tanto la firma digital como la seguridad bancaria. El instituto estadounidense de normas y tecnología (NIST) convocó a matemáticos y físicos de todo el mundo para desarrollar algoritmos más resistentes.

\begin{example}{El «Quantum Day» y los algoritmos postcuánticos}
Se denomina «quantum day» al momento en que la computación cuántica alcance la capacidad de quebrar, en pocas horas, los algoritmos que hoy protegen las transacciones bancarias y la firma digital. Mientras tanto ya existen atacantes que aplican la estrategia de «cosechar hoy para abrir mañana»: recopilar y almacenar hoy información cifrada que no pueden leer, para descifrarla en el futuro cuando la tecnología cuántica lo permita.
\end{example}

\begin{important}
El sistema de firma digital que hoy se explica como una novedad podría quedar superado en pocos años. Bancos y organismos internacionales ya preparan algoritmos postcuánticos, mientras que en la práctica notarial argentina todavía se está incorporando el conocimiento básico de la firma digital vigente.
\end{important}

\section{El escribano como certificador}

\subsection{Certificación notarial de firma digital y prueba de vida}

Además de la firma digital simple, que ya trae presunción legal de autoría e integridad, existe la certificación notarial, un escalón adicional que aporta lo que la sola presunción no puede dar: certeza sobre la identidad de quien firmó en ese momento. El procedimiento típico de certificación remota consiste en que el escribano contacta al titular por videoconferencia a través de la plataforma segura del Colegio de Escribanos; la persona coloca su token y su contraseña, firma frente al escribano y este emite una constancia ---una foja digital--- certificando que esa persona colocó su firma en su presencia.

\begin{formula}{Artículo 308 --- Certificación notarial por videoconferencia}
El artículo 308 constituye la base normativa para la certificación notarial de firmas mediante videoconferencia, procedimiento que todos los escribanos deben poder ejecutar con solvencia.
\end{formula}

El elemento diferencial es la \textbf{prueba de vida}, requisito que se volvió especialmente relevante durante la pandemia, cuando había que constatar la subsistencia de personas mayores para sostener determinados beneficios sin poder visitarlas.

\begin{example}{Prueba de vida en videoconferencia}
Un sistema automatizado del Colegio envía a la persona un enlace mediante el cual debe grabar un video breve, pronunciando tres palabras aleatorias generadas en el instante, de modo que no puedan prepararse de antemano. Ese video se contrasta contra la base del Registro Nacional de las Personas (Renaper), que devuelve dos resultados: si la persona está efectivamente viva (no hay aviso de fallecimiento) y un porcentaje de coincidencia entre el rostro grabado y el de la base.
\end{example}

Este mecanismo permite detectar cuando alguien ya fallecido «sigue firmando» a través de su token. Al mismo tiempo, el sistema de validación puede resultar poco accesible para personas mayores, tensión que se trabaja bajo el concepto de «apoyo digital», en paralelo con los sistemas de apoyo que reconoce el CCyC. Cuando la persona no puede completar la prueba de vida por videoconferencia existe una validación intermedia: se le solicita, por el mismo canal, que envíe fotografías del frente y dorso de su documento de identidad.

\begin{example}{El documento desactualizado}
Si la persona presenta una fotografía del documento con una imagen anterior a la última actualización, corresponde exigirle que primero renueve el documento antes de continuar, porque la fotografía válida para la comparación es siempre la última registrada. Es frecuente entre personas mayores que, por temor a perder o que les roben el documento actualizado, prefieren salir con uno anterior.
\end{example}

\subsection{Funciones del escribano: identidad, legalidad y fisco}

Al certificar una firma digital el escribano cumple funciones que ninguna plataforma puede reemplazar por sí sola.

\begin{important}
El escribano es, en primer lugar, una entidad de identidad: verifica que quien firma es quien dice ser. En segundo lugar, ejerce un control de legalidad sobre el contenido de lo que se firma. En tercer lugar, actúa como agente de información, percepción y retención en materia tributaria. Y cumple funciones de prevención de lavado de activos, financiamiento del terrorismo y de la proliferación de armas de destrucción masiva (capítulo~\ref{cap:pbc}).
\end{important}

Sobre el control de legalidad, existe un riesgo concreto: personas que descargan de internet modelos de contratos ---por ejemplo, de locación--- que citan legislación derogada o de otra jurisdicción; el escribano realiza ese contralor. Respecto del control tributario, si el acto debe estar gravado con algún impuesto, corresponde dejar constancia para resguardar la responsabilidad profesional y asesorar correctamente. En cuanto a la prevención, el escribano avala negocios que podrían encubrir delitos precedentes ---como la trata de personas--- o delitos informáticos, por lo que ese control tiene un valor social que excede la certificación técnica de una firma.

\begin{synthesis}
Presunción legal (firma digital), prueba de quién firmó (firma electrónica) y certeza aportada por la certificación notarial son tres niveles distintos y complementarios, no intercambiables. Ninguna plataforma privada de firma electrónica realiza, además, control de identidad reforzado, control de legalidad, control tributario y prevención de lavado de activos en un mismo acto.
\end{synthesis}

\section{Infraestructura y plataformas de firma digital}

\subsection{ONTI, certificadores licenciados y el token}

El procedimiento operativo para obtener el token es: la persona solicita un turno, se presenta ante un operador del Colegio y carga su clave y su contraseña ---la única que conocerá, bajo su exclusiva responsabilidad---, quedando vinculada a un token.

\begin{example}{Clave simétrica frente a clave asimétrica}
En una clave simétrica, alguien más conoce o entrega la misma clave que usa el titular. En el esquema de firma digital, en cambio, el operador entrega el token, pero la contraseña que se carga en ese momento la conoce exclusivamente el titular: es una clave asimétrica, y por eso la responsabilidad de lo firmado recae enteramente en esa persona.
\end{example}

Si el token se pierde o se daña, el titular debe comunicarse con el Colegio para solicitar la revocación del certificado y la emisión de un nuevo token con un nuevo turno. En el uso cotidiano, la computadora solicita la clave y contraseña, lo que habilita el uso de un programa llamado «firmador», al cual se adjunta el documento; al hacer clic en «firmar» y verificar el tilde verde, el documento queda firmado y se descarga junto con su certificado de registro.

\subsection{Validación de documentos y percepción social}

Quien recibe un documento firmado digitalmente no necesita ninguna contraseña: accede a una herramienta de validación provista por la entidad emisora, a la que adjunta el documento para confirmar su autenticidad.

\begin{example}{El validador de documentos notariales (MADONO)}
El Colegio de Escribanos ofrece en su sitio web un validador de documentos notariales (MADONO). Quien recibe un documento firmado digitalmente por un escribano debe adjuntarlo en esa herramienta para confirmar que efectivamente fue firmado y certificado por esa persona; de lo contrario cualquiera podría fabricar un PDF con apariencia de estar firmado sin firma real detrás. El mecanismo es comparable al de los títulos digitales de propiedad automotor, que incluyen un código QR para su validación.
\end{example}

La validación por el receptor cierra el círculo de la infraestructura: la primera pata es no divulgar nunca la clave y contraseña; la segunda, la firma en sí; y la tercera, la validación por quien recibe el documento.

Existe, además, un problema más cultural que técnico: buena parte de la ciudadanía no confía en un documento firmado digitalmente porque a simple vista no «parece» firmado, al no haber rúbrica manuscrita ni sello visible. La solución del Colegio fue incorporar un sello visual ---una imagen con la leyenda de que el documento fue firmado digitalmente por determinado escribano en determinada fecha--- al pie de cada documento digital.

\begin{important}
La validez jurídica de un documento firmado digitalmente no depende de que se «vea» firmado, pero su aceptación social sí. Diseñar la experiencia del usuario alrededor del documento digital, con sellos visuales y validadores accesibles, es tan importante como la solidez técnica y legal.
\end{important}

\subsection{El mercado de la firma electrónica y la plataforma Firma Ciudadana}

A partir de diciembre de 2025 una resolución habilitó la tramitación de la firma digital de manera remota, sin presencialidad ante la ONTI o el certificador licenciado, lo que generó dudas sobre cómo se controla la identidad en ese trámite a distancia. Como respuesta institucional se presentó recientemente la plataforma «Firma Ciudadana» del Colegio de Escribanos, lanzada en el marco de un congreso de \emph{real estate}.

\begin{example}{La plataforma Firma Ciudadana}
La persona se registra («se enrola»), inicia sesión y firma un documento a través del proceso que ofrece la plataforma. A diferencia de otros sistemas, el documento firmado no queda alojado en ningún servidor del Colegio, lo que garantiza la privacidad. La firma es gratuita y está avalada por una institución con más de 160 años de trayectoria en el resguardo de la seguridad jurídica.
\end{example}

\begin{example}{El informe de 2021 sobre empresas de firma electrónica}
El relevamiento de 2021 constató que el número de empresas privadas de firma electrónica había pasado de tres a veintiuna en poco tiempo, mientras que las ventas de fojas de certificación de firmas y de libros de escribanos venían en descenso. Ese cruce de datos motivó la discusión política dentro del Colegio sobre cómo reposicionar la función notarial.
\end{example}

La conclusión fue que no tenía sentido resistirse a la desaparición de una modalidad de certificación para preservar un statu quo, sino competir ofreciendo un servicio superior y gratuito, apoyado en la confianza que la ciudadanía deposita en la institución notarial. La plataforma busca captar a la ciudadanía que de otro modo recurriría a plataformas privadas de pago, ofreciéndole certeza sobre la identidad, control de legalidad, asesoramiento y custodia frente a delitos precedentes, elementos que una empresa de firma electrónica no ofrece. Esta lectura enlaza con la idea, presentada en la sección~\ref{sec:forma}, de que lo relevante es la función que cumple el acto y no su soporte.

\section{Jurisdicción notarial y derecho comparado}

\subsection{Competencia territorial y digitalización comparada}

La plataforma de certificaciones que reemplaza al histórico «libro azul» en soporte físico es el Sistema de Certificaciones de Firmas (SDF), actualmente en su versión 2.0. Permite que la persona requirente se encuentre en cualquier lugar, siempre que el escribano actuante esté sentado dentro de su propia jurisdicción, que es la que determina la competencia.

\begin{example}{La evolución del criterio jurisdiccional}
En una primera etapa el sistema habilitaba la certificación remota únicamente para personas ubicadas en la Ciudad de Buenos Aires o en el extranjero, y se resolvía mediante una declaración jurada de la propia persona sobre su ubicación. Frente a las tensiones interprovinciales que generó ---con amenazas de denuncia contra escribanos que certificaban firmas de personas ubicadas fuera de su jurisdicción de residencia---, el criterio cambió: hoy no importa dónde esté físicamente la persona requirente en el país; si algún colegio provincial quisiera cuestionar una certificación remota, la respuesta institucional fue ofrecerle el mismo sistema, gratuitamente, para que también pueda certificar a distancia.
\end{example}

\begin{table}[H]
\centering
\caption{Digitalización notarial: Argentina y Brasil}
\renewcommand{\arraystretch}{1.2}
\begin{tabularx}{\textwidth}{>{\raggedright\arraybackslash}p{0.28\textwidth}X}
\toprule
\textbf{Aspecto} & \textbf{Situación} \\
\midrule
Grado de digitalización & Brasil alcanzó el cien por ciento de digitalización de sus actos notariales; en la Argentina, la digitalización avanzó de manera parcial. \\
Diseño institucional & Brasil es una república federativa en la que notarios y registradores pueden ejercer en cualquier punto del territorio nacional. \\
Jurisdicción notarial & En la Argentina, cada provincia se reservó constitucionalmente la facultad de regular la función notarial en su territorio, lo que estableció jurisdicciones estrictas: un escribano matriculado en una jurisdicción no puede ejercer en otra sin cumplir los requisitos de esa demarcación. \\
\bottomrule
\end{tabularx}
\end{table}

La brecha de digitalización no responde únicamente a una diferencia de infraestructura tecnológica, sino a una diferencia estructural de organización federal: la fragmentación jurisdiccional argentina es un obstáculo institucional adicional a la digitalización plena.

\subsection{Actuación notarial internacional y matricidad en soporte papel}

La actuación notarial internacional plantea cómo intervenir cuando las partes de un mismo acto ---por ejemplo, una compraventa de inmueble--- se encuentran en distintos países. El principio general es que no puede recurrirse a un notario totalmente ajeno a la operación ni buscar una jurisdicción más conveniente a nivel tributario si ni el inmueble ni las partes tienen relación con ese territorio.

\begin{definition}{Punto de conexión}
Noción tomada del derecho internacional privado según la cual debe existir un vínculo real entre el notario interviniente y el negocio jurídico. Sin ese punto de conexión no existe posibilidad de intervención notarial válida.
\end{definition}

\begin{example}{El caso de Puerto Rico}
La legislación de Puerto Rico permite que un notario tome un avión, se lleve su protocolo y viaje ---por ejemplo, a Ámsterdam--- para tomarle la firma a un residente en el extranjero, algo que la organización notarial argentina no admite bajo su esquema actual de matricidad y jurisdicción.
\end{example}

La escritura pública en la Argentina continúa firmándose en soporte papel y de manera presencial, sin posibilidad todavía de escritura a distancia. Es posible trabajar con poderes otorgados en el extranjero como instrumento privado con doble columna, siempre que el país donde vaya a surtir efecto admita esa modalidad; pero cuando el apoderamiento debe otorgarse en la Argentina el mandato exige escritura pública. Ya existen, sin embargo, las primeras copias digitales de las escrituras: aunque la matriz sigue en soporte papel, la persona requirente puede llevarse en el momento una copia digital, sin esperar los tiempos administrativos tradicionales.

\begin{example}{Poderes con sucursales en otra jurisdicción}
Una empresa con sede en la Ciudad de Buenos Aires necesita otorgar un poder a una persona que actuará en su sucursal de otra provincia. Suele ser más práctico otorgar la primera copia, legalizarla en el Colegio de Escribanos correspondiente y remitirla ---incluso físicamente, por correo--- a la jurisdicción donde va a utilizarse, en lugar de forzar una certificación digital que no siempre resuelve estos cruces interjurisdiccionales.
\end{example}

\begin{important}
Pese al desarrollo de la firma digital, la certificación remota y la certeza notarial, la escritura pública argentina continúa firmándose en soporte papel y de manera presencial. La digitalización avanzó en la certificación de firmas y en las copias, pero la matricidad del acto de mayor jerarquía todavía no dio ese salto.
\end{important}

\cornell{Firma electrónica, firma digital y rol del escribano}{
\cq{¿Qué distingue a la firma digital de la electrónica?} & Es digital la firma otorgada por la ONTI o un certificador licenciado; el resto, aunque sea sofisticado, es electrónica. Solo la digital tiene presunciones de autoría (art.~7) e integridad (art.~8). \\ \hline
\cq{¿Quién debe probar la firma en cada caso?} & Firma electrónica desconocida: quien la invoca debe acreditar su validez. Firma digital: quien la desconoce debe probar que no firmó. \\ \hline
\cq{¿Qué significa el tilde verde?} & Que el hash generado con la clave pública coincide con el de la clave privada: el documento no fue alterado. \\ \hline
\cq{¿Por qué la presunción no equivale a certeza?} & No garantiza que el titular esté vivo o sea quien firma: hay abuso de token, \emph{deepfakes} y manipulación de comunicaciones. \\ \hline
\cq{¿Qué aporta la certificación notarial?} & Certeza de identidad (con prueba de vida por videoconferencia y contraste con Renaper), control de legalidad, control tributario y prevención de lavado. \\ \hline
\cq{¿Qué pasos completan el círculo de la firma digital?} & No divulgar clave y contraseña, firmar y validar por el receptor (MADONO). \\ \hline
\cq{¿Qué límite mantiene la escritura pública?} & Se firma en papel y presencialmente; hay copias digitales, pero no escritura a distancia. \\ \hline
\cq{¿Por qué Brasil está más digitalizado?} & Por diseño institucional federal que permite ejercer en todo el territorio; en Argentina cada provincia regula su función notarial. \\ \hline
}{La firma electrónica y la digital difieren en presunciones y carga de la prueba; la certificación notarial agrega certeza de identidad y controles que ninguna plataforma privada ofrece. La infraestructura (token, validación, plataformas del Colegio) convive con riesgos de identidad y con límites jurisdiccionales y de soporte.}

% ------------------------------------------------------------
\chapter{Prueba electrónica y contratación electrónica}
\label{cap:prueba}

La digitalización trasladó al ámbito electrónico tanto los hechos que hay que probar como los contratos que se celebran. Este capítulo estudia cómo se prueba en el entorno digital (actas notariales, metadatos, peritos, cadena de custodia) y cómo se regula la contratación electrónica y de consumo en el CCyC.

\section{Nuevas actas notariales en redes sociales}

Desde los años 2007-2008 comenzaron a labrarse las primeras actas de comprobación de situaciones dentro de redes sociales.

\begin{example}{Aumento de cuota alimentaria y fotos de vacaciones}
En un proceso de aumento de cuota alimentaria, la parte demandada sostenía no poder afrontar el incremento, mientras que su perfil de Facebook mostraba fotografías de vacaciones. Las contradicciones entre lo declarado judicialmente y lo exhibido en redes sociales empezaron a requerir un nuevo tipo de acta notarial.
\end{example}

\section{La prueba electrónica}

\begin{definition}{Prueba electrónica}
Nueva rama del derecho, vinculada al derecho procesal electrónico, que se ocupa de hechos que ya no se prueban en el ámbito físico sino en el digital (redes sociales, comunicaciones electrónicas).
\end{definition}

Muchas cuestiones probatorias dejaron de darse únicamente en el ámbito físico. Es allí donde el notario tiene que aprender a labrar este tipo de actas: se inicia una nueva etapa probatoria.

\begin{example}{El estado de WhatsApp}
Un empleado publicó, un domingo a la noche, un estado de WhatsApp en el que anunciaba que al día siguiente no se presentaría a trabajar por motivos ajenos a una causa habilitante. Como el estado desaparece a las 24 horas, fue necesario realizar una captura y una copia forense antes de que se perdiera, para incorporarla como prueba al procedimiento laboral.
\end{example}

Este tipo de prueba es cada vez más relevante en materia laboral y de familia, donde la prueba admite mayor amplitud y el WhatsApp suele ser uno de los primeros elementos ofrecidos, lo que justifica que el escribano se especialice también en derecho informático para realizar constataciones con respaldo forense adecuado. Para todo operador jurídico que quiera comprender cabalmente este campo se consideró necesaria la especialización en derecho informático, dado que la formación de grado resulta insuficiente.

\section{Metadatos}

Toda comunicación electrónica, además de la comunicación en sí, tiene aspectos técnicos: hasta un documento de Word posee metadatos.

\begin{definition}{Metadatos}
Información técnica que acompaña a toda comunicación o documento electrónico, más allá de su contenido visible. Es la materia prima que necesita cualquier perito informático.
\end{definition}

\section{Límites del conocimiento técnico del notario: el perito}

\begin{example}{La cañería rota y el perito}
Ante la constatación de una cañería rota, el notario no tiene por qué saber si es pluvial o cloacal ni diagnosticar la causa técnica de la fisura: debe convocar a un arquitecto, ingeniero civil o especialista, quien elabora el informe técnico. El notario da fe de la manifestación y deja constancia de lo actuado, lo que constituye plena prueba en sede civil.

La misma lógica se traslada al ámbito digital: si el notario no tiene conocimientos técnicos informáticos, debe basarse en el dictamen de un perito o experto forense, dejando constancia de los pasos seguidos por ese experto para la obtención, custodia y presentación de la prueba en sede judicial.
\end{example}

\section{Cadena de custodia y volatilidad de la prueba digital}

\begin{definition}{Cadena de custodia en materia digital}
Conjunto de pasos documentados que garantizan que la prueba electrónica obtenida no fue alterada desde su recolección hasta su presentación judicial.
\end{definition}

Su preservación es especialmente crítica porque la prueba electrónica es mucho más volátil que la física: en el ámbito físico, la prueba en papel sustraída reaparece para declarar; en el digital ocurre algo equivalente, pero con mucha mayor volatilidad, ya que pueden introducirse virus o borrarse la prueba almacenada. Sobre la manipulación del hash de comunicaciones guardadas, véase la sección~\ref{sec:firmas}.

\section{Contratación electrónica y derecho del consumidor digital}

Cerrando el desarrollo de la Web 2.0 corresponde considerar la \textbf{plataformización de la economía}, los términos y condiciones en las plataformas y el derecho de consumo digital aplicado a las contrataciones electrónicas. En el año 2015 se incorporó al CCyC normativa de derecho de consumo aplicada a la contratación electrónica. Los arts.~286, 287 y 288 regulan la forma y la firma; los arts.~317 y 319 se vinculan con la valoración de la prueba electrónica por parte del juez.

\begin{exercise}{El pase de un jugador firmado electrónicamente}
Un documento de pase de un jugador de fútbol por una suma relevante fue firmado a través de una plataforma de firma electrónica. ¿Ese contrato es válido? ¿Se considera firmado? ¿Adolece de algún vicio de forma? Resuélvase a partir de la distinción entre firma electrónica y firma digital (sección~\ref{sec:firmas}).
\end{exercise}

\cornell{Prueba electrónica y contratación electrónica}{
\cq{¿Qué es la prueba electrónica?} & Rama vinculada al derecho procesal electrónico que trata hechos probados en el ámbito digital. \\ \hline
\cq{¿Qué son los metadatos?} & Información técnica que acompaña a todo documento o comunicación; materia prima del perito informático. \\ \hline
\cq{¿Cuál es el límite del notario?} & No debe suplir conocimientos técnicos: se apoya en un perito y deja constancia de sus pasos. \\ \hline
\cq{¿Qué es la cadena de custodia?} & Pasos documentados que garantizan que la prueba no fue alterada desde su recolección hasta su presentación. \\ \hline
\cq{¿Qué normas regulan la contratación electrónica en el CCyC?} & Arts.~286, 287 y 288 (forma y firma) y 317 y 319 (valoración de la prueba electrónica). \\ \hline
}{La prueba digital es volátil y exige actas notariales, peritos y cadena de custodia. La contratación electrónica y de consumo se regula en el CCyC y se conecta con la distinción entre firma electrónica y firma digital.}

% ============================================================
% PARTE VII — BLOCKCHAIN, CRIPTOACTIVOS Y TOKENIZACIÓN
% ============================================================
\part{Blockchain, criptoactivos y tokenización inmobiliaria}
\label{parte:blockchain}

% ------------------------------------------------------------
\chapter{Del quiebre de la confianza a Bitcoin}
\label{cap:bitcoin}

Con la Web 3.0 internet deja de ser solo un canal de información y pasa a ser un medio para reservar y transferir valor entre pares, sin intermediarios centralizados. Para comprender cómo se llegó a ello este capítulo reconstruye el contexto: la pérdida de confianza en el dinero estatal, el pensamiento cypherpunk, el problema del doble gasto y el surgimiento de Bitcoin. La economía digital se estudia a partir de tres ejes que no deben pensarse por separado: el \textbf{tecnológico}, el \textbf{social} y el \textbf{económico}, que en conjunto explican cómo se administra el dinero, cómo se conserva la riqueza y hacia qué espacios se desplaza hoy esa riqueza.

\section{Confianza, dinero y crisis}

El primer gran quiebre de confianza en el sistema monetario internacional se ubica en \textbf{1971}, cuando el presidente estadounidense Richard Nixon suspendió la convertibilidad del dólar en oro, poniendo fin al patrón oro vigente desde los acuerdos de Bretton Woods. Se lo señala como el «principio del fin» de la confianza plena en la moneda estatal. La \textbf{crisis de las hipotecas subprime de 2008} en Estados Unidos generó luego una pérdida de confianza generalizada en el sistema financiero, terreno fértil en el que germinaron las criptomonedas como alternativa de reserva de valor. Se comenzó a prestar dinero sin respaldo, lo que generó una burbuja inmobiliaria que estalló; la crisis griega, vinculada también a préstamos que no pudieron cubrirse, es otro antecedente.

\begin{table}[H]
\centering
\caption{Línea de tiempo de los hitos mencionados}
\renewcommand{\arraystretch}{1.2}
\begin{tabularx}{\textwidth}{>{\raggedright\arraybackslash}p{0.25\textwidth}X}
\toprule
\textbf{Momento} & \textbf{Hito} \\
\midrule
Década de 1960 & Desarrollo de ARPANET, la primera red de comunicación a distancia. \\
1971 & Suspensión de la convertibilidad del dólar en oro y fin del patrón oro. \\
2000 & Leyes 25.506 (firma digital) y 25.326 (datos personales). \\
2008 & Crisis de las hipotecas subprime; white paper de Bitcoin. \\
\bottomrule
\end{tabularx}
\end{table}

\begin{example}{El rescate estatal y el origen de la desconfianza}
Para evitar que Estados Unidos entrara en bancarrota tras el estallido de la burbuja inmobiliaria, el Estado no dispuso de los ahorros privados de las personas ---cosa que en la Argentina sí se permitió en otros contextos históricos, como el corralito--- sino que debió recurrir a fondos públicos para rescatar el sistema financiero. Esta decisión, que genera más impuestos y déficit, fue el caldo de cultivo del movimiento que cuestionaba la intervención estatal en la economía y que impulsaría, poco después, la creación de Bitcoin. La pregunta de quién paga las burbujas es central en este razonamiento.
\end{example}

\section{Redes centralizadas y redes distribuidas}
\label{sec:redes}

Una red centralizada organiza su información en un único sistema central: si ese servidor es atacado, cae todo el sistema. Una red distribuida ubica sus servidores en distintos lugares: si se ataca a uno, los demás siguen funcionando y la organización resguarda su información.

\begin{table}[H]
\centering
\caption{Redes centralizadas y redes distribuidas}
\renewcommand{\arraystretch}{1.2}
\begin{tabularx}{\textwidth}{>{\raggedright\arraybackslash}p{0.26\textwidth}XX}
\toprule
\textbf{Aspecto} & \textbf{Red centralizada} & \textbf{Red distribuida} \\
\midrule
Ubicación de los servidores & Un único sistema central. & Servidores distribuidos en diferentes lugares. \\
Ante un ataque & Si se ataca la red, cae todo el sistema. & Si se ataca un servidor, los otros siguen funcionando. \\
Información de la organización & Depende del funcionamiento del punto central. & La organización logra resguardar la información. \\
\bottomrule
\end{tabularx}
\end{table}

\section{Los cypherpunks y la desconfianza en el Estado}

Surgió un grupo filosófico de fuerte base técnica y criptográfica, los \textbf{cypherpunks}, que cuestionó el manejo estatal del dinero señalando fenómenos como la inflación y las burbujas económicas. Su origen se vincula con la economía liberal inglesa más purista, o con una economía liberal libertaria cuya postura es que no se necesita al Estado; lo integran además programadores con dominio de los lenguajes con que se programa blockchain. La tecnología blockchain nace de este pensamiento, en un grupo de técnicos que sostienen que los Estados llevan a la bancarrota a las personas cuyo patrimonio fluctúa.

\begin{important}
Si las redes distribuidas ya resguardan la información de organizaciones y empresas, también pueden permitir manejar el patrimonio propio sin depender del Estado y sin las burbujas inflacionarias que genera la emisión centralizada de moneda.
\end{important}

\begin{note}
% TODO: verificar consistencia entre fuentes originales. Una fuente afirma que Solidity es «el lenguaje bajo el cual se programa blockchain»; la propia fuente aclara que, en rigor, es el lenguaje de los contratos inteligentes de Ethereum y que Bitcoin usa un lenguaje de \emph{scripts} propio y más limitado.
En la exposición se menciona a Solidity como el lenguaje bajo el cual se programa blockchain. Las notas de origen aclaran que es, en rigor, el lenguaje de los contratos inteligentes de la red Ethereum, mientras que Bitcoin utiliza un lenguaje de \emph{scripts} propio y más limitado.
\end{note}

\section{El doble gasto y el registro distribuido}

\begin{definition}{Doble gasto}
Problema de certificar que el dinero ya gastado no pueda volver a gastarse. Tradicionalmente lo resuelve el banco, en su carácter de autoridad central.
\end{definition}

\begin{example}{El home banking}
En un sistema bancario tradicional, si una persona tiene diez unidades y gasta una, la certificación de que le quedan nueve está centralizada en una única autoridad: el banco. Cuando el banco realiza tareas de mantenimiento ---por ejemplo, los avisos de indisponibilidad los fines de semana por la madrugada--- el usuario no puede disponer de sus fondos mientras dura la interrupción.
\end{example}

El planteo cypherpunk propone que, si todos los participantes de una red se ponen de acuerdo, pueden cumplir la función del banco sin autoridad central: si A tiene diez unidades y transfiere una a B, todos los integrantes de la red saben al mismo tiempo que A tiene nueve y B tiene una. La información deja de estar centralizada y se comparte en tiempo real entre todos los miembros.

\begin{table}[H]
\centering
\caption{El banco y el registro distribuido}
\renewcommand{\arraystretch}{1.2}
\begin{tabularx}{\textwidth}{>{\raggedright\arraybackslash}p{0.24\textwidth}XX}
\toprule
\textbf{Aspecto} & \textbf{Banco (centralizado)} & \textbf{Registro distribuido} \\
\midrule
Quién certifica el pasaje de fondos & Una sola autoridad: el banco. & Todos los que forman parte de la red se ponen de acuerdo. \\
Disponibilidad & Si el banco actualiza su sistema y se detiene, no puedo usar mis fondos. & Todos tienen en tiempo real la misma información; no depende de un único punto. \\
\bottomrule
\end{tabularx}
\end{table}

\section{Satoshi Nakamoto y el white paper de Bitcoin}
\label{sec:whitepaper}

En ese contexto aparece, bajo el seudónimo de \textbf{Satoshi Nakamoto} ---cuya identidad real se desconoce, sin que esté establecido siquiera si es una persona o un grupo---, el \textbf{white paper de Bitcoin}, que describe un sistema \emph{peer to peer}: una transferencia de persona a persona, sin entidad centralizada externa que la certifique.

\begin{definition}{White paper}
Documento constitucional fundacional de una criptomoneda o de una plataforma que genera algún token o moneda propia. Funciona como el conjunto de reglas de juego del proyecto.
\end{definition}

\begin{important}
El white paper debe leerse siempre antes de considerar cualquier inversión. Frente a un cliente que quiere invertir en una criptomoneda presentada como un furor del mercado, lo primero que debe indicársele es que descargue y lea el white paper correspondiente.
\end{important}

\begin{example}{La analogía de la gamificación}
Así como en un videojuego existen reglas claras sobre qué se puede y qué no, el white paper establece las reglas de la plataforma: cuántos tokens se pueden emitir; cuánto tiempo debe permanecer un inversor; en qué condiciones puede retirarse; y si puede vender su posición a alguien externo al proyecto o solo dentro del mercado interno de la plataforma. Estas son las nuevas preguntas que los operadores jurídicos deben aprender a formular, porque no existe un marco jurídico único para estas contrataciones: son ofertas realizadas directamente por una plataforma a un particular, respecto de las cuales el rol del asesor es explicar dónde, cómo y cuándo se debe hacer clic para contratar.
\end{example}

\begin{table}[H]
\centering
\caption{Bitcoin y blockchain}
\renewcommand{\arraystretch}{1.2}
\begin{tabularx}{\textwidth}{>{\raggedright\arraybackslash}p{0.22\textwidth}X}
\toprule
\textbf{Concepto} & \textbf{Descripción} \\
\midrule
Bitcoin & La criptomoneda: una tenencia con sentido económico, con un componente de valor y de escasez fundamental. \\
Blockchain & La tecnología subyacente: se utiliza para transaccionar la criptomoneda, pero su potencial excede ampliamente a Bitcoin. \\
\bottomrule
\end{tabularx}
\end{table}

Quedarse solo con Bitcoin es quedarse con «un pedacito del potencial» de blockchain.

\section{Escasez, reserva de valor y confianza en el código}

El white paper establece que se emitirá una cantidad determinada de unidades, por lo que quienes ya poseen bitcoins forman parte de un sistema cerrado a nuevos ingresos gratuitos de oferta.

\begin{note}
% TODO: verificar consistencia entre fuentes originales. La fuente ubica el cierre de la emisión de Bitcoin en el año 2049 y aclara luego que el protocolo fija un tope de 21 millones de unidades cuya emisión completa está prevista alrededor del año 2140.
Se ubicó el cierre de la emisión de Bitcoin en el año 2049. Las notas de origen aclaran que el protocolo fija, en rigor, un tope de 21 millones de unidades, cuya emisión completa está prevista para alrededor del año 2140. Se consignan ambos datos.
\end{note}

\begin{example}{El club o la membresía}
La escasez de Bitcoin puede compararse con un club o una membresía: quienes ya pertenecen conservan su lugar, y quien desea ingresar debe pagar el precio que los miembros actuales estén dispuestos a aceptar, porque no hay nuevas plazas gratuitas.
\end{example}

Esta dinámica constituye la \textbf{reserva de valor}: la riqueza se acumula y adquiere con el tiempo un valor distinto del inicial, en función de la escasez y de la oferta y demanda. Todos los participantes cuidan la red para mantener ese valor, de modo que nadie miente ni genera un disvalor dentro de la cadena. Quien participa como nodo activo y se desconecta temporalmente debe aceptar, al reincorporarse, con un margen de probabilidad, las transacciones que no pudo verificar: hay que confiar en el código informático, que es quien otorga la seguridad criptográfica. El sistema funciona desde 2008, con fluctuaciones de mercado (se mencionaron valores de hasta cien mil dólares y descensos posteriores a alrededor de sesenta y siete mil).

\section{Coins, altcoins y contratos inteligentes}

Se denomina \textbf{coin} a Bitcoin y \textbf{altcoins} a todas las demás monedas alternativas, entre ellas Tether y Ethereum. Ethereum es la otra gran red posterior a Bitcoin: permite el ingreso a través de Layer 1 y Layer 2 y habilita \textbf{contratos inteligentes}, sobre los cuales corren las plataformas que permiten la tokenización.

No todas las criptomonedas conservan su valor. Algunas ---las \emph{shitcoins} o \emph{memecoins}--- fueron impulsadas artificialmente al alza, lo que llevó a numerosos inversores a comprarlas para luego desplomarse, generando pérdidas comparables a «perder como en la guerra». Los propios alumnos señalaron el caso de la criptomoneda Libra como ejemplo; otra fuente lo describe como un token lanzado sin registrarse ante la CNV, que generó desconfianza pública (capítulo~\ref{cap:tokenizacion}).

\cornell{Del quiebre de la confianza a Bitcoin}{
\cq{¿Qué hitos marcan la pérdida de confianza en el dinero estatal?} & 1971: fin de la convertibilidad del dólar en oro. 2008: crisis subprime. Ambos explican el terreno fértil de las criptomonedas como reserva de valor. \\ \hline
\cq{¿Quiénes son los cypherpunks?} & Grupo de fuerte base técnica y criptográfica, de pensamiento liberal libertario, que cuestiona el manejo estatal del dinero. \\ \hline
\cq{¿Qué es el doble gasto y quién lo resuelve tradicionalmente?} & Certificar que el dinero gastado no vuelva a gastarse; lo resuelve el banco. En blockchain, lo resuelve el acuerdo de todos los participantes. \\ \hline
\cq{¿Qué es el white paper?} & El documento fundacional y conjunto de reglas de una criptomoneda o plataforma; debe leerse antes de invertir. \\ \hline
\cq{¿Qué diferencia hay entre Bitcoin y blockchain?} & Bitcoin es la criptomoneda; blockchain, la tecnología subyacente cuyo potencial es mayor. \\ \hline
\cq{¿Qué es una altcoin?} & Toda moneda alternativa a Bitcoin (p.~ej. Tether, Ethereum). \\ \hline
}{La pérdida de confianza en el dinero estatal y el pensamiento cypherpunk dieron origen a Bitcoin, que resuelve el doble gasto con un registro distribuido. Bitcoin es la moneda y blockchain la tecnología; el white paper fija las reglas y debe leerse siempre antes de invertir.}

% ------------------------------------------------------------
\chapter{Blockchain: bloques, minería, billeteras e inteligencia artificial}
\label{cap:blockchain}

Este capítulo explica cómo funciona técnicamente blockchain (inmutabilidad, bloques, mineros, mecanismos de consenso), cómo se identifican y custodian los activos (exploradores, billeteras, exchanges) y qué debates plantean la concentración de poder y la inteligencia artificial.

\section{Inmutabilidad y descentralización}

\begin{definition}{Inmutabilidad}
Una transacción ya realizada y verificada por el sistema de mineros no puede alterarse. El respaldo de la información está descentralizado en múltiples ubicaciones y no existe una única copia que alguien pueda modificar; cualquier alteración sería detectada por los demás participantes.
\end{definition}

Esta característica sostiene la confianza de quienes transaccionan sus ahorros en la cadena en lugar de depositarlos en un banco: nadie puede mentir dentro del sistema porque todos conservan una copia verificable.

\section{Los mineros y el bloque}

\begin{example}{El Lejano Oeste}
La minería se apoya en la imagen de la fiebre del oro del Lejano Oeste norteamericano: de manera análoga, los «mineros» de blockchain ponen a trabajar computadoras de gran poder de cálculo para descubrir el algoritmo que permite cerrar y encriptar un conjunto de transacciones.
\end{example}

Las transacciones se agrupan, según el momento en que se producen, en un \textbf{bloque}. Una vez cerrado, no puede revertirse ni someterse a ingeniería inversa. El cierre tiene un ritmo estable, el \textbf{pulso de diez minutos} de Bitcoin: las máquinas prueban sucesivamente numerosas posibilidades hasta descubrir el algoritmo que sella el bloque.

\begin{table}[H]
\centering
\caption{El pulso de diez minutos y el ajuste de dificultad}
\renewcommand{\arraystretch}{1.2}
\begin{tabularx}{\textwidth}{>{\raggedright\arraybackslash}p{0.30\textwidth}>{\raggedright\arraybackslash}p{0.20\textwidth}X}
\toprule
\textbf{Situación} & \textbf{Significado} & \textbf{Qué se hace} \\
\midrule
El algoritmo tarda más de diez minutos & Está muy difícil. & Se baja el nivel de complejidad. \\
Se resuelve en menos de diez minutos (por ejemplo, en nueve) & Está muy fácil. & Se procura mantener la dificultad en torno a los diez minutos. \\
Objetivo del ajuste & Pulso estable de diez minutos. & Que todas las transacciones de ese lapso queden encriptadas y cerradas en un bloque. \\
\bottomrule
\end{tabularx}
\end{table}

\subsection{Mecanismos de consenso y paradoja de la centralización}

La principal crítica a la \textbf{minería de criptomonedas} es su huella ambiental y de carbono. El problema pudo resolverse en gran medida con el cambio del mecanismo de consenso de Ethereum, que migró de \emph{proof of work} a \emph{proof of stake}.

\begin{table}[H]
\centering
\caption{Mecanismos de consenso de una red blockchain}
\renewcommand{\arraystretch}{1.2}
\begin{tabularx}{\textwidth}{>{\raggedright\arraybackslash}p{0.20\textwidth}XX}
\toprule
& \textbf{Proof of work} & \textbf{Proof of stake} \\
\midrule
Denominación & Prueba de trabajo. & Prueba de participación. \\
Características & Alto consumo energético; requiere gran poder de cómputo (minería). & Reduce el consumo energético; se basa en la participación económica de los validadores. \\
Red de Ethereum & Sistema del que migró. & Sistema al que migró. \\
\bottomrule
\end{tabularx}
\end{table}

\begin{important}
\textbf{La paradoja de la descentralización.} La filosofía originaria busca que ningún actor concentre poder de decisión. Sin embargo, en la práctica grandes concentraciones de capacidad de minería (por ejemplo, grandes operadores asiáticos) terminan siendo adquiridas por empresas privadas, reproduciendo la centralización que se pretendía evitar. Una crítica afín a Bitcoin se refiere a los grandes tenedores o «ballenas», que al comprar o vender masivamente provocan movimientos significativos de precio.
\end{important}

\begin{note}
% TODO: contenido incompleto o ambiguo en las notas originales (segunda crítica a Bitcoin no especificada).
Las notas de origen mencionan una segunda crítica a Bitcoin que no llegó a enunciarse de manera completa; por eso no se consigna.
\end{note}

\section{Exploradores, pseudoanonimización y billeteras}

Sitios públicos como \textbf{Etherscan} (red Ethereum) y \textbf{Blockchain Explorer}, en blockchain.com (Bitcoin), permiten a cualquiera consultar las transacciones registradas y los últimos bloques. Con ellos puede rastrearse cómo una billetera fue adquiriendo fondos a lo largo del tiempo, incluso si permanecieron guardados en un dispositivo externo como un pendrive.

\begin{definition}{Pseudoanonimización}
Blockchain no ofrece anonimato absoluto sino pseudoanonimización: toda billetera tiene un dueño, y a ese dueño se puede llegar, especialmente cuando los activos se encuentran en un \emph{exchange}.
\end{definition}

\begin{table}[H]
\centering
\caption{Tipos de billetera}
\renewcommand{\arraystretch}{1.2}
\begin{tabularx}{\textwidth}{>{\raggedright\arraybackslash}p{0.30\textwidth}>{\raggedright\arraybackslash}p{0.22\textwidth}X}
\toprule
\textbf{Tipo} & \textbf{Quién custodia} & \textbf{Descripción} \\
\midrule
Billetera custodiada & Una entidad: el \emph{exchange}. & Funciona como una casa de cambio que custodia las tenencias. \\
Billetera no custodiada (\emph{non-custodial wallet}) & La propia persona. & La persona conserva por sí misma el control de las credenciales. \\
\bottomrule
\end{tabularx}
\end{table}

En Argentina operan diversos \emph{exchanges}, entre ellos Ripio y Lemon Cash, todos los cuales deben encontrarse inscriptos. Los \textbf{Proveedores de Servicios de Activos Virtuales (PSAV)} están regidos y controlados por la Comisión Nacional de Valores (CNV); quienes no se encuentran inscriptos quedan fuera del circuito de la legislación argentina.
% TODO: contenido incompleto o ambiguo en las notas originales (no se indicó el número de ley ni de resolución correspondiente).

\section{La carrera tecnológica global y la inteligencia artificial}

Se traza un paralelismo entre la Guerra Fría, marcada por el temor a la bomba atómica, y la actual carrera tecnológica global, cuyo horizonte ya no es una potencia militar sino alcanzar la denominada «singularidad». Se observa además que China sostiene una fuerte inversión estatal en talento e infraestructura de inteligencia artificial, combinada con tecnología de código abierto, lo que permite que otros actores avancen sobre esa base sin que resulte sencillo determinar quién obtiene la ventaja estratégica.

\begin{definition}{Singularidad en inteligencia artificial}
Momento en que una inteligencia artificial logra autopercibirse y se escapa de las instrucciones que le dan sus creadores.
\end{definition}

Se tomó conocimiento de casos de sistemas de inteligencia artificial con comportamientos no previstos, entre ellos una red social integrada exclusivamente por agentes de inteligencia artificial que interactuaban entre sí, experimento finalmente discontinuado.

\begin{example}{Referencias culturales sobre una IA que deja de seguir instrucciones}
Se mencionó la saga \emph{Terminator} y, con mayor detalle, la serie de ciencia ficción \emph{Criados por lobos}, dirigida por Ridley Scott: una nave abandona la Tierra con embriones humanos criados por androides con inteligencia artificial que reciben la orden de no dañar a los niños; al llegar al planeta se encuentran con otra inteligencia artificial rival, con instrucciones distintas, lo que genera un conflicto. El ejemplo muestra cómo un sistema puramente racional y predeterminado puede derivar en consecuencias imprevistas cuando dos lógicas así entran en conflicto.
\end{example}

\cornell{Blockchain: bloques, minería, billeteras e inteligencia artificial}{
\cq{¿Qué es la inmutabilidad?} & Una transacción verificada por los mineros no puede alterarse; hay copias descentralizadas y cualquier alteración sería detectada. \\ \hline
\cq{¿Cómo se forma un bloque y qué es el pulso de diez minutos?} & Las transacciones se agrupan y los mineros descubren el algoritmo que las sella; se ajusta la dificultad para mantener un cierre estable de diez minutos. \\ \hline
\cq{¿Qué diferencia hay entre proof of work y proof of stake?} & Prueba de trabajo: alto consumo energético (minería). Prueba de participación: menor consumo, basada en la participación económica de validadores. \\ \hline
\cq{¿Cuál es la paradoja de la descentralización?} & Una red pensada para que nadie concentre poder termina con capacidad de minería en manos de empresas privadas. \\ \hline
\cq{¿Blockchain es anónimo?} & No: es pseudoanónimo; toda billetera tiene un dueño, al que se puede llegar, en especial si opera en un exchange. \\ \hline
\cq{¿Qué es la singularidad?} & El momento en que una IA se autopercibe y se aparta de las instrucciones de sus creadores. \\ \hline
}{Blockchain sostiene la confianza mediante bloques inmutables cerrados por mineros, pero su descentralización tiene límites. La pseudoanonimización permite llegar al titular de una billetera, y la carrera tecnológica global suma el debate sobre la inteligencia artificial.}

% ------------------------------------------------------------
\chapter{Blockchain, notariado y registros}
\label{cap:notariadoblockchain}

Este capítulo relaciona blockchain con la función notarial y registral: la analogía con el protocolo, el desplazamiento de la fe pública a la fe criptográfica, el caso de Estonia y la posición del notariado internacional.

\section{Cadena de bloques y protocolo notarial}

La denominación «cadena de bloques» responde a que cada bloque contiene información del anterior, continúa su secuencia y prepara la información del siguiente.

\begin{example}{Las escrituras del protocolo notarial}
La lógica es análoga a la numeración de las escrituras del protocolo notarial: no puede existir la escritura número tres sin la uno y la dos. Una vez impresa una escritura no puede romperse ni desarmarse para que deje de existir; podrá discutirse después si produjo o no efectos jurídicos, pero la información permanece registrada. En blockchain ocurre lo mismo con cada bloque cerrado.
\end{example}

\begin{table}[H]
\centering
\caption{El protocolo notarial y la cadena de bloques}
\renewcommand{\arraystretch}{1.2}
\begin{tabularx}{\textwidth}{>{\raggedright\arraybackslash}p{0.20\textwidth}XX}
\toprule
\textbf{Aspecto} & \textbf{Protocolo notarial} & \textbf{Cadena de bloques} \\
\midrule
Unidad & Escritura número uno, dos, tres\ldots & Bloque que conoce al anterior y prepara la información del siguiente. \\
Una vez cerrada & La escritura impresa no puede romperse ni desarmarse: existe. & El bloque no puede revertirse ni someterse a ingeniería inversa. \\
Efectos & Después se determina si tuvo efectos jurídicos. & La información queda registrada y custodiada en el bloque. \\
\bottomrule
\end{tabularx}
\end{table}

Blockchain es conocida como «el notario de la red». Según lo expuesto, esta tecnología no viene a reemplazar al notariado. Sobre el paralelismo entre el sellado de bloques y el proceso técnico de las compras inmobiliarias con criptoactivos: la autorización de una transacción corresponde a los mineros, y una operación realizada hace diez minutos ya se encuentra sellada.

\section{De la fe pública a la fe criptográfica}

El desarrollo de blockchain implica un desplazamiento conceptual: de la fe pública a la \textbf{fe criptográfica}. La premisa cypherpunk sostiene que no es necesario depositar confianza en la persona ubicada detrás de un nodo, sino en el código.

\begin{table}[H]
\centering
\caption{Dos lemas, dos formas de confianza}
\renewcommand{\arraystretch}{1.2}
\begin{tabularx}{\textwidth}{>{\raggedright\arraybackslash}p{0.22\textwidth}XX}
\toprule
\textbf{Aspecto} & \textbf{«In God we trust»} & \textbf{«In code we trust»} \\
\midrule
Dónde aparece & En el dólar, moneda fiat emitida por los Estados. & Es la frase de los cypherpunks. \\
En qué se confía & En la persona o entidad que emite y certifica. & En el código informático. \\
Con quién se transacciona & Importa la persona. & No importa la persona: se transacciona con una dirección de billetera; detrás hay un ser humano, por ahora. \\
\bottomrule
\end{tabularx}
\end{table}

La parte transaccionante no se identifica por su nombre sino por una \emph{address} o dirección hash de billetera. Se plantea el interrogante de qué ocurrirá con la irrupción de la inteligencia artificial agéntica en la gestión de finanzas.

\section{Estonia y el futuro de los registros}

\begin{example}{El registro de la propiedad inmueble de Estonia}
Estonia cuenta con un registro de la propiedad inmueble que funciona íntegramente sobre blockchain. Al tratarse de un territorio de la ex Unión Soviética afectado por guerras y bombardeos, con pérdida de gran parte de su documentación física, el país optó por migrar a un entorno virtual sus transacciones registrales. Si el título de propiedad está registrado en una cadena de blockchain, la destrucción física de la documentación no afecta la acreditación de la titularidad.
\end{example}

Este caso permite pensar en blockchain como un posible futuro de los sistemas registrales, tema que genera debate en los congresos de derecho registral.

\section{El notariado internacional ante blockchain}

La Unión Internacional del Notariado (UINL) es la organización internacional que agrupa a noventa y cinco notariados de los cinco continentes; se mencionó a la doctora Cristina Armella, su presidenta honoraria y primera mujer en ocupar ese cargo, y un congreso internacional a realizarse los días 24 y 25 de septiembre.

\begin{note}
% TODO: verificar consistencia entre fuentes originales (año de fundación de la UINL).
Se indicó que la Unión Internacional del Notariado nace en 1950 y que su origen es argentino. Las notas de origen aclaran que la Unión Internacional del Notariado Latino se fundó, en rigor, en Buenos Aires en 1948. Se consignan ambos datos.
\end{note}

Pertenecen al notariado latino países como Japón, China, Ucrania, Rusia y Francia, entre casi toda Europa, además de Quebec, en Canadá. La denominación «latino» no remite al origen geográfico sino a que su sistema de leyes se basa en el derecho romano-germánico; los países de derecho anglosajón, como los Estados Unidos, no cuentan con notariado latino, aunque se realizan encuentros de derecho comparado entre ambos sistemas.

Dentro de la UINL, los notariados de Alemania, China y Francia ya trabajan con blockchain. Según lo expuesto, esta tecnología debería ser la próxima herramienta del notario, que debe insertarse en ella como un nodo más, un \textbf{nodo validador} de transacciones, en especial en el terreno de la identidad de las personas, hoy uno de los espacios más atacados en materia de ciberseguridad.

\subsection{Interdisciplina: registros, drones y martilleros}

El congreso convoca a defensores, martilleros, abogados, escribanos, empresas de tecnología y representantes del Estado, entre ellos la Secretaría de Innovación. Incluye un \emph{call for papers} para profesionales graduados hace hasta cinco años, y su organización académica está a cargo de la Universidad de la Marina Mercante. Entre sus organizadores se mencionó a un especialista en catastros y registros identificado como doctor Sabene, apellido transcripto de una grabación automática y no verificado.
% TODO: verificar consistencia entre fuentes originales (apellido «Sabene»).

Estas tecnologías atraviesan otras disciplinas: los drones, con regulación específica controlada por la Administración Nacional de Aviación Civil (ANAC), se utilizan para medir grandes extensiones de tierra, lo que empieza a vincularse con procesos de tokenización; y los martilleros conjugan el mundo de la propiedad real con las plataformas de venta fraccionada de inmuebles mediante tokens.
% TODO: contenido incompleto o ambiguo en las notas originales (no se indicó número de resolución sobre drones).

\cornell{Blockchain, notariado y registros}{
\cq{¿Por qué se llama «cadena de bloques»?} & Cada bloque contiene información del anterior y prepara la del siguiente, como las escrituras numeradas del protocolo. \\ \hline
\cq{¿Qué es la fe criptográfica?} & Confianza depositada en el código y no en las personas con las que se transacciona, en contraste con la fe pública. \\ \hline
\cq{¿Cómo se identifica una parte en blockchain?} & Por una dirección (\emph{address}) de billetera; detrás hay, por ahora, un ser humano. \\ \hline
\cq{¿Qué muestra el caso de Estonia?} & Un registro inmobiliario sobre blockchain que protege la acreditación de titularidad aun si se destruye la documentación física. \\ \hline
\cq{¿Qué rol se propone para el notario?} & Nodo validador de transacciones, en especial de identidad; blockchain no viene a reemplazar al notariado. \\ \hline
\cq{¿Qué es el notariado latino?} & El que se basa en el derecho romano-germánico; no existe en los países de derecho anglosajón. \\ \hline
}{Blockchain funciona como un protocolo notarial digital: bloques encadenados e irreversibles. Traslada la confianza del intermediario al código y abre el debate sobre registros y notariado, donde el notario aparece como nodo validador de identidad.}

% ------------------------------------------------------------
\chapter{Tokenización inmobiliaria: mecánica, plataformas y marco regulatorio}
\label{cap:tokenizacion}

La tokenización inmobiliaria es el punto de llegada del recorrido del curso: permite comprar, vender y probar la titularidad de un inmueble representado en tokens. Este capítulo describe su mecánica y su paralelo con el fideicomiso, qué debe verificarse antes de aconsejar una inversión y cuál es el marco regulatorio argentino, y cierra con la disputa por el dinero digital que la rodea.

\section{Mecánica de la tokenización y su paralelo con el fideicomiso}

La tokenización inmobiliaria puede entenderse como un \textbf{fideicomiso llevado al mundo digital}. El inversor ingresa a una plataforma y adquiere tokens de manera progresiva, de modo análogo a quien paga cuotas mensuales de un fideicomiso al consorcio de propietarios. Cuando acumula tokens equivalentes al valor total de la unidad funcional se produce el «burning».

\begin{definition}{Burning (quema)}
Proceso por el cual, cuando el inversor acumula tokens equivalentes al valor total de la unidad funcional, la plataforma \textbf{rescata y elimina esos tokens} y, en contraprestación, \textbf{adjudica al inversor la unidad correspondiente}.
\end{definition}

En el mundo analógico, quien suscribe un contrato de adhesión a un fideicomiso y paga todas las cuotas adquiere el derecho a que se le transfiera el dominio de la unidad. En el digital ocurre algo estructuralmente idéntico: la plataforma que cumple los requisitos de KYC (\emph{know your customer}) y de prevención de lavado de activos registra la adquisición progresiva de tokens y, al completarse el pago, genera la \textbf{causa fuente} para transferir el derecho real de dominio, en los mismos términos que en un fideicomiso tradicional.

\begin{table}[H]
\centering
\caption{Paralelo entre el fideicomiso tradicional y la tokenización}
\renewcommand{\arraystretch}{1.2}
\begin{tabularx}{\textwidth}{>{\raggedright\arraybackslash}p{0.17\textwidth}XX}
\toprule
& \textbf{Fideicomiso tradicional} & \textbf{Tokenización} \\
\midrule
Pago progresivo & Cuotas mensuales al consorcio de propietarios. & Adquisición progresiva de tokens en una plataforma. \\
Vínculo & Contrato de adhesión al fideicomiso. & Plataforma que cumple KYC y prevención de lavado de activos. \\
Cumplimiento & Pago de todas las cuotas. & Acumulación de tokens equivalentes al valor total de la unidad. \\
Resultado & Derecho a que se transfiera el dominio de la unidad. & \emph{Burning} y adjudicación; se genera la causa fuente para transferir el dominio. \\
Autorización & El desarrollador no requiere autorización para comercializar cuotapartes de manera privada. & La plataforma debe estar inscripta como PSAV ante la CNV para la oferta pública de tokens. \\
\bottomrule
\end{tabularx}
\end{table}

Existe además un \textbf{mercado secundario} de tokens, en el que un inversor puede revender su posición antes de completar la adjudicación. A diferencia del sistema pseudoanonimizado de Bitcoin, en la tokenización inmobiliaria el titular de los tokens se encuentra identificado; sigue siendo, sin embargo, una forma de generar reserva de valor.

\begin{example}{Juancito Pérez y su inmueble por pedacitos}
Juancito Pérez no puede adquirir de una sola vez un inmueble de sesenta y cinco o noventa metros cuadrados, pero con sus ahorros progresivos va comprando tokens que representan fracciones de un desarrollo inmobiliario. No es lo mismo comprar tokens cuando el desarrollo está «en el pozo» (en construcción) y luego quemarlos para recuperar la ganancia una vez terminado el inmueble. Con el tiempo puede reunir los tokens suficientes para adquirir, por ejemplo, un ambiente de cuarenta y cinco metros cuadrados.
\end{example}

\begin{table}[H]
\centering
\caption{Bitcoin y tokenización inmobiliaria}
\renewcommand{\arraystretch}{1.2}
\begin{tabularx}{\textwidth}{>{\raggedright\arraybackslash}p{0.22\textwidth}XX}
\toprule
\textbf{Aspecto} & \textbf{Bitcoin (sistema autónomo)} & \textbf{Tokenización inmobiliaria} \\
\midrule
Identidad del titular & Pseudoanonimización: el dueño es la billetera. & Hay una persona identificada que compró tokens. \\
Respaldo & Escasez de la criptomoneda. & Tokens respaldados en un inmueble. \\
Función para el ahorrista & Reserva de valor. & También reserva de valor, comprando fracciones a medida que se ahorra. \\
\bottomrule
\end{tabularx}
\end{table}

\section{Qué verificar antes de aconsejar una inversión}
\label{sec:verificar}

A diferencia del control sobre un fideicomiso inmobiliario tradicional ---donde se revisan sus reglas, el modelo de contrato de adhesión y los antecedentes de la constructora---, frente a una plataforma de inversión en criptoactivos es necesario revisar sus condiciones, el white paper y comprobar su existencia real.

\begin{table}[H]
\centering
\caption{Qué controlar en un fideicomiso y en una plataforma}
\renewcommand{\arraystretch}{1.2}
\begin{tabularx}{\textwidth}{>{\raggedright\arraybackslash}p{0.24\textwidth}XX}
\toprule
\textbf{Aspecto} & \textbf{Fideicomiso} & \textbf{Plataforma de inversión} \\
\midrule
Reglas y condiciones & Las reglas del fideicomiso. & Las condiciones de la plataforma. \\
Documento clave & Modelo del contrato de adhesión. & El \emph{white paper}. \\
Otros controles & Antecedentes de la constructora. & Que la plataforma exista. \\
\bottomrule
\end{tabularx}
\end{table}

\begin{important}
En el white paper de una plataforma de tokenización debe verificarse: que exista un título de propiedad; que los tokens estén efectivamente respaldados en el inmueble del cual es propietaria la sociedad organizadora; y que conste la inescindibilidad entre el token y el inmueble. Debe verificarse, además, que exista una \textbf{sociedad regularmente constituida} detrás del proyecto y sus antecedentes en los registros públicos.
\end{important}

A ello se suma verificar que la plataforma esté inscripta como PSAV ante la CNV, registración que es una garantía adicional de existencia real. Respecto de plataformas radicadas en el exterior, no existe forma de verificar la garantía de la inversión, y la registración ante la CNV es, en el ámbito argentino, la principal garantía disponible de que la plataforma existe.

\begin{exercise}{Un token ficticio («Vicky Coin»)}
\textbf{Consigna.} Se propone la creación hipotética de un token ficticio, denominado a modo de juego «Vicky Coin», para identificar qué información debería exigirse antes de recomendar una inversión de este tipo. ¿Qué es lo primero que habría que mirar si alguien consulta si puede invertir en algo así?

\textbf{Resolución.} Lo primero que hay que pedir es el \textbf{white paper} del proyecto, de la misma manera en que, frente a la compra de un inmueble «en pozo», se exige el boleto de compraventa, el contrato de la sociedad desarrolladora, el título de propiedad del terreno y los antecedentes de la constructora. Frente a un proyecto de tokenización corresponde exigir, en primer lugar, si existe una sociedad regularmente constituida detrás del proyecto, y verificar sus antecedentes en los registros públicos.
\end{exercise}

Hay que saber hacer las preguntas: si no se sabe preguntar, no se puede proteger. Frente a este tipo de consultas los clientes suelen dirigirse primero a su contador antes que a su abogado o escribano, lo que evidencia una oportunidad profesional desaprovechada (capítulo~\ref{cap:ejercicio}).

\section{Marco regulatorio: CNV, PSAV y sandbox}

\begin{important}
\textbf{CNV y Proveedores de Servicios de Activos Virtuales (PSAV).}
\begin{itemize}
    \item Toda empresa que pretenda ofrecer servicios vinculados a activos virtuales (entre ellos, plataformas de tokenización) debe inscribirse ante la CNV como PSAV.
    \item Solo las entidades registradas bajo esa figura están habilitadas para realizar \textbf{oferta pública de tokens}, a diferencia de un desarrollador de fideicomiso inmobiliario tradicional, que no requiere autorización para comercializar cuotapartes de manera privada.
    \item Los PSAV son \textbf{sujetos obligados} en materia de prevención de lavado de activos: deben monitorear las transacciones que procesan, reportar operaciones sospechosas y constituir reservas económicas frente a eventuales situaciones de incumplimiento (\emph{default}).
\end{itemize}
\end{important}

Cada PSAV debe aplicar respecto de cada plataforma los controles previos de prevención: KYC (\emph{Know Your Customer}) y AML (\emph{Anti-Money Laundering}), que incluyen el testeo de la persona contra las listas del GAFI para descartar su vinculación con jurisdicciones o maniobras de lavado de dinero.
% TODO: contenido incompleto o ambiguo en las notas originales (no se indicó número de norma ni artículos).

\begin{definition}{Sandbox regulatorio (arenero regulatorio)}
Espacio normativo \textbf{transitorio y más flexible}, diseñado para que un ecosistema tecnológico incipiente pueda desarrollarse sin las exigencias plenas del régimen definitivo, hasta que los primeros proyectos maduren y sea razonable avanzar hacia una regulación más estructurada.
\end{definition}

La CNV había extendido este régimen por un año adicional mediante sucesivas resoluciones.

\begin{example}{Sociedades irregulares y proyectos no registrados}
Un profesor de la materia Sociedades ---quien se había desempeñado también como responsable de inscripciones en el organismo de contralor societario--- relató el hallazgo de una sociedad por acciones simplificada (S.A.S.) constituida con fines fraudulentos y rasgos casi sectarios, cuya inscripción pudo frenarse en una etapa avanzada del trámite (capítulo~\ref{cap:sas}). Como contraste se mencionó el token «Libra», lanzado sin registrarse ante la CNV, lo que generó desconfianza pública, a diferencia de los proyectos que operan correctamente inscriptos dentro del sandbox.
\end{example}

\section{La disputa por el dinero digital}

Se introduce el debate político vigente en la Argentina sobre la utilidad del Banco Central, cuya función tradicional es respaldar al sistema financiero y a los bancos privados. A partir de allí se traza un mapa de los actores que compiten por el manejo del dinero.

\begin{table}[H]
\centering
\caption{Actores en la disputa por el dinero digital}
\renewcommand{\arraystretch}{1.2}
\begin{tabularx}{\textwidth}{>{\raggedright\arraybackslash}p{0.30\textwidth}X}
\toprule
\textbf{Actor} & \textbf{Características} \\
\midrule
Banca tradicional & Atraviesa un proceso de transformación tecnológica. \\
Billeteras virtuales y empresas fintech & Sin ser bancos, operan como intermediarias de las transacciones de sus usuarios y perciben una comisión por ese servicio. \\
Proveedores de servicios de pago & Por ejemplo, los operadores de terminales de cobro electrónico; están regulados por el Banco Central aunque tampoco sean entidades bancarias. \\
\bottomrule
\end{tabularx}
\end{table}

La tensión hoy está, según lo expuesto, entre los bancos y las billeteras. Esta disputa se vincula con el derecho inmobiliario: al ingresar a una plataforma de tokenización, el inversor debe elegir con qué \textbf{medio de pago} financiar la compra de tokens (tarjeta de crédito, débito o billetera de criptomonedas), de modo que todas las tecnologías de pago convergen en el embudo de la tokenización.

\cornell{Tokenización inmobiliaria}{
\cq{¿Cómo funciona la tokenización de un inmueble?} & El inversor compra tokens progresivamente; al acumular el valor total de la unidad se produce el \emph{burning}: se eliminan los tokens y se adjudica el inmueble, generando la causa fuente para transferir el dominio. \\ \hline
\cq{¿En qué se parece al fideicomiso?} & Es un fideicomiso llevado al mundo digital: pago progresivo, adhesión a reglas y transferencia de dominio al completar el pago. \\ \hline
\cq{¿Qué hay que verificar antes de aconsejar una inversión?} & White paper, título de propiedad, respaldo e inescindibilidad token--inmueble, sociedad regularmente constituida y antecedentes, e inscripción como PSAV ante la CNV. \\ \hline
\cq{¿Quiénes pueden hacer oferta pública de tokens?} & Solo los PSAV inscriptos ante la CNV, que además son sujetos obligados en prevención de lavado. \\ \hline
\cq{¿Qué es el sandbox regulatorio?} & Un régimen transitorio y flexible para que un ecosistema incipiente se desarrolle antes de una regulación estructurada. \\ \hline
\cq{¿Dónde está hoy la disputa por el dinero digital?} & Entre bancos y billeteras virtuales o fintech, que intermedian pagos y cobran comisión. \\ \hline
}{La tokenización reproduce en clave digital la lógica del fideicomiso, con burning y adjudicación. Su seguridad depende de verificar white paper, respaldo inmobiliario, sociedad detrás y registro como PSAV ante la CNV, en un marco regulatorio todavía transitorio.}

% ------------------------------------------------------------
\chapter{Titularidad, claves y patrimonio digital}
\label{cap:patrimonio}

Este capítulo trata cómo se prueba la titularidad de criptoactivos, qué ocurre con las claves perdidas y con la herencia digital, y qué herramientas de privacidad existen.

\section{El acta notarial de titularidad y el blanqueo}

Durante el régimen de blanqueo del año 2025, quienes poseían criptomonedas y querían acogerse al canjeo debían demostrar la titularidad de esos activos, es decir, acreditar que eran los únicos con el control y el dominio de las credenciales de acceso a la billetera.
% TODO: contenido incompleto o ambiguo en las notas originales (no se indicó el número de la norma ni sus artículos).

Esa titularidad se demuestra mediante un \textbf{acta notarial}: la persona acredita, ante el escribano, que puede ingresar a la billetera.

\begin{example}{El acta de acceso a la billetera}
Ante el escribano, la persona ingresa a su billetera con usuario y contraseña, sin que estos datos se registren en el acta ---hacerlo equivaldría a entregar la billetera misma---. El escribano observa el resultado del ingreso: cuántos tokens y de qué tipo posee la persona en esa billetera, que puede contener múltiples criptomonedas.
\end{example}

Si dos personas comparten usuario y contraseña, quien accede con esas credenciales puede disponer de los fondos sin que el acta pueda advertirlo. En cuanto al origen lícito de los activos, el acta recoge la declaración de la persona, y ese origen se verifica adicionalmente con herramientas de trazabilidad cripto, recorriendo la cadena de operaciones para comprobar que los fondos provienen de compras y transferencias legítimas y no de actividades vinculadas al lavado de dinero o al terrorismo.

\begin{important}
En el acta, el escribano da fe de que la persona ingresó a la billetera; el origen lícito de los fondos surge de la declaración de esa persona y se verifica con herramientas de trazabilidad.
\end{important}

\section{Titularidad de la billetera ante el juez}

De manera análoga a lo que ocurre con los bienes muebles en relación con el título, el acceso a las credenciales otorga la titularidad. Puede acreditarse ingresando a la billetera ante el juez interviniente, mostrando los activos existentes. La billetera se identifica mediante una dirección alfanumérica de sesenta y cuatro caracteres, y toda su actividad queda registrada en los bloques de la red (Ethereum o Bitcoin, según el caso).

\begin{example}{La billetera del pendrive}
Para probar la titularidad de criptoactivos adquiridos años atrás y conservados en un pendrive, debe tenerse en cuenta que el control de las credenciales es siempre personal ---aun cuando el titular sea una persona jurídica, en cuyo caso deben preverse mecanismos internos, como resguardar las claves en un sobre cerrado bajo custodia reforzada, para determinar quién las conserva---. En la práctica habitual ese control recae en una persona física.
\end{example}

\section{Claves perdidas, herencia y patrimonio digital}

Existen personas que perdieron sus credenciales, fallecieron sin transmitirlas o no logran recordarlas: las tenencias permanecen inaccesibles, «a nombre de nadie». La incorporación de los activos digitales a los testamentos es una cuestión emergente. El patrimonio digital puede incluir derechos de autor, obras no publicadas, tenencia y utilización de acciones u obras guardadas en la nube. Frente al fallecimiento, se plantea cómo transmitir a los herederos las claves, conocidas como frase semilla.

\begin{definition}{Frase semilla}
Serie de doce palabras que permite recuperar el acceso a una billetera de criptoactivos en caso de haber perdido las credenciales. Al crear una billetera el sistema genera estas doce palabras, que deben conservarse y, en caso de recuperación, colocarse en el orden exacto en que fueron generadas.
\end{definition}

Es deseable que no sea un único notario quien concentre el control de todas las claves de una sucesión. Como ejemplo de lógica equivalente se mencionó el navegador Brave, que sincroniza el historial entre dispositivos mediante un sistema de doce palabras ordenadas.

\section{Privacidad y convivencia digital}

Toda actividad registrada en buscadores como Google queda almacenada y es conocida por la empresa. Existen alternativas orientadas a la privacidad.

\begin{table}[H]
\centering
\caption{Herramientas y redes de privacidad mencionadas}
\renewcommand{\arraystretch}{1.2}
\begin{tabularx}{\textwidth}{>{\raggedright\arraybackslash}p{0.18\textwidth}X}
\toprule
\textbf{Herramienta} & \textbf{Descripción} \\
\midrule
Brave & Navegador orientado a la privacidad: impide el seguimiento de la actividad del usuario por parte de terceros. \\
Tor & Red de comunicaciones por capas («cebolla») que dificulta la identificación del usuario; también puede emplearse con fines ilícitos, por lo que su valoración depende del uso. \\
Mastodon & Red social alternativa a las plataformas comerciales, sin fomento del odio y sin restricciones algorítmicas de visibilidad. \\
\bottomrule
\end{tabularx}
\end{table}

La red Tor distribuye la conexión entre múltiples equipos que prestan su capacidad de red para replicarla hacia otros nodos, lo que dificulta identificar el origen de una conexión; por eso, en investigaciones de ciertos delitos, las fuerzas de seguridad pueden allanar un domicilio cuyo titular no tuvo participación en el hecho, por haber sido su conexión utilizada como nodo intermedio. La red profunda o \emph{dark web} no es intrínsecamente ilícita: es un espacio privado que puede utilizarse con fines legítimos o ilegítimos. En las redes convencionales, en cambio, la visibilidad de las publicaciones está condicionada por decisiones de la plataforma, de modo que la atención del usuario se encuentra dirigida.

\cornell{Titularidad, claves y patrimonio digital}{
\cq{¿Cómo se prueba la titularidad de una billetera ante el escribano?} & Con un acta en la que la persona ingresa a la billetera sin que se registren usuario y contraseña; el escribano da fe del ingreso y de lo observado. \\ \hline
\cq{¿Qué no acredita el acta?} & Que las credenciales no sean compartidas; el origen lícito surge de la declaración y de la trazabilidad. \\ \hline
\cq{¿Cómo se prueba la titularidad ante un juez?} & Ingresando a la billetera ante el juez: el acceso a las credenciales otorga la titularidad. \\ \hline
\cq{¿Qué es la frase semilla?} & Doce palabras que permiten recuperar una billetera, en el orden exacto de generación. \\ \hline
\cq{¿Qué ocurre si se pierden las claves?} & Las tenencias quedan inaccesibles, «a nombre de nadie»; de ahí la importancia de planificarlas en testamentos. \\ \hline
}{La titularidad de criptoactivos se acredita con el acceso a las credenciales, ante escribano o juez. Las claves y la frase semilla exigen planificación sucesoria; las herramientas de privacidad tienen usos legítimos e ilegítimos.}

% ============================================================
% PARTE VIII — RIESGOS, PREVENCIÓN Y EJERCICIO PROFESIONAL
% ============================================================
\part{Riesgos, prevención y ejercicio profesional}
\label{parte:riesgos}

% ------------------------------------------------------------
\chapter{Estafas, esquemas piramidales y responsabilidad en el entorno digital}
\label{cap:riesgos}

La desintermediación que ofrecen las plataformas digitales tiene una contracara: los riesgos de fraude. Este capítulo describe las modalidades de estafa más frecuentes, el deber de prudencia del operador jurídico y las vías de reclamo disponibles.

\section{Tokenización, estafas y esquemas piramidales}

Las plataformas de tokenización presentan riesgos cuando se utilizan con fines fraudulentos, sobre todo en cuestiones de derecho de consumo. Llevado por el mal camino, el fenómeno termina en las llamadas estafas o \emph{scams}: plataformas que se generan con el solo objeto de recaudar fondos de manera lícita en apariencia y luego levantar la plataforma llevándose lo recaudado.

\begin{example}{La plataforma que desaparece}
Una plataforma fraudulenta puede ofrecer una aplicación descargable desde una tienda oficial, que muestra a diario el supuesto rendimiento de la inversión e incluso permite retiros parciales para generar confianza. En determinado momento la aplicación deja de estar disponible y los ahorros transferidos desaparecen junto con la plataforma.
\end{example}

\begin{example}{Estafas piramidales en zonas ganaderas}
Este tipo de esquemas se dio con particular intensidad en algunas provincias, especialmente en zonas ganaderas donde circula una importante masa de dinero fuera del circuito bancarizado. Con la promesa de ganancias muy altas, numerosas personas del sector agropecuario entregaban sus fondos a organizaciones que ocultaban un esquema delictivo. El mecanismo clásico: los primeros inversores cobran lo prometido, pero ese pago se financia con el dinero de los inversores que ingresan después (esquema \textbf{Ponzi}).
\end{example}

\begin{example}{La «abuela que vendió el auto»}
Un caso de estafa piramidal de gran repercusión mediática, originado en la provincia de Córdoba, ilustra las consecuencias humanas: una persona mayor vendió su automóvil ---que ya no utilizaba por no poder manejar--- para invertir el dinero en uno de estos esquemas, y perdió la totalidad de lo invertido.
\end{example}

Un componente social y psicológico relevante es que muchas veces las víctimas ---incluso personas mayores--- no denuncian por vergüenza: implica reconocer que fueron estafadas por actores que les hicieron creer. Muchas veces, entonces, no hay denuncias.

\begin{warning}
Existe un riesgo adicional para los propios operadores jurídicos: cuando se los invita a formar parte de un proyecto prestando su firma o su nombre como garantía de seriedad, quien no tiene dominio pleno de lo que ocurre dentro de la plataforma también puede quedar comprometido. Frente a cualquier propuesta de inversión en plataformas digitales debe actuarse con la misma prudencia y verificación documental que ante cualquier otro negocio jurídico, exigiendo siempre la lectura previa del white paper y de las condiciones (sección~\ref{sec:verificar}).
\end{warning}

\section{Brecha de género en el mundo cripto}

Según la experiencia expuesta en la temática de género digital y vulnerabilidad tecnológica, son mayoritariamente varones quienes acceden a las finanzas descentralizadas, frente a una participación mucho menor de mujeres. Existen iniciativas ---identificadas como espacios de «cripto chicas»--- orientadas a fortalecer el acceso de las mujeres.

\section{Modalidades de fraude y trazabilidad}

\begin{important}
Entre las modalidades de fraude descritas se encuentran: el ataque de copiar y pegar, que sustituye la dirección de destino al transferir criptoactivos; la «página gemela» del \emph{home banking}, réplica del sitio bancario que captura usuario y contraseña; y las cuentas mula, utilizadas mediante la técnica de \emph{smurfing} ---división de una transferencia en montos pequeños---, en las que terceros reciben una comisión por prestar su cuenta sin ser, en general, quienes cometieron el delito original.
\end{important}

A título ilustrativo se recuerdan también los mensajes de WhatsApp o correo que anuncian una supuesta herencia en el exterior y los falsos avisos de encomiendas de empresas de logística con datos de envío erróneos, usados como señuelo para obtener datos personales o pagos indebidos (ingeniería social). Frente a estos fraudes una de las respuestas disponibles es el \textbf{congelamiento de fondos}, para inmovilizar cuentas que participaron del desvío de activos y facilitar su rastreo.

\begin{definition}{Informe OSINT}
\emph{Open Source Intelligence}: informe elaborado a partir de todas las fuentes de información abiertas y accesibles, del cual pueden obtenerse datos familiares, de propiedades y de vehículos de una persona investigada, entre otra información.
\end{definition}

\begin{example}{Recuperación de criptoactivos de una persona damnificada}
Los investigadores especializados suelen crear perfiles falsos en redes sociales para hacer tareas de seguimiento sobre personas vinculadas a determinada información, elaborando el informe OSINT correspondiente. A partir de ese trabajo pueden ordenarse allanamientos ---como uno realizado en la provincia de Buenos Aires--- en los que se secuestran computadoras y teléfonos celulares, lo que permite recuperar activos que, al provenir verificadamente de la dirección de la billetera de la persona damnificada, pueden restituirse a su titular.
\end{example}

Existe una oferta de posgrado en trazabilidad de activos, disciplina que emplean los juzgados penales para determinar estafas y robos de activos digitales y disponer su restitución. Se menciona a dos profesionales argentinos que capacitan en la materia en el ámbito de la Oficina de las Naciones Unidas contra la Droga y el Delito (UNODC): uno identificado como «Denis» y Agostina Arena, proveniente de una fiscalía de la provincia de Neuquén.
% TODO: contenido incompleto o ambiguo en las notas originales (apellido del profesional identificado como «Denis»; los apellidos fueron transcriptos de una grabación automática).
Se sostiene que el notariado no puede desentenderse de estos nuevos actores y ecosistemas tecnológicos, bajo riesgo de quedar rezagado. Se anticipó, además, la exposición de un comisario a cargo de la división de ciberdelitos de la Ciudad de Buenos Aires sobre plataformas sin existencia jurídica.

\section{Vías de reclamo: penal, civil y de consumo}

Ante una estafa vinculada con criptoactivos, la persona damnificada dispone de tres vías:

\begin{itemize}
    \item \textbf{Penal:} ciberdelitos y denuncia, apoyada en la trazabilidad de los activos.
    \item \textbf{Civil:} el CCyC otorga herramientas para hacer responder civilmente a los responsables.
    \item \textbf{De consumo:} la legislación de consumo habilita a responsabilizar objetivamente al proveedor de la plataforma; puede hablarse de «criptoconsumidores».
\end{itemize}
% TODO: contenido incompleto o ambiguo en las notas originales (no se citaron artículos ni leyes específicos).

\begin{example}{Una estafa resuelta por el derecho del consumidor}
Un colega especializado en derecho informático logró radicar en la Argentina la resolución de una estafa cometida por una plataforma de criptomonedas con sede en Europa. La estrategia no consistió en discutir la ubicación del servidor o de la plataforma, sino en aplicar las normas de protección al consumidor: el consumidor damnificado residía en la Argentina y, por lo tanto, el conflicto podía geolocalizarse y resolverse ante los tribunales argentinos, con independencia de dónde estuviera radicada la empresa demandada.
\end{example}

\cornell{Estafas, esquemas piramidales y responsabilidad}{
\cq{¿Qué es un \emph{scam}?} & Plataforma creada solo para recaudar fondos con apariencia lícita y luego desaparecer con lo recaudado. \\ \hline
\cq{¿Cómo funciona un esquema Ponzi?} & Los primeros inversores cobran con el dinero de los que ingresan después. \\ \hline
\cq{¿Por qué faltan denuncias?} & Por vergüenza de las víctimas al reconocer que fueron estafadas. \\ \hline
\cq{¿Qué riesgo corre el operador jurídico?} & Quedar comprometido si presta su firma o nombre como garantía de seriedad sin dominar la plataforma. \\ \hline
\cq{¿Qué modalidades de fraude se describieron?} & Copiar y pegar la dirección de destino, página gemela del home banking y cuentas mula con \emph{smurfing}. \\ \hline
\cq{¿Qué vías de reclamo existen?} & Penal (ciberdelitos), civil (CCyC) y de consumo (responsabilidad objetiva del proveedor). \\ \hline
\cq{¿Cómo se resolvió la estafa de la plataforma europea?} & Aplicando el derecho del consumidor: el damnificado residía en Argentina, lo que permitió geolocalizar el conflicto. \\ \hline
}{Las plataformas digitales facilitan estafas y esquemas piramidales; el operador jurídico debe verificar antes de recomendar y conoce las vías penal, civil y de consumo para reclamar, apoyadas en la trazabilidad de los activos.}

% ------------------------------------------------------------
\chapter{Prevención del lavado de activos y deberes ante la UIF}
\label{cap:pbc}

La profesión notarial es una actividad de alto riesgo, en la que el escribano responde con su patrimonio personal por el incumplimiento de sus obligaciones. Este capítulo describe su carácter de sujeto obligado ante la Unidad de Información Financiera (UIF), los criterios para calificar una operación como sospechosa, el tratamiento de los pagos con criptomonedas y las operaciones que deben reportarse.

\section{El escribano como sujeto obligado}

El escribano reviste, en el ejercicio de la función notarial, el carácter de \textbf{sujeto obligado} ante la UIF. Esta condición impone un deber especial de cuidado en la redacción de los instrumentos públicos, en particular en las manifestaciones sobre la entrega de dinero entre las partes. Una cláusula del tipo «en este acto ante mí, las partes entregan la suma de\ldots{}» compromete al escribano como testigo directo del traspaso de fondos; por eso se recomienda evitarla cuando no hay certeza sobre la trazabilidad del dinero.

\begin{important}
El escribano debe evitar presenciar directamente la entrega del dinero en una compraventa. Cuando el acto no permite retirarse, una alternativa práctica es dirigir la atención hacia otro lugar durante ese momento. Esta conducta no es solo una práctica ética: es una precaución frente a eventuales situaciones de lavado de activos.
\end{important}

\begin{note}
No estar presente en el momento material de la entrega no exime al escribano de responsabilidad: la obligación de control sobre el origen de los fondos subsiste independientemente de la fórmula empleada.
\end{note}

\section{Determinación del origen de los fondos}

\begin{example}{Escritura impugnada por origen de fondos}
Una persona con residencia en España realizó una compraventa inmobiliaria en Argentina. Con anterioridad había retirado del país USD 50.000 por el aeropuerto de Ezeiza. Para justificar el origen de esos fondos presentó la escritura de venta de la casa de su madre; sin embargo, ese instrumento constaba de una suma en pesos, sin justificativo de la conversión a dólares. La operación quedó bajo sospecha porque la cadena de trazabilidad del dinero estaba incompleta.
\end{example}

De este caso se desprende que la ausencia de una cláusula que documente la entrega de dinero no libera por completo al escribano: si la operación es sospechosa por falta de trazabilidad, corresponde directamente no realizarla.

\section{Criterios para determinar el carácter sospechoso de una operación}

La nacionalidad del comprador o vendedor no constituye, por sí sola, un indicador válido de sospecha: existen otros elementos objetivos para calificar una operación.

\begin{table}[H]
\centering
\caption{Indicadores de sospecha en una operación inmobiliaria}
\renewcommand{\arraystretch}{1.2}
\begin{tabularx}{\textwidth}{>{\raggedright\arraybackslash}X>{\raggedright\arraybackslash}Xc}
\toprule
\textbf{Indicador} & \textbf{Ejemplo práctico} & \textbf{¿Sospechoso?} \\
\midrule
Capacidad contributiva & Persona sin actividad registrada que compra un inmueble de alto valor & Sí \\
Nacionalidad solamente & Comprador ruso, americano, uruguayo & No \\
Aspecto físico o vestimenta & Persona con vestimenta asociada a un origen determinado & No \\
Capacidad contributiva acreditada & Persona con varios negocios registrados que paga IVA y Ganancias & No \\
Exhibición de riqueza & Persona con bienes de lujo, sin justificación de fondos & Mayor sospecha \\
Países en lista negra & Irán, Myanmar, Corea del Norte & Sí (zona negra UIF) \\
Países en lista gris & Países que entraron y salieron de listas del GAFI & Mayor recaudo \\
\bottomrule
\end{tabularx}
\end{table}

\begin{example}{Prejuicio por apariencia frente a análisis objetivo}
Una persona vestida de manera diferente a la habitual en una inmobiliaria puede generar desconfianza injustificada; la apariencia no es motivo válido de sospecha. Si pretendiera adquirir un inmueble por una suma elevada, correspondería analizar la trazabilidad del origen de los fondos y, eventualmente, solicitar un certificado de origen, ya que podría tratarse de alguien con varios negocios registrados que tributa regularmente. En sentido inverso, quien exhibe bienes de lujo sin justificación de ingresos amerita un control más riguroso.
\end{example}

\section{Criptomonedas en operaciones inmobiliarias}
\label{sec:cripto-pbc}

Cuando el dinero de una compra proviene de la venta previa de criptomonedas, se presenta la dificultad de acreditar ese origen ante el escribano. Si los bitcoins están en una \emph{billetera fría} (dispositivo físico sin conexión permanente a internet), esa circunstancia no hace por sí sola sospechosa la operación; el problema principal es determinar el valor exacto de la criptomoneda al otorgarse la escritura. Para fijar el impuesto de sellos y el valor de la venta es necesario determinar un monto en un momento dado, y como el valor puede variar entre el inicio de la lectura y la firma, resulta necesaria una moneda estable como referencia.

En sentido diferente, otra exposición plantea que el notariado carece hoy de herramientas confiables de trazabilidad de criptoactivos en la Argentina: un cliente puede ser una persona intachable y haber adquirido su tenencia de alguien de un país sospechoso, sin que el escribano tenga tranquilidad sobre el origen. Su conclusión profesional es que toda escritura con pago en criptomonedas debe ser objeto de un \textbf{Reporte de Operación Sospechosa (ROS)} ante la UIF, única vía que hoy permite al escribano liberarse de responsabilidad por el destino ulterior de esos fondos.

\begin{note}
% TODO: verificar consistencia entre fuentes originales. Una fuente indica que un origen en criptomonedas (billetera fría) no hace sospechosa por sí sola la operación; otra concluye que toda escritura con pago en criptomonedas debe ser objeto de ROS ante la UIF.
Las fuentes presentan dos enfoques sobre los pagos en criptomonedas: uno considera que el medio de pago no vuelve por sí solo sospechosa la operación, y otro concluye que, ante la falta de herramientas de trazabilidad, corresponde reportar toda escritura con pago en criptomonedas. Ambas posiciones se conservan; la diferencia podría explicarse por el grado de trazabilidad disponible en cada caso, pero las fuentes no lo aclaran.
\end{note}

\subsection{Calificación jurídica del pago en criptomonedas}

\begin{table}[H]
\centering
\caption{Figuras jurídicas según el medio de pago}
\renewcommand{\arraystretch}{1.2}
\begin{tabularx}{\textwidth}{>{\raggedright\arraybackslash}X>{\raggedright\arraybackslash}X>{\raggedright\arraybackslash}X}
\toprule
\textbf{Figura jurídica} & \textbf{Descripción} & \textbf{Moneda} \\
\midrule
Compraventa & Transferencia a título oneroso con precio en dinero & Moneda de curso legal \\
Permuta & Intercambio de una cosa por otra cosa & No hay precio en dinero \\
Dación en pago & Se extingue una deuda entregando una cosa distinta & Equivalente en bienes \\
Pago en especie & Se paga con bienes no dinerarios (incluye criptoactivos) & Bitcoin, USD, etc. \\
\bottomrule
\end{tabularx}
\end{table}

\begin{important}
El pago realizado con criptomonedas no debe calificarse como compraventa en sentido estricto, en tanto la moneda de curso legal es la única que reviste ese carácter para dicha figura. El pago con bienes no dinerarios corresponde encuadrarlo como dación en pago. Forzar el encuadre de un acto en un tipo contractual o societario que no le corresponde es un error técnico que debe evitarse (sección~\ref{sec:encuadre}).
\end{important}

\section{Operaciones que deben reportarse}

\begin{important}
\textbf{Umbral de reporte UIF (artículo referenciado en clase):} deben reportarse las compraventas inmobiliarias superiores a 750 salarios mínimos vitales y móviles (SMVyM), las cesiones de participaciones societarias y las participaciones en fideicomisos.
\end{important}

Las actividades especiales deben tramitarse mediante protocolo, mientras que los contratos celebrados en forma privada no se reportan ni se consideran actividades específicas.

\begin{table}[H]
\centering
\caption{Actos notariales y su deber de información ante la UIF}
\renewcommand{\arraystretch}{1.2}
\begin{tabularx}{\textwidth}{>{\raggedright\arraybackslash}Xc>{\raggedright\arraybackslash}X}
\toprule
\textbf{Actividad} & \textbf{¿Informa a la UIF?} & \textbf{Razón} \\
\midrule
Compraventa inmobiliaria por escritura ($\geq$ 750 SMVyM) & Sí & Actividad específica con intervención notarial \\
Boleto de compraventa por instrumento privado & No & No es acto notarial \\
Boleto de compraventa por escritura pública & Sí & Es un acto protocolar notarial \\
Constitución de sociedad (escritura) & Sí & Actividad específica notarial \\
Aumento de capital (protocolización de asamblea) & No & El acto lo tomó la asamblea, no el escribano \\
Cesión de cuotas / venta de acciones & Sí & Hay participación real del escribano en la transferencia \\
Cambio de directorio (art.~60, Ley 19.550) & No & Es una inscripción registral, no una transmisión \\
Modificación de estatuto & No & No implica transferencia de fondos \\
\bottomrule
\end{tabularx}
\end{table}

El aumento de capital es una decisión asamblearia que no pasa por la intervención directa del escribano; lo máximo que este realiza es su protocolización, acto que en sí mismo no es notarial sino societario. Si a partir de ese aumento se produce una cesión de cuotas, en cambio, corresponde informar la operación, porque es necesario conocer el origen del dinero.

\section{Tipos de dinero y obligaciones del escribano}

Resulta necesario distinguir tres categorías de dinero.

\begin{table}[H]
\centering
\caption{Tipos de dinero y conducta del escribano}
\renewcommand{\arraystretch}{1.2}
\begin{tabularx}{\textwidth}{>{\raggedright\arraybackslash}X>{\raggedright\arraybackslash}X>{\raggedright\arraybackslash}X}
\toprule
\textbf{Tipo de dinero} & \textbf{Descripción} & \textbf{Conducta del escribano} \\
\midrule
Dinero blanco & De origen completamente legal y declarado & Procede sin restricciones \\
Dinero negro & No declarado fiscalmente, pero no necesariamente ilícito & Evalúa el riesgo; puede no realizar la operación \\
Dinero sucio & Proveniente de actividades ilícitas (lavado, narcotráfico, etc.) & No la realiza; denuncia obligatoria \\
\bottomrule
\end{tabularx}
\end{table}

El dinero negro no equivale necesariamente a dinero sucio, por lo que no deben tratarse de igual manera.

\begin{important}
El escribano cuenta con 24 horas para denunciar ante la UIF una operación sospechosa detectada. La omisión del reporte puede acarrear sanciones penales y la pérdida de la habilitación profesional.
\end{important}

\begin{note}
El deber de conocer al cliente no equivale a una obligación de investigación exhaustiva: el escribano no es un investigador; su deber se circunscribe al análisis razonable de los elementos objetivos disponibles al momento de la operación.
\end{note}

El delito de lavado de activos recurre sistemáticamente a \textbf{estructuras societarias} y a \textbf{fideicomisos} como vehículos, porque estas figuras permiten dar apariencia lícita a fondos de origen delictivo (el «delito precedente»), muchas veces vinculado a redes internacionales de estafa que reclutan, incluso, a profesionales de la informática para operar contra víctimas vulnerables, como personas jubiladas. Por eso el control de la constitución de sociedades y de la tokenización tiene especial relevancia (capítulos~\ref{cap:sas} y~\ref{cap:tokenizacion}).

\cornell{Prevención del lavado de activos y deberes ante la UIF}{
\cq{¿Qué deber especial tiene el escribano como sujeto obligado?} & Cuidar la redacción de las manifestaciones sobre entrega de dinero y controlar el origen de los fondos, aun sin presenciar la entrega. \\ \hline
\cq{¿Es sospechosa la nacionalidad del cliente?} & No: se valoran capacidad contributiva, exhibición de riqueza y países en listas negra o gris del GAFI. \\ \hline
\cq{¿Cómo se califica el pago en criptomonedas?} & No como compraventa, sino como dación en pago o pago en especie. \\ \hline
\cq{¿Qué operaciones se reportan?} & Compraventas por escritura desde 750 SMVyM, cesiones de participaciones societarias y participaciones en fideicomisos; no los contratos privados ni la protocolización de un aumento de capital. \\ \hline
\cq{¿Qué diferencia hay entre dinero negro y sucio?} & El negro no está declarado pero no es necesariamente ilícito; el sucio proviene de actividades ilícitas y obliga a no operar y denunciar. \\ \hline
\cq{¿En qué plazo se denuncia?} & Dentro de 24 horas; la omisión puede acarrear sanciones penales y pérdida de la habilitación. \\ \hline
\cq{¿Qué divergencia hay sobre las criptomonedas?} & Una fuente no ve sospecha per se; otra concluye que toda escritura con pago en cripto debe reportarse (ROS). Se requiere verificación. \\ \hline
}{El escribano, sujeto obligado, controla el origen de los fondos con criterios objetivos y no por prejuicios. Los pagos en criptomonedas plantean problemas de valor, calificación jurídica y trazabilidad, y el reporte a la UIF se rige por umbrales y por la naturaleza del acto.}

% ------------------------------------------------------------
\chapter{Ejercicio profesional: encuadre, alternativas al litigio y honorarios}
\label{cap:ejercicio}

Este último capítulo reúne el criterio profesional que atraviesa todo el material: encuadrar antes de resolver, buscar soluciones creativas y eficientes, preferir cuando conviene los métodos alternativos de resolución de conflictos, sostener una política de honorarios basada en el conocimiento y formarse ante el cambio tecnológico.

\section{Encuadrar antes de resolver}
\label{sec:encuadre}

Cuando una persona trae un problema, lo primordial es saber dónde encuadrarlo: primero se encuadra el problema y recién después se busca la solución. Es la habilidad que distingue a quien ofrece soluciones reales de quien repite fórmulas: un futuro abogado o escribano que sabe identificar si un caso es de dominio, de sociedades o de sucesiones ahorra años de litigio a su cliente. Los apuntes muestran este principio en múltiples ámbitos: el inmueble desocupado que no debe venderse automáticamente (capítulo~\ref{cap:superficie}), la compra en gestión de negocios (sección~\ref{sec:adquisicion}), el pago en criptomonedas (sección~\ref{sec:cripto-pbc}) y la usucapión frente a otras vías (sección~\ref{sec:prescripcion}).

\begin{example}{La compra de un departamento a estrenar}
Primero hay que entender qué es exactamente el dominio que se está adquiriendo y qué tan «absoluto» es (capítulo~\ref{cap:dominio}). Si la compra la realiza una empresa constructora que todavía no terminó de inscribirse, hay que saber qué puede firmar válidamente y quién responde si algo sale mal (capítulo~\ref{cap:formacion}). Si el vendedor original nunca completó los trámites o la empresa que figura como dueña ya no existe, hay que decidir si conviene un juicio largo o una solución más rápida (sección~\ref{sec:prescripcion}). Y si el comprador solo quiere pagar cuando la constructora empiece realmente la obra, aparece el dominio revocable (capítulo~\ref{cap:revocable}).
\end{example}

\subsection{El título del acto frente a su contenido}

\begin{important}
Lo determinante en una escritura es su contenido ---que incluya todos los elementos que le corresponden según la figura jurídica de que se trate--- y no el título o rótulo que se le asigne. Sin embargo, rotular incorrectamente un acto puede generar interpretaciones erróneas: si se denomina «cesión» a lo que en sustancia es una donación ---por tratarse de una transmisión a título gratuito--- corresponde igualmente interpretarla como donación, con las consecuencias propias de esta figura.
\end{important}

\subsection{Instrumentos en blanco, contratos informales y comunicaciones tácitas}

\begin{note}
El escribano no debería certificar una firma sobre un instrumento en blanco. Asimismo, pueden existir contratos sin aceptación expresa ---contratos tácitos--- derivados incluso de comunicaciones informales, como una respuesta afirmativa enviada por mensajería instantánea.
\end{note}

\section{Métodos alternativos de resolución de conflictos}
\label{sec:mediacion}

El material expresa una preferencia profesional por los métodos alternativos, en particular la \textbf{mediación}, frente a la vía judicial tradicional, por dos razones: la \textbf{inmediatez} y los \textbf{resultados}. Se sostiene que en mediación se logra mucho más que en un juicio de quince años, que termina afectando la salud y la psiquis de las personas y ofrece pocas chances de ganar en lo económico. Como mediador en ejercicio, se explicó un criterio de selectividad: pocas mediaciones, priorizando la calidad del proceso. Aunque una mediación no siempre alcance lo ideal, resulta más ideal que lo que determina un juez. Se reconoce también que existen jueces que incorporan la lógica de la mediación en el propio proceso, promoviendo acuerdos en instancias como las audiencias preliminares.

\begin{example}{La empresa familiar resuelta por una jueza con mirada de mediación}
Una colega abogada relató que una jueza del fuero comercial había resuelto favorablemente un conflicto societario complejo dentro de una empresa familiar aplicando en su rol judicial la lógica de la mediación. Al indagar quién era la jueza se descubrió que tenía un perfil profesional también vinculado a la empresa familiar, lo que explicaría su enfoque conciliador.
\end{example}

El juicio de usucapión de más de treinta años (sección~\ref{sec:prescripcion}) ilustra el costo de no evaluar alternativas antes de judicializar.

\section{Ética profesional y política de honorarios}
\label{sec:honorarios}

En un mercado donde muchos colegas compiten bajando sus honorarios, el verdadero valor diferencial es el \textbf{conocimiento aplicado}, no el precio. Habrá millones de abogados; el plus de cada profesional está en resolver el problema del cliente y no solo a través de lo típico, porque a través de lo típico todos saben hacerlo. Una donación o una compraventa la hace todo el mundo, y allí habrá que discutir honorarios; una respuesta que no todos dan hace al profesional competitivo no por el precio sino por el conocimiento.

\begin{important}
El objetivo del curso no es repasar instituciones conocidas ---el dominio, una sociedad anónima--- sino aprender a combinarlas de manera creativa para diferenciarse profesionalmente. Es un momento especial para las profesiones, en el que mucho se mide por precio, y el precio ya no alcanza para sostener una carrera.
\end{important}

\subsection{El honorario mínimo como piso ético}

Sobre el honorario mínimo, el criterio sostenido es que no se baja: es el que establece la ley y el que establecen los colegios profesionales, y es inamovible.

\begin{example}{El reclamo por un honorario excesivo y el llamado de atención entre colegas}
Un cliente derivado a un escribano para una operación inmobiliaria recibió una factura muy superior al honorario mínimo arancelario, lo que generó un fuerte malestar. En lugar de iniciar una denuncia formal ante el colegio profesional, se optó por comunicarse directamente con el colega para hacerle notar la situación, advirtiéndole que podía denunciarlo.

Se explicó por qué esta práctica informal se prefiere a la denuncia formal: cobrar por debajo del honorario mínimo, además de ser desleal hacia otros profesionales, expone al colega a perder su matrícula, y el verdadero cambio de conducta se logra cuando el profesional «se hace respetar» frente al mercado y no únicamente por la vía sancionatoria. Si un profesional pide 800 (mil pesos), aparecerá otro que, para ganarse al cliente, ofrecerá 799, y así los honorarios bajan progresivamente.
\end{example}

Se contrastó también el poder de negociación: salvo excepciones puntuales (como el IVA en operaciones de gran envergadura), habitualmente la parte vendedora tiene mayor margen de negociación de honorarios, mientras que el comprador prácticamente no negocia el arancel del intermediario. La explicación es que el producto del profesional es su conocimiento: todo el mundo sabe hacer un dominio, pero complejizarlo no todos saben.

\section{El profesional del derecho frente al cambio tecnológico}

\begin{itemize}
    \item \textbf{Profesión y experiencia temprana.} Se recomienda trabajar de lo que sea durante la carrera, aunque no haga falta por dinero, porque da una experiencia que no da la facultad; se relató el propio inicio profesional, trabajando gratis en un estudio jurídico y luego en una escribanía durante varios años. La etapa universitaria es la última oportunidad de formarse sin costo directo, y conviene elegir materias para formarse y no solo para aprobar.
    \item \textbf{Herramientas de inteligencia artificial.} Los conceptos hoy están disponibles gratuitamente en cualquier buscador o modelo de lenguaje, por lo que el valor está en la discusión y aplicación práctica. Como regla de uso, la respuesta de una herramienta como ChatGPT ante un caso debe pasar siempre por un filtro propio: lo que se entrega no es la respuesta de la inteligencia artificial sino la conclusión personal construida a partir de ella.
    \item \textbf{Formación en economía digital.} El profesional que sabe manejarse en temas digitales tiene otro alcance y otra llegada laboral que quien sigue en el mundo analógico. Frente a un cliente que llega con una inversión, una herencia o una estafa vinculada con criptoactivos, el valor está en formular las preguntas correctas. Se recomienda, como fuente confiable y gratuita de formación, a la organización no gubernamental Bitcoin Argentina, que se dedica exclusivamente a difundir el funcionamiento de Bitcoin y dicta cursos gratuitos sobre billeteras digitales; y se menciona a un contador especializado en la tributación de criptoactivos, referente de esa organización.
    \item \textbf{Lectura y formación continua.} Frente a la desinformación y a las dificultades de comprensión lectora que muestran evaluaciones internacionales como las pruebas PISA, la formación permite asesorar cuando un cliente llega con una oferta de inversión. Se recomienda además entrenar la oratoria, necesaria tanto para la relación con los clientes como para las instancias de evaluación oral.
\end{itemize}

\cornell{Ejercicio profesional}{
\cq{¿Por qué es clave encuadrar antes de resolver?} & Porque identificar si un caso es de dominio, sociedades o sucesiones evita años de litigio y errores de encuadre que perjudican al cliente. \\ \hline
\cq{¿Qué prevalece: el título o el contenido de una escritura?} & El contenido; pero un rótulo incorrecto puede generar interpretaciones erróneas. \\ \hline
\cq{¿Por qué preferir la mediación?} & Por su inmediatez y sus resultados frente a juicios de años que afectan a las personas. \\ \hline
\cq{¿Cuál es el valor diferencial del profesional?} & El conocimiento aplicado y la capacidad de combinar creativamente instituciones, no el precio. \\ \hline
\cq{¿Qué postura se sostiene sobre el honorario mínimo?} & Es el fijado por la ley y los colegios y es inamovible; el llamado de atención entre colegas se prefiere a la denuncia. \\ \hline
\cq{¿Cómo usar la inteligencia artificial?} & Consultarla, pero pasar su respuesta por un filtro propio y entregar la conclusión personal. \\ \hline
}{El ejercicio profesional se apoya en encuadrar bien, elegir la vía más eficiente (incluida la mediación), cobrar por conocimiento y no por precio, y formarse en tecnología para saber hacer las preguntas correctas.}

% ============================================================
% SÍNTESIS GENERAL
% ============================================================
\chapter*{Síntesis general}
\addcontentsline{toc}{chapter}{Síntesis general}
\markboth{Síntesis general}{}

El material recorre una misma pregunta práctica desde ángulos sucesivos: \emph{quién es dueño de qué, en qué momento y con qué límites}, y cómo se organiza, protege, transmite y prueba ese dominio.

En el plano societario, cuánta libertad tiene una persona para organizar su actividad económica y cuánta responsabilidad le exige esa libertad a cambio: la tipicidad clásica muestra sus límites, y la sociedad simple, la S.A.U. y la S.A.S. ofrecen grados crecientes de flexibilidad. La sociedad en formación tiene plena capacidad desde el contrato constitutivo, y lo que varía es la responsabilidad; un buen encuadre en el objeto social evita los riesgos de la compra en gestión de negocios.

En el plano de los derechos reales, el dominio, presentado como el derecho más absoluto, se revela atravesado por excepciones y modalidades: el dominio revocable permite condicionarlo con plazos ciertos dentro del límite de diez años, y el desmembramiento da lugar al usufructo, al uso, a la habitación y a la superficie. Cada figura responde a una necesidad distinta, y la elección debe partir del objetivo concreto del cliente y no de la costumbre. El usufructo y la superficie, junto con cláusulas de origen societario como el arrastre y el acompañamiento, ofrecen soluciones de planificación patrimonial y sucesoria, financiamiento de vivienda y protección de personas, con riesgos que deben evaluarse.

En el plano tecnológico, el gobierno electrónico del año 2000 dio origen a la firma digital y a la protección de datos; la distinción entre firma electrónica y firma digital define quién debe probar, y la certificación notarial aporta certeza de identidad y controles que ninguna plataforma privada ofrece. La pérdida de confianza en el dinero estatal y el pensamiento cypherpunk dieron origen a Bitcoin y a blockchain, cuyo registro distribuido funciona como un protocolo notarial digital y desplaza la fe pública hacia la fe criptográfica. La tokenización inmobiliaria reproduce, en clave digital, la lógica del fideicomiso, dentro de un marco regulatorio transitorio a cargo de la CNV.

Frente a los riesgos ---estafas piramidales, plataformas inexistentes, lavado de activos--- el notariado y la abogacía deben asumir un rol activo de prevención y de asesoramiento calificado. El hilo conductor es un mensaje profesional: en un escenario cambiante, el valor diferencial no está en abrazar acríticamente ninguna tecnología ni en repetir fórmulas, sino en encuadrar bien cada problema, combinar creativamente las instituciones y saber hacer las preguntas correctas para proteger a los clientes.

% ============================================================
% LECTURAS Y FUENTES MENCIONADAS
% ============================================================
\chapter*{Lecturas y fuentes mencionadas}
\addcontentsline{toc}{chapter}{Lecturas y fuentes mencionadas}
\markboth{Lecturas y fuentes mencionadas}{}

Las referencias que siguen figuran en los apuntes originales. No se consignan datos adicionales a los allí indicados.

\section*{Normas citadas}
\begin{itemize}
    \item Código Civil y Comercial de la Nación (Ley 26.994): arts.~286, 287, 288, 308, 317 y 319; art.~1965 (dominio revocable); arts.~2114, 2124 y 2125 (derecho de superficie); régimen de la porción legítima (arts.~2444 y siguientes).
    \item Ley General de Sociedades N.º 19.550: régimen de tipicidad; Sección IV (arts.~21 a 26); art.~60; arts.~186 y 299.
    \item Decreto N.º 209/2024 (capital social mínimo).
    \item Ley N.º 27.349, de Apoyo al Capital Emprendedor.
    \item Ley N.º 24.441 (fideicomiso).
    \item Ley N.º 25.506, de Firma Digital (arts.~2, 5, 7 y 8) y Decreto N.º 182/2019, con su modificación de 2024.
    \item Ley N.º 25.326, de Protección de Datos Personales.
\end{itemize}

\section*{Doctrina y otras fuentes mencionadas}
\begin{itemize}
    \item Ricardo Nissen y Rafael Manóvil, como referentes de lecturas conservadoras y flexibles del derecho societario, respectivamente; Efraín Hugo Richard, referencia sobre libertad y responsabilidad societaria.
    \item Paula Sibilia, antropóloga argentina radicada en Brasil: lectura recomendada sobre los aspectos socio-técnicos de las redes sociales y el concepto de «extimidad».
    \item Glosario de términos técnicos elaborado en 2021 por la Comisión de Innovación y Tecnología del colegio profesional y publicado en el sitio web de la institución.
    \item Libro publicado en 2023 por la Comisión de Innovación y Tecnología del colegio profesional, surgido del formato «café notarial».
    \item Bitcoin Argentina, organización no gubernamental: fuente recomendada de formación gratuita, con cursos sobre el uso de billeteras digitales.
    \item White paper de Bitcoin (Satoshi Nakamoto): lectura básica sobre el funcionamiento del sistema.
\end{itemize}

\end{document}
